\documentclass[11pt]{article}

\usepackage[a4paper,margin=29mm]{geometry}
\usepackage[T1]{fontenc}
\usepackage{lmodern}
\usepackage{amsmath,amssymb,amsthm,mathtools}
\usepackage{booktabs,longtable,array}
\usepackage{enumitem}
\usepackage{microtype}
\usepackage{xcolor}
\usepackage[colorlinks=true,linkcolor=blue!55!black,citecolor=blue!55!black,
  urlcolor=blue!55!black]{hyperref}
\usepackage[nameinlink,noabbrev]{cleveref}

\newtheorem{theorem}{Theorem}[section]
\newtheorem{lemma}[theorem]{Lemma}
\newtheorem{proposition}[theorem]{Proposition}
\newtheorem{corollary}[theorem]{Corollary}
\theoremstyle{definition}
\newtheorem{definition}[theorem]{Definition}
\newtheorem{construction}[theorem]{Construction}
\theoremstyle{remark}
\newtheorem{remark}[theorem]{Remark}

\newcommand{\R}{\mathbb{R}}
\newcommand{\PP}{\mathbb{P}}
\newcommand{\cP}{\mathcal{P}}
\newcommand{\cC}{\mathcal{C}}
\newcommand{\cL}{\mathcal{L}}
\newcommand{\B}{\mathcal{B}}
\newcommand{\ind}{\mathbf{1}}
\newcommand{\floor}[1]{\left\lfloor #1\right\rfloor}
\newcommand{\ceil}[1]{\left\lceil #1\right\rceil}
\DeclareMathOperator{\supp}{supp}

\title{The minimum number of proper circles determined by a planar point set}
\author{Anonymous review manuscript}
\date{5 August 2026}

\begin{document}

\maketitle

\begin{abstract}
For a finite set \(P\subset\R^2\), let \(C(P)\) be the number of distinct
proper Euclidean circles containing at least three points of \(P\).  We
determine the minimum of \(C(P)\) over all \(n\)-point sets which are neither
collinear nor concyclic.  The three exceptional values are
\[
 f(6)=8,\qquad f(7)=11,\qquad f(8)=17,
\]
whereas
\[
 f(n)=1+\binom{n-1}{2}-\floor{\frac{n-1}{2}}
\]
for \(n=4,5\) and for every \(n\ge9\).

The proof is organized around unique ownership of triples by maximal line
and circle blocks.  Inversion converts the local block structure into real
line arrangements, allowing Melchior-type inequalities and their weighted
consequences.  A uniform half-size cap and one nonnegative-slack identity
settle all \(n\ge15\), with a separate two-case argument at \(n=14\).
The finite window is reduced to named real-projective rigidity lemmas and
small displayed certificates.  The proof uses no computerized search
through point configurations, block supports, or symmetry orbits.
\end{abstract}

\section{Introduction}

Erd\H{o}s asked for the least number of circles determined by \(n\) planar
points not all on a circle~\cite{Erdos1961}.  If collinear sets are admitted,
the minimum is trivially zero, so one must specify a nondegeneracy
convention.  We follow Elliott~\cite{Elliott1967}: the points are required
to be neither all collinear nor all concyclic.  A determined circle is
counted once, independently of how many selected points it contains.
This differs from the ordinary-circle problem, in which exactly three
selected points are required~\cite{LinEtAl2018}.

Purdy and Smith corrected the previously claimed large-\(n\) formula and
identified the natural extremal construction consisting of a circle and
its centre~\cite{PurdySmith2010}.  The present argument proves the exact
minimum at every finite value under Elliott's convention.

\begin{theorem}\label{thm:main}
For every integer \(n\ge4\),
\[
 (f(4),f(5),f(6),f(7),f(8))=(3,5,8,11,17),
\]
and
\[
 \boxed{\displaystyle
 f(n)=1+\binom{n-1}{2}-\floor{\frac{n-1}{2}}}
 \qquad(n\ge9).
\]
\end{theorem}

\subsection*{Architecture of the proof}

Every selected triple is contained in a unique maximal carrier: its line
when it is collinear, and its proper circle otherwise.  This elementary
partition is the common source of all counting identities in the paper.
Inverting about a selected point turns the carriers through that point into
spanned lines.  Melchior's inequality, applied before and after restoring
the inversion centre, then yields local slack variables and two global
weighted inequalities.

The large range is conceptually the shortest part.  A rich-carrier pencil
first bounds every block by \(\floor{n/2}\).  A single functional
\(S_M\), whose coefficients are visibly nonnegative at that cap, gives the
target for every \(n\ge15\).  At \(n=14\), the same functional handles the
absence of a seven-point circle and one selected-circle inequality handles
its presence.

For \(4\le n\le13\), the universal identities leave a finite number of
equality mechanisms.  We isolate their geometry as reusable lemmas:
the real Fano and Orchard walls, pentagon and golden-axis rigidity,
circle-lift trace and conic-event bounds, and the grid and gallery
obstructions.  The remaining arithmetic consists of fixed coefficient
identities or small structural certificates displayed in the text.

\subsection*{Classical and computational boundaries}

The external incidence inputs are stated with their exact hypotheses at
first use.  They are Melchior's inequality~\cite{Melchior1940}, the
non-strict Kelly--Moser ordinary-line estimate~\cite{KellyMoser1958},
Langer's inequality in the form recorded by de Zeeuw
~\cite{Langer2003,DeZeeuw2018}, Lenchner's affine ordinary-point theorem
~\cite{Lenchner2007}, and Ungar's direction theorem~\cite{Ungar1982}.
Inversion, the
circle lift, polarity, and the local conic arguments are proved within the
paper rather than treated as black boxes.

Exact-arithmetic checks used during preparation replay displayed integer
identities, rational coefficients, and fixed determinants.  They do not
discover cases and do not supply a geometric premise.  Two small
finite residues remain explicit: a nine-row coefficient table in the
ten-point six-circle case, and a four-state corner propagation in the
twelve-line gallery.  These are hand-checkable certificates, not
enumerations of configurations or realizations.

\subsection*{Organization}

\Cref{sec:foundations-small} develops the block calculus, records the
classical inputs and upper constructions, and proves the cases
\(4\le n\le9\).  \Cref{sec:finite-n10-n12} proves the three most geometric
finite cases \(n=10,11,12\).  \Cref{sec:n13-infinity} treats \(n=13\), the
two-case endpoint \(n=14\), and the uniform range \(n\ge15\), and records
the fixed-certificate boundary.

\section{Foundations and the cases \texorpdfstring{$4\leq n\leq9$}{4 <= n <= 9}}
\label{sec:foundations-small}

This section fixes the conventions used throughout the paper, derives the
universal counting rows, and proves all values up to nine.  A line is used as
a bookkeeping carrier, but is never counted as a circle.

\subsection{The problem and the block language}

\begin{definition}[The V1 problem]\label{def:v1-admissible}
Let $P\subset\R^2$ be a set of $n\geq4$ distinct points.  We call $P$
\emph{admissible} if it is neither collinear nor contained in one proper
Euclidean circle.  A \emph{determined circle} is a proper circle containing a
noncollinear triple of $P$.  Write $C(P)$ for the number of distinct
determined circles and
\[
  f(n):=\min\{C(P): |P|=n\text{ and }P\text{ is admissible}\}.
\]
\end{definition}

For a line or a proper circle $X$, its \emph{full trace} is $P\cap X$.
A maximal generalized $s$-block is a full trace of size $s\geq3$.  We write
\[
 L_s:=\#\{\text{$s$-point line blocks}\},\qquad
 C_s:=\#\{\text{$s$-point circle blocks}\},\qquad
 B_s:=L_s+C_s.
\]
Thus
\begin{equation}\label{eq:circle-objective}
 C(P)=\sum_{s\geq3}C_s.
\end{equation}
For $p\in P$, let $\ell_s(p)$, $c_s(p)$, and
$d_s(p)=\ell_s(p)+c_s(p)$ count, respectively, the line, circle, and all
$s$-blocks through $p$.

\begin{lemma}[Unique ownership]\label{lem:triple-ownership}
Every triple of $P$ belongs to exactly one maximal generalized block.
Consequently
\begin{align}
 \mathsf T
   &:=\sum_{s\geq3}\binom{s}{3}B_s=\binom n3,
       \label{eq:row-T}\\
 \sum_{p\in P}d_s(p)&=sB_s,
 &\sum_{p\in P}\ell_s(p)&=sL_s.                  \label{eq:block-incidence}
\end{align}
If $L_2$ denotes the number of ordinary, two-point full line traces, then
\begin{equation}\label{eq:line-pair-partition}
 L_2+\sum_{s\geq3}\binom{s}{2}L_s=\binom n2.
\end{equation}
\end{lemma}

\begin{proof}
A collinear triple has its unique carrier line; a noncollinear triple has its
unique circumcircle.  Passing to the full trace gives a unique maximal owner.
Equation~\eqref{eq:row-T} counts triples by that owner, and
\eqref{eq:block-incidence} counts point--block incidences.  Finally, every
pair of distinct points has one carrier line, which proves
\eqref{eq:line-pair-partition}.
\end{proof}

\begin{lemma}[The eight-point Orchard bound]\label{lem:orchard-eight}
Let \(Q\) be a set of eight noncollinear points of the real projective plane,
with no four collinear.  If \(t_i\) denotes the number of lines containing
exactly \(i\) points of \(Q\), then
\[
 t_3\leq 7.
\]
\end{lemma}

\begin{proof}
Pair counting gives \(t_2+3t_3=28\).  Melchior's inequality excludes
\(t_3\geq9\), so it remains only to rule out \(t_3=8\).  In that case
\(t_2=4\).  At a point \(v\in Q\), let \(r_v\) and \(s_v\) be the numbers
of three-point and ordinary lines through \(v\).  Then
\[
 7=2r_v+s_v.
\]
Thus \(s_v\) is a positive odd integer, while
\(\sum_v s_v=2t_2=8\).  Consequently \(s_v=1\) and \(r_v=3\) for every
\(v\).  The four ordinary lines are therefore a perfect matching; denote
its pairs by
\[
 A=\{A_0,A_1\},\quad B=\{B_0,B_1\},\quad
 C=\{C_0,C_1\},\quad D=\{D_0,D_1\}.
\]
Every three-point line uses three different pairs.  Since the eight such
lines account for all \(24\) cross-pair incidences, exactly two omit each
of \(A,B,C,D\).  Explicitly, if \(n_A,n_B,n_C,n_D\) are the four
omission counts, then the four edges between, say, \(A\) and \(B\) give
\(n_C+n_D=4\); the six analogous equations force
\(n_A=n_B=n_C=n_D=2\).

We first determine the resulting incidence structure.  Two triples on the
same three pairs cannot differ in only one entry, because then their two
common points would determine both lines.  They also cannot differ in
exactly two entries.  Indeed, after relabelling suppose the two \(ABC\)
triples are
\[
 A_0B_0C_0,\qquad A_1B_1C_0.
\]
The unused cross edges force the remaining triples to have the form
\[
\begin{array}{ll}
 A_0B_1D_x,&A_1B_0D_y,\\
 A_0C_1D_z,&A_1C_1D_{1-z},\\
 B_0C_1D_w,&B_1C_1D_{1-w}.
\end{array}
\]
The \(AD\), \(BD\), and \(CD\) edges successively give
\(y=z=1-x\) and \(w=x\), after which one of the two \(C_1D_i\) edges is
repeated.  Hence the two triples on any fixed three pairs are complementary.
After fixing the \(ABC\) pair as \(000,111\), the unused \(AB,AC,BC\)
edges are the three cross-matchings.  A flip of the \(D\)-labels fixes one
\(ABD\) triple, and the unused \(AD\) and \(BD\) edges then force all
remaining bits.  Up to relabelling, the unique incidence list is
\begin{equation}\label{eq:orchard-eight-incidence}
\begin{split}
&A_0B_0C_0,\quad A_1B_1C_1,\quad
 A_0B_1D_0,\quad A_1B_0D_1,\\
&A_0C_1D_1,\quad A_1C_0D_0,\quad
 B_0C_1D_0,\quad B_1C_0D_1.
\end{split}
\end{equation}

Choose the projective frame
\[
 A_0=(1,0,0),\quad A_1=(0,1,0),\quad
 B_0=(0,0,1),\quad B_1=(1,1,1).
\]
The first four incidences in \eqref{eq:orchard-eight-incidence} allow the
normalization
\[
 C_0=(1,0,u),\quad C_1=(1,v,1),\quad
 D_0=(w,1,1),\quad D_1=(0,1,z),
\]
where the relevant parameters are nonzero.  The next incidences give
\[
 vz=uw=vw=1.
\]
Thus \(u=v\) and \(w=z=1/u\).  The last incidence,
\(B_1C_0D_1\), now gives
\[
 u^2-u+1=0,
\]
which has no real root.  Hence \(t_3=8\) is impossible.
\end{proof}

\begin{lemma}[Radical-axis cross-block bound]
\label{lem:radical-axis-cross-block}
Let \(A\) and \(B\) be marked points on two distinct real circles, after
deleting any marked point common to the circles.  A \emph{cross-block} is a
line or circle containing two points of \(A\) and two points of \(B\).
If \((|A|,|B|)=(5,4)\), there are at most \(12\) cross-blocks.  If
\((|A|,|B|)=(4,4)\), there are at most \(10\).
\end{lemma}

\begin{proof}
For each chord of the first circle and each chord of the second, the four
endpoints are concyclic precisely when the two supporting lines meet on the
radical axis.  A coincident pair of supporting lines gives a line
cross-block and is assigned to its intersection with the same radical
axis.  (The common supporting line cannot itself be the radical axis,
because the marked common points were deleted.)

At a point \(q\) of the radical axis, let \(a_q\) and \(b_q\) be the
numbers of selected chords from \(A\) and \(B\) through \(q\).  Chords
through a fixed external point form a matching, so
\[
 a_q\leq\floor{|A|/2},\qquad b_q\leq\floor{|B|/2}.
\]
Moreover, \(\sum_q a_q=\binom{|A|}{2}\),
\(\sum_q b_q=\binom{|B|}{2}\), and the number of cross-blocks is
\[
 M=\sum_q a_qb_q.
\]
For \(5+4\), this gives \(M\leq2\sum_qb_q=12\).

For \(4+4\), both multiplicity sums are \(6\), and \(M\leq12\).  Equality
\(M=11\) is impossible: from
\[
 12-M=\sum_q a_q(2-b_q)=1
\]
there would be one point with \(a_q=b_q=1\), while every other point with
\(a_q>0\) had \(b_q=2\).  If \(r\) and \(s\) count the remaining points
with \(a_q=1\) and \(a_q=2\), respectively, then
\[
 1+r+2s=6,\qquad 1+2r+2s\leq6,
\]
forcing \(r=0\) and then \(1+2s=6\), a contradiction.  Equality \(M=12\)
would force exactly three points \(q\), with
\(a_q=b_q=2\) at each.  These are the three diagonal points of each of the
two complete quadrangles, yet all three lie on the radical axis.  The
diagonal points of a nondegenerate complete quadrangle are noncollinear
(in the standard frame their determinant is \(-2\)).  Thus \(M\leq10\).
\end{proof}

\subsection{Sharp upper constructions}

\begin{construction}[A circle and its centre]\label{con:circle-centre}
Put $r=n-1$ distinct points on a proper circle $\Gamma$ and add its centre
$o$.  Choose exactly $\floor{r/2}$ antipodal boundary pairs; if $r$ is odd,
choose the remaining boundary point without its antipode.  The boundary
triples determine only $\Gamma$.  A triple consisting of $o$ and a boundary
pair is collinear precisely for an antipodal pair.  Every other boundary
pair gives a different proper circle, since a circle different from
$\Gamma$ meets $\Gamma$ in at most two points.  Thus this admissible set
determines exactly
\begin{equation}\label{eq:circle-centre-count}
 F(n):=1+\binom{n-1}{2}-\floor{\frac{n-1}{2}}
\end{equation}
 circles.  In particular it gives the sharp upper targets $3,5$ for
 $n=4,5$ and $25$ for $n=9$.
\end{construction}

 The three exceptional upper bounds have equally explicit fixed certificates.
For points $p_i=(x_i,y_i)$, write
\begin{equation}\label{eq:concyclic-determinant}
 \Delta(i,j,k,l):=
 \det\begin{pmatrix}
 x_i^2+y_i^2&x_i&y_i&1\\
 x_j^2+y_j^2&x_j&y_j&1\\
 x_k^2+y_k^2&x_k&y_k&1\\
 x_l^2+y_l^2&x_l&y_l&1
 \end{pmatrix}.
\end{equation}
For a noncollinear quadruple, $\Delta=0$ is equivalent to concyclicity.

\begin{construction}[Six rational points and eight circles]
\label{con:six-upper}
Take
\[
\begin{array}{c|cccccc}
i&0&1&2&3&4&5\\ \hline
p_i&(0,0)&(4,0)&(1,0)&(1,3)&
 (\tfrac25,\tfrac65)&(2,2).
\end{array}
\]
The collinear triples are exactly
\[
 012,\qquad034,\qquad135.
\]
Exact substitution in $x^2+y^2+Dx+Ey+F=0$ gives the following eight
circle traces and coefficient triples $(D,E,F)$:
\[
\begin{array}{c|c}
013&(-4,-2,0)\\
0145&(-4,0,0)\\
0235&(-1,-3,0)\\
024&(-1,-1,0)\\
1234&(-5,-3,4)\\
125&(-5,-1,4)\\
245&(-3,-2,2)\\
345&(-2,-4,4).
\end{array}
\]
The three four-circles own four triples each and the other five circles own
one triple each.  Together with the three collinear triples this accounts for
all $\binom63=20$ triples.  Thus the configuration is admissible and
determines exactly $8$ circles.
\end{construction}

\begin{construction}[Seven points and eleven circles]
\label{con:seven-upper}
Let $u=\sqrt3$ and take
\[
\begin{array}{c|ccccccc}
i&0&1&2&3&4&5&6\\ \hline
p_i&(-1,0)&(1,0)&(0,u)&(0,0)&
(-\tfrac12,\tfrac u2)&(\tfrac12,\tfrac u2)&(0,\tfrac u3).
\end{array}
\]
The collinear triples are exactly
\[
 013,\ 024,\ 056,\ 125,\ 146,\ 236,
\]
where a string denotes the corresponding set of indices.  Direct expansion
of \eqref{eq:concyclic-determinant}, using the dihedral symmetry of the
equilateral triangle, gives exactly the six proper four-circles
\[
 0145,\ 0235,\ 0346,\ 1234,\ 1356,\ 2456.
\]
Their $24$ triples are disjoint.  The five remaining noncollinear triples
are
\[
 012,\ 016,\ 026,\ 126,\ 345.
\]
The determinant list also shows that they have five different owners and
that no five-point circle occurs.  Hence the configuration determines
$6+5=11$ circles.
\end{construction}

\begin{construction}[Eight rational points and seventeen circles]
\label{con:eight-upper}
Take
\[
\begin{array}{c|cccccccc}
i&0&1&2&3&4&5&6&7\\ \hline
p_i&(0,0)&(1,0)&(\tfrac13,0)&(\tfrac23,0)&
(\tfrac23,\tfrac13)&(\tfrac25,\tfrac15)&
(\tfrac13,\tfrac13)&(\tfrac35,\tfrac15).
\end{array}
\]
Its collinear triples are exactly
\[
 012,\ 013,\ 023,\ 123,\ 045,\ 167.
\]
Exact rational evaluation of \eqref{eq:concyclic-determinant} gives exactly
the eleven proper four-circles
\[
\begin{gathered}
0146,\ 0157,\ 0247,\ 0256,\ 0367,\ 1245,\\
1347,\ 1356,\ 2346,\ 2357,\ 4567.
\end{gathered}
\]
No five-point subset has all its quadruples on this list.  The six remaining
noncollinear triples are
\[
 034,\ 035,\ 126,\ 127,\ 267,\ 345.
\]
Thus the eleven four-circles own $44$ triples and the last six triples have
six distinct owners, for a total of $17$ circles.
\end{construction}

\subsection{The elementary cases \texorpdfstring{$n=4,5$}{n=4,5}}

\begin{proposition}\label{prop:n4-n5}
One has $f(4)=3$ and $f(5)=5$.
\end{proposition}

\begin{proof}
Let first $|P|=4$.  If three points lie on a line, the fourth point together
with each of the three pairs on that line gives a proper circle; the three
circles are distinct.  If no triple is collinear, all four triples are
noncollinear, and two of them could share a circle only if all four points
were concyclic.  Hence $C(P)\geq3$.  Construction
\ref{con:circle-centre} attains three.

Now let $|P|=5$.  A four-point line and its outsider give six distinct
circles, so assume that no line has four points.  Distinct three-point lines
use disjoint pairs; therefore $3L_3\leq\binom52$ and $L_3\leq3$.  If no
four-point circle occurs, every noncollinear triple has a different owner,
and
\[
 C(P)=\binom53-L_3\geq7.
\]
If a four-point circle $\Gamma$ occurs, it is unique.  Every three-point line
contains the outsider and a chord of $\Gamma$, and these chords form a
matching, so there are at most two such lines.  At least eight triples are
noncollinear; four lie on $\Gamma$, while the other four have mutually
different owners, since a repeated owner would be a second four-point
circle.  Thus $C(P)\geq1+4=5$, attained by
\cref{con:circle-centre}.
\end{proof}


We repeatedly use, without further comment, that two distinct proper circles
meet in at most two points, two distinct lines in at most one point, and a
line meets a proper circle in at most two points.

\subsection{Classical incidence inputs and their exact scope}
\label{subsec:classical-inputs}

For later reference we record every external incidence theorem used in the
proof.  The hypotheses below are part of the statements; in particular, none
of these results assumes general position.

\begin{theorem}[Melchior]\label{thm:melchior}
Let $\mathcal A$ be a finite arrangement of distinct real projective lines,
not all concurrent, and let $t_r(\mathcal A)$ count vertices incident with
exactly $r$ lines.  Then
\begin{equation}\label{eq:melchior}
 t_2\geq3+\sum_{r\geq4}(r-3)t_r.
\end{equation}
Equivalently, the same inequality holds for the multiplicities of lines
determined by a finite set of distinct noncollinear real projective points.
All projective vertices, including those at infinity, are counted.
\end{theorem}

This is the point-dual form of Melchior's theorem~\cite{Melchior1940}.  We do
not import a classification of its equality cases.

\begin{theorem}[Kelly--Moser]\label{thm:kelly-moser}
If $Q$ is a finite set of $q$ distinct noncollinear real points and $t_2$ is
the number of its ordinary determined lines, then
\begin{equation}\label{eq:kelly-moser}
 3q\leq7t_2,
 \qquad
 t_2\geq\ceil{\frac{3q}{7}}.
\end{equation}
The inequality is non-strict; no equality classification is assumed.
\end{theorem}

This is Theorem~3.6 of Kelly and Moser~\cite{KellyMoser1958}.

\begin{theorem}[Langer--de Zeeuw form]\label{thm:langer}
Let $Q$ be a finite set of $N$ distinct affine points over $\mathbb C$.
Assume that every affine line contains at most $2N/3$ points of $Q$.  If
$t_r$ counts the spanned lines containing exactly $r$ points, then
\begin{equation}\label{eq:langer}
 \sum_{r\geq2}r t_r\geq\frac{N(N+3)}{3}.
\end{equation}
In particular the result applies to real point sets.  Both the occupancy cap
and the conclusion are non-strict.
\end{theorem}

We use the formulation of de Zeeuw~\cite{DeZeeuw2018}, derived from
Langer's orbifold inequality~\cite{Langer2003}.

\begin{theorem}[Lenchner]\label{thm:lenchner}
Let $\mathcal A$ be an affine arrangement of $q$ distinct real lines, neither
all parallel nor all concurrent.  Apart from one specified exceptional
six-line arrangement, the number $o_{\mathrm{fin}}(\mathcal A)$ of finite
ordinary intersections satisfies
\begin{equation}\label{eq:lenchner}
 7o_{\mathrm{fin}}(\mathcal A)\geq2q-3.
\end{equation}
Consequently \eqref{eq:lenchner} is available under the simpler guard
$q\neq6$.
\end{theorem}

The word \emph{finite} is essential.  This is the corollary of Lenchner's
Theorem~8~\cite{Lenchner2007} used later in the paper.

\begin{theorem}[Ungar, six-point consequence]\label{thm:ungar-six}
Six distinct noncollinear real affine points determine at least six
directions, where a direction is a parallelism class of determined lines.
\end{theorem}

This is the only consequence of Ungar's direction theorem that we use
\cite{Ungar1982}.

\begin{theorem}[Hall]\label{thm:hall}
Let $G$ be a finite bipartite graph with parts $X,Y$.  If
$|S|\leq|N(S)|$ for every $S\subseteq X$, then $G$ has a matching saturating
$X$.
\end{theorem}

We use Hall's theorem~\cite{Hall1935} only after the relevant neighbourhood
inequalities have been established geometrically.

Inversion, projective duality, directed Menelaus, and power of a point are
not treated as external incidence black boxes.  The inversion statements
needed for the universal rows are proved next; later geometric uses are
proved at their point of application.

\subsection{Inversion and the universal rows}

Fix $p\in P$ and $\rho>0$.  On $\R^2\setminus\{p\}$ put
\[
 I_{p,\rho}(x)=p+\frac{\rho}{\|x-p\|^2}(x-p).
\]

\begin{lemma}[Inversion dictionary]\label{lem:inversion-dictionary}
The map $I_{p,\rho}$ is an injective involution on
$\R^2\setminus\{p\}$.  A line through $p$ maps to itself, while a proper
circle through $p$ maps to a line avoiding $p$, and conversely.  Hence every
maximal $s$-block through $p$ becomes a full line containing exactly $s-1$
points of the inverted set
\[
 Q_p:=I_{p,\rho}(P\setminus\{p\}).
\]
The set $Q_p$ is noncollinear.  It remains noncollinear after the inversion
centre $o=p$ is added.
\end{lemma}

\begin{proof}
The formula immediately gives $I_{p,\rho}^2=\mathrm{id}$.  Translating
$p$ to the origin and substituting the formula into a linear or quadratic
equation gives the stated line--circle dictionary.  Full traces are
preserved because the map is injective.

If $Q_p$ lay on a line through $o$, then all original points would lie on
that line.  If it lay on a line avoiding $o$, the inverse carrier would be a
proper circle through $p$ containing all of $P$.  Both possibilities violate
admissibility.  Adding $o$ cannot make a noncollinear set collinear.
\end{proof}

\begin{corollary}[The Lenchner pivot bridge]\label{cor:pivot-lenchner}
Let \(p\in P\), put \(q=n-1\), and assume \(q\neq6\).  Then at least
\[
 \ceil{\frac{2q-3}{7}}
\]
proper three-point circles of \(P\) pass through \(p\).  Equivalently,
\[
 c_3(p)\geq\ceil{\frac{2n-5}{7}}.
\]
\end{corollary}

\begin{proof}
Invert \(P\setminus\{p\}\) at \(p\), obtaining the \(q\) distinct,
noncollinear image points \(Q_p\) from
\cref{lem:inversion-dictionary}.  Work projectively, dualize those points,
and declare the dual \(p^*\) of the inversion centre to be the line at
infinity.  This produces \(q\) distinct real affine lines.

The dual lines are not all parallel: otherwise the points of \(Q_p\)
would be collinear on a line through \(p\).  They are not all concurrent
at a finite point: otherwise \(Q_p\) would be collinear on a line avoiding
\(p\).  Thus the arrangement satisfies the geometric hypotheses of
\cref{thm:lenchner}; the guard \(q\neq6\) also removes its sole exceptional
cardinality.  It consequently has at least
\(\ceil{(2q-3)/7}\) finite ordinary intersections.

Under projective duality, such an intersection is an ordinary determined
line of \(Q_p\).  Finiteness says exactly that this line avoids \(p\):
ordinary intersections on \(p^*\) correspond instead to image lines
through \(p\), and are deliberately not counted.  By inversion, an ordinary
image line avoiding \(p\) is a proper circle whose full trace on \(P\)
consists of \(p\) and the two preimages.  Distinct image lines give distinct
circles, proving the claim.
\end{proof}

The inversion dictionary verifies every noncollinearity hypothesis in the next
two applications of \cref{thm:melchior}.

\begin{proposition}[Universal Melchior rows]\label{prop:universal-melchior-rows}
For every admissible $n$-point set and every $p\in P$, define
\begin{equation}\label{eq:pivot-slack}
 \sigma_p:=d_3(p)-3-\sum_{s\geq5}(s-4)d_s(p).
\end{equation}
Then $\sigma_p\geq0$, and
\begin{align}
 \mathsf P
   &:=\sum_{s\geq3}s(4-s)B_s\geq3n,
      \label{eq:row-P}\\
 \sum_{p\in P}\sigma_p&=\mathsf P-3n.             \label{eq:sum-pivot-slack}
\end{align}
If $\ell_2(p)$ counts the ordinary original lines through $p$, then
\begin{align}
 \ell_2(p)+\sum_{s\geq3}(s-1)\ell_s(p)&=n-1,
      \label{eq:pivot-arms}\\
 \sum_{s\geq3}s\ell_s(p)&\leq n-1+\sigma_p.
      \label{eq:added-centre-local}
\end{align}
After summing \eqref{eq:added-centre-local}, one obtains
\begin{equation}\label{eq:row-D}
 \mathsf D
 :=\sum_{s\geq3}s(s-4)C_s
   +\sum_{s\geq3}2s(s-2)L_s
 \leq n(n-4).
\end{equation}
Finally, direct application of Melchior to $P$ gives
\begin{equation}\label{eq:global-line-row}
 3L_3+\sum_{s\geq4}
   \left(\binom{s}{2}+s-3\right)L_s
 \leq\binom n2-3.
\end{equation}
\end{proposition}

\begin{proof}
Apply \cref{thm:melchior} to the $n-1$ inverted points.  By
\cref{lem:inversion-dictionary}, its $(s-1)$-point lines are exactly the
$s$-blocks through $p$.  Melchior is therefore precisely
$\sigma_p\geq0$.  Summing over $p$ and using
\eqref{eq:block-incidence} gives \eqref{eq:row-P} and
\eqref{eq:sum-pivot-slack}.

The original lines through $p$ partition $P\setminus\{p\}$ into arms,
which proves \eqref{eq:pivot-arms}.  Add the inversion centre $o$.  The
line multiplicities of the resulting $n$-point set are exactly
\[
 t_2=c_3(p)+\ell_2(p),\qquad
 t_i=c_{i+1}(p)+\ell_i(p)\quad(i\geq3).
\]
The slack in Melchior is then
\[
 \sigma_p+\ell_2(p)-\sum_{s\geq3}\ell_s(p)\geq0.
\]
Using \eqref{eq:pivot-arms} yields
\eqref{eq:added-centre-local}.  Sum it over $p$, use
\eqref{eq:block-incidence} and \eqref{eq:sum-pivot-slack}, and subtract
$\mathsf P$ from the resulting inequality.  The circle coefficient is
$-s(4-s)=s(s-4)$ and the line coefficient is
$s^2-s(4-s)=2s(s-2)$, giving \eqref{eq:row-D}.

For the last assertion, substitute \eqref{eq:line-pair-partition} in
\eqref{eq:melchior} applied directly to $P$.
\end{proof}

\begin{proposition}[Conditional Langer row]\label{prop:universal-langer-row}
Suppose that, for every pivot $p$, every line determined by $Q_p$ contains
at most $2(n-1)/3$ image points.  Then
\begin{equation}\label{eq:row-L}
 \mathsf L:=\sum_{s\geq3}s(s-1)B_s
 \geq\frac{n(n-1)(n+2)}{3}.
\end{equation}
\end{proposition}

\begin{proof}
Apply \cref{thm:langer} to the $N=n-1$ image points at each pivot.  Its
incidence count is $\sum_{s\geq3}(s-1)d_s(p)$.  Summing over $p$ and using
\eqref{eq:block-incidence} proves \eqref{eq:row-L}.  Notice that the
occupancy hypothesis must be verified before this row is used.
\end{proof}

Equations \eqref{eq:row-T}, \eqref{eq:row-P}, \eqref{eq:row-D}, and,
under its stated cap, \eqref{eq:row-L}, are the four universal rows imported
by the later sections.

\subsection{Two pencil bounds}

\begin{lemma}[Rich-line pencil]\label{lem:rich-line-pencil}
If a line contains $m\geq3$ points of an admissible $n$-point set and
$q=n-m\geq1$ points lie off it, then
\begin{equation}\label{eq:rich-line-pencil}
 C(P)\geq q\binom m2-\binom q2\floor{\frac m2}.
\end{equation}
\end{lemma}

\begin{proof}
For a fixed outsider, its circles with pairs on the line are proper and
mutually distinct.  For two fixed outsiders, common circles use disjoint
pairs of line points: two pairs sharing a point would give two circles with
three common points.  Thus two outsider pencils overlap in at most
$\floor{m/2}$ circles.  The two-term Bonferroni inequality gives
\eqref{eq:rich-line-pencil}; higher intersections do not invalidate this
lower bound.
\end{proof}

\begin{lemma}[Rich-circle pencil]\label{lem:rich-circle-pencil}
If a proper circle $\Gamma$ contains $m\geq3$ selected points and
$q=n-m\geq1$ points lie outside it, then
\begin{equation}\label{eq:rich-circle-pencil}
 C(P)\geq
 1+q\left(\binom m2-\floor{\frac m2}\right)
   -\binom q2\floor{\frac m2}.
\end{equation}
\end{lemma}

\begin{proof}
For an outsider $x$, the chords of $\Gamma$ passing through $x$ form a
matching on the $m$ marked points, so at most $\floor{m/2}$ marked pairs are
collinear with $x$.  Every other pair gives a distinct proper circle through
$x$.  As in the line case, two outsider pencils overlap in at most
$\floor{m/2}$ circles.  Add $\Gamma$ and apply Bonferroni.
\end{proof}

\subsection{Six points}

\begin{proposition}\label{prop:n6}
One has $f(6)=8$.
\end{proposition}

\begin{proof}
Construction~\ref{con:six-upper} gives the upper bound.  Suppose, for a
contradiction, that an admissible six-point set satisfies $C(P)\leq7$, and
let $k$ be the maximum number of selected points on a proper circle.

If $k=5$, \cref{lem:rich-circle-pencil} gives $C(P)\geq9$.  If $k=3$,
every noncollinear triple has a different circle owner.  Let a largest line
contain $m$ points and put $r=6-m$.  For $m\geq4$, a collinear triple not
contained in that line uses at most one of its points, and an outsider pair
can be collinear with at most one point on the line.  Thus the number $T_0$
of collinear triples satisfies
\[
 T_0\leq\binom m3+\binom r3+\binom r2.
\]
This leaves at least $10$ noncollinear triples for $m=5$ and at least $15$
for $m=4$.  If all lines have size at most three, their collinear triples use
disjoint sets of three pairs, so there are at most
$\floor{\binom62/3}=5$ of them.  Again at least $15$ circles remain.

It remains that $k=4$.  The rich-line bound is $10$ for both $m=4$ and
$m=5$, so every line has size at most three.  All maximal generalized blocks
therefore have size three or four, and every four-block is a circle.  Put
\[
 b=C_4,\qquad L=L_3.
\]
Triple ownership, the pivot Melchior row, the exact circle count, and the
global line row give, respectively,
\begin{align}
 B_3+4b&=20,                                      \label{eq:n6-triples}\\
 0\leq\sum_{p\in P}(d_3(p)-3)&=42-12b,           \label{eq:n6-pivot}\\
 C(P)&=20-3b-L,                                   \label{eq:n6-circles}\\
 L&\leq4.                                         \label{eq:n6-lines}
\end{align}
Hence $b\leq3$, while $C(P)\leq7$ forces
\begin{equation}\label{eq:n6-terminal-rows}
 b=3,\qquad L=4.
\end{equation}

At any $p\in P$, ownership of the ten triples containing $p$ reads
\[
 d_3(p)+3d_4(p)=10.
\]
The first Melchior row gives $d_3(p)\geq3$, and $d_3(p)\equiv1\pmod3$.
Since $\sum_p d_4(p)=4b=12$, necessarily
\begin{equation}\label{eq:n6-local-profile}
 d_3(p)=4,\qquad d_4(p)=2,\qquad \sigma_p=1
 \quad\text{for every }p.
\end{equation}
The added-centre inequality \eqref{eq:added-centre-local} becomes
\begin{equation}\label{eq:n6-kappa}
 \kappa_p:=6-3\ell_3(p)\geq0.
\end{equation}

Since $L=4$, summing \eqref{eq:n6-kappa} gives zero, so every point lies on
exactly two of the four three-point lines.  No three of those lines are
concurrent, and their six pairwise intersections are the selected points:
the line pattern is a complete quadrilateral.  Label its four lines
\[
 012,\qquad034,\qquad135,\qquad245.
\]
Apply a Euclidean similarity, which preserves lines and circles, so that
\begin{equation}\label{eq:n6-pasch-coordinates}
\begin{gathered}
p_0=(0,0),\quad p_1=(1,0),\quad p_3=(u,v),\qquad v\ne0,\\
p_2=(a,0),\quad p_4=b(u,v),\quad
p_5=t(1,0)+(1-t)(u,v),\qquad
t=\frac{a(b-1)}{b-a}.
\end{gathered}
\end{equation}
Here distinctness gives $a,b\notin\{0,1\}$.  Moreover $a\neq b$: if
$a=b$, then the line $p_2p_4$ is parallel to $p_1p_3$ (their direction
vectors differ by the nonzero scalar $a$), and the two lines are distinct
because $a\neq1$; they therefore cannot share the affine point $p_5$.
Put
$R=u^2+v^2>0$.

By \eqref{eq:n6-local-profile}, the complements of the three four-circles
are three disjoint pairs partitioning $P$.  Each complement meets every one
of the four collinear triples.  In a complete quadrilateral the only such
pairs are the three diagonal pairs, so the circle supports are forced to be
\[
 0145,\qquad0235,\qquad1234.
\]
Substitution in \eqref{eq:concyclic-determinant} gives, in the same order,
\[
\begin{split}
 \Delta(0,1,4,5)
  &=\frac{vb^2(a-1)(b-1)}{(a-b)^2}(-Rb+2au-a),\\
 \Delta(0,2,3,5)
  &=\frac{va^2(a-1)(b-1)}{(a-b)^2}(-Rb-a+2bu),\\
 \Delta(1,2,3,4)
  &=v(a-1)(b-1)(-Rb+a).
\end{split}
\]
All prefactors displayed here are nonzero.  Thus concyclicity gives
\[
 a(1-2u)+bR=0,\qquad
 a+b(R-2u)=0,\qquad
 a-bR=0.
\]
The third equation gives $a=bR$.  Substitution in the second gives
$2b(R-u)=0$, hence $R=u$.  The first equation then becomes
$2bu(1-u)=0$.  Since $b\ne0$ and $u=R>0$, one has $u=1$.
But $R=u=1$ and $R=u^2+v^2$ force $v=0$, contradicting
\eqref{eq:n6-pasch-coordinates}.  Thus $C(P)\geq8$.
\end{proof}

\subsection{Seven points and the planar Fano wall}

\begin{lemma}[Planar Fano wall]\label{lem:fano-seven-wall}
There do not exist seven distinct real points together with seven proper
four-point circles whose three-point complements form a Steiner triple
system on the seven points.
\end{lemma}

\begin{proof}
We first derive, rather than quote, the unique form of a Steiner triple
system on seven labels.  Name one triple $123$.  The two remaining triples
through $1$ partition $\{4,5,6,7\}$ into two pairs, so after relabelling they
are $145$ and $167$.  The triples through $2$ use neither pair $45$ nor pair
$67$; after possibly interchanging $6,7$, they are $246$ and $257$.  The
unused pairs on $\{4,5,6,7\}$ are then $47$ and $56$, forcing the final two
triples through $3$.  Thus the system is
\begin{equation}\label{eq:fano-triples}
 123,\quad145,\quad167,\quad246,\quad257,\quad347,
 \quad356.
\end{equation}

Invert at point $1$.  The four triples in \eqref{eq:fano-triples} avoiding
$1$ have complementary four-circles through $1$, which become the lines
\begin{equation}\label{eq:fano-inverted-lines}
 357,\qquad346,\qquad256,\qquad247.
\end{equation}
The first three form a nondegenerate triangle and the fourth is a
transversal.  Put
\[
 A=3,\quad B=5,\quad C=6,\qquad
 X=7\in AB,\quad Y=4\in AC,\quad Z=2\in BC.
\]
The complements of the three Fano triples through $1$ become the proper
circles
\begin{equation}\label{eq:fano-inverted-circles}
 ABYZ,\qquad ACXZ,\qquad BCXY.
\end{equation}

Use directed affine parameters
\[
 AB:A=0,B=1,X=x,\quad
 AC:A=0,C=1,Y=y,\quad
 BC:B=0,C=1,Z=z.
\]
All six quantities $x,y,z,1-x,1-y,1-z$ are nonzero.  Put
$a^2=|BC|^2$, $b^2=|CA|^2$, and $c^2=|AB|^2$.  Power of a point applied at
the opposite vertices to the three circles in
\eqref{eq:fano-inverted-circles} gives
\[
 b^2(1-y)=a^2(1-z),\qquad
 c^2(1-x)=a^2z,\qquad
 c^2x=b^2y.
\]
Multiplying the resulting ratios yields
\begin{equation}\label{eq:fano-power-product}
 \frac{x}{1-x}\frac{z}{1-z}\frac{1-y}{y}=+1.
\end{equation}
But $X,Y,Z$ are collinear by the fourth line in
\eqref{eq:fano-inverted-lines}, so directed Menelaus in triangle $ABC$
gives the same left-hand side equal to $-1$.  This contradiction proves the
lemma.
\end{proof}

\begin{proposition}\label{prop:n7}
One has $f(7)=11$.
\end{proposition}

\begin{proof}
Construction~\ref{con:seven-upper} gives $f(7)\leq11$.  Suppose that an
admissible seven-point set satisfies $C(P)\leq10$.  The two pencil bounds
give
\[
 \begin{array}{c|ccc}
 \text{circle size }m&4&5&6\\ \hline
 \text{lower bound}&7&15&13
 \end{array}
 \qquad
 \begin{array}{c|ccc}
 \text{line size }m&4&5&6\\ \hline
 \text{lower bound}&12&18&15.
 \end{array}
\]
The entry $m=4$ for circles need not be excluded; all larger circles and
all lines of size at least four are excluded.  Thus maximal blocks have size
three or four, and every four-block is a circle.

Let $L=L_3$, let $Q=C_4$, and let $R=C_3$.  Triple ownership gives
\begin{equation}\label{eq:n7-count}
 35=L+4Q+R,\qquad C(P)=Q+R=35-L-3Q.
\end{equation}
Since $21=L_2+3L$, Melchior gives $L_2\geq3$ and hence $L\leq6$.
The complements of two different four-circles intersect in at most one
point: the two circle supports intersect in at most two.  Therefore no pair
of labels occurs in two complementary triples, and $3Q\leq21$, so $Q\leq7$.

The assumption $C(P)\leq10$ and \eqref{eq:n7-count} imply
$L+3Q\geq25$.  If $Q\leq6$, the left side is at most $6+18=24$.
Consequently $Q=7$.  The seven complements contain
$7\binom32=21$ distinct pairs, hence every pair exactly once: they form a
Steiner triple system.  This contradicts \cref{lem:fano-seven-wall}.
\end{proof}

\subsection{Eight points}

\begin{lemma}[Even-arrangement slack]\label{lem:even-arrangement-slack}
Let $\mathcal A$ be a non-pencil arrangement of an even number of distinct
real projective lines.  Its Melchior slack
\[
 \kappa:=t_2-3-\sum_{r\geq4}(r-3)t_r
\]
cannot equal $1$.
\end{lemma}

\begin{proof}
Let $V,E,F$ be the numbers of vertices, edges, and faces of the projective
cell decomposition.  Euler's formula and the incidence count
$E=\sum_r r t_r$ give
\[
 \kappa=\sum_{\text{faces }G}(|G|-3).
\]
Because the number of lines is even, the sign of the product of homogeneous
linear equations for the lines is well-defined on the real projective plane
and two-colours the faces.  Every edge borders one face of each colour.
If $F_i$ is the number of faces of colour $i$, the sum of their boundary
lengths is $E$, and their excess is
\[
 \kappa_i=E-3F_i\geq0.
\]
Thus $\kappa_0\equiv\kappa_1\pmod3$ and
$\kappa=\kappa_0+\kappa_1\neq1$.
\end{proof}

\begin{proposition}\label{prop:n8}
One has $f(8)=17$.
\end{proposition}

\begin{proof}
Construction~\ref{con:eight-upper} gives the upper bound.  Suppose that an
admissible eight-point set has $C(P)\leq16$, and let $k$ be the maximum
number of selected points on a proper circle.  For $k=7,6,5$,
\cref{lem:rich-circle-pencil} gives, respectively,
\[
 C(P)\geq19,\qquad22,\qquad19.
\]
If $k=3$, every noncollinear triple has a different owner.  If a largest
line has $m\geq4$ points and $r=8-m$, the number of collinear triples is at
most
\[
 \binom m3+\binom r3+\binom r2\leq35.
\]
If all lines have size at most three, their triples use disjoint triples of
pairs and there are at most $\floor{28/3}=9$ of them.  In either event
$C(P)\geq21$.  The case $k=8$ is inadmissible, so only $k=4$ remains.

A line of size $m=5,6,7$ gives, by
\cref{lem:rich-line-pencil}, respectively $24,27,21$ circles.  Hence every
maximal generalized block has size three or four.  Put
\[
 b=B_4=C_4+L_4,\qquad
 L=L_3+L_4,\qquad h=L_4.
\]
Triple ownership, the summed pivot row, the original Melchior row, the
added-centre row, and the circle count are
\begin{align}
 B_3+4b&=56,                                         \label{eq:n8-triples}\\
 K:=\sum_{p\in P}\sigma_p&=144-12b\geq0,             \label{eq:n8-K}\\
 3L+4h&\leq25,                                       \label{eq:n8-global-line}\\
 9L+7h&\leq200-12b,                                  \label{eq:n8-added}\\
 C(P)&=56-3b-L.                                      \label{eq:n8-circle-count}
\end{align}
Indeed, \eqref{eq:n8-added} is
$9L_3+16L_4\leq56+K$ rewritten in terms of $L,h$.
Equations \eqref{eq:n8-K}, \eqref{eq:n8-global-line}, and
\eqref{eq:n8-circle-count} imply
\begin{equation}\label{eq:n8-b-range}
 b\in\{11,12\}.
\end{equation}

Suppose first that $b=11$.  Then $K=12$ and
\eqref{eq:n8-circle-count} gives $L\geq7$.  The two line inequalities force
\[
 L=7,\qquad h=0.
\]
At $p$, write $r_p=d_4(p)$.  Ownership of triples through $p$ and the
definition of $\sigma_p$ give
\[
 d_3(p)+3r_p=21,\qquad \sigma_p=18-3r_p.
\]
The local added-centre slack is
\[
 \kappa_p:=7+\sigma_p-3\ell_3(p).
\]
It is nonnegative and congruent to $1$ modulo $3$, hence at least one.
On the other hand,
\[
 \sum_{p\in P}\kappa_p=56+K-9L_3=56+12-63=5,
\]
which cannot be the sum of eight positive integers.

Now let $b=12$.  Then $K=0$, so $\sigma_p=0$ for every $p$, and
\begin{equation}\label{eq:n8-local-kappa}
 \kappa_p=7-3\ell_3(p)-4\ell_4(p)\geq0.
\end{equation}
The added-centre configuration has eight points and no line containing more
than four of them.  Its dual is therefore a non-pencil arrangement of eight
real projective lines, and $\kappa_p$ is its Melchior slack.  By
\cref{lem:even-arrangement-slack},
\begin{equation}\label{eq:n8-kappa-not-one}
 \kappa_p\neq1.
\end{equation}
Equations \eqref{eq:n8-local-kappa}--\eqref{eq:n8-kappa-not-one} imply
$\ell_3(p)\leq1$: if $\ell_4(p)=0$, the only additional nonnegative case is
$\ell_3(p)=2$, which gives slack one; if $\ell_4(p)\geq1$, then necessarily
$\ell_4(p)=1$ and already $\ell_3(p)\leq1$.  Hence
$3L_3\leq8$ and $L_3\leq2$.

Equation \eqref{eq:n8-circle-count} gives $L\geq4$.  Combining this with
\eqref{eq:n8-global-line}, \eqref{eq:n8-added}, and $L_3\leq2$ leaves
\begin{equation}\label{eq:n8-final-line-row}
 L_3=2,\qquad L_4=2,\qquad C(P)=16.
\end{equation}
Moreover,
\[
 \sum_p\kappa_p=56-9L_3-16L_4=6.
\]
The four possible local pairs and slacks are
\[
\begin{array}{c|cccc}
(\ell_3(p),\ell_4(p))&(0,0)&(1,0)&(0,1)&(1,1)\\ \hline
\kappa_p&7&4&3&0.
\end{array}
\]
Since $\sum_p\ell_4(p)=4L_4=8$, every point lies on one four-line; the slack
sum then forces two points of type $(0,1)$ and six of type $(1,1)$.  Thus the
two four-lines are disjoint and partition $P$.  Each of the two three-lines
contains three points, so two of them lie on the same four-line.  The
three-line and that four-line would then be the same geometric line, contrary
to maximality of their different supports.  This final contradiction proves
$C(P)\geq17$.
\end{proof}

\subsection{Nine points}

\begin{proposition}\label{prop:n9}
One has \(f(9)=25\).
\end{proposition}

\begin{proof}
Construction~\ref{con:circle-centre}, with four antipodal boundary pairs,
gives \(f(9)\leq1+\binom82-4=25\).  For the reverse inequality, suppose
that an admissible nine-point set \(P\) satisfies
\begin{equation}\label{eq:n9-counterassumption}
 C(P)\leq24,
\end{equation}
and let \(k\) be the largest number of selected points on a proper circle.
Then \(3\leq k\leq8\).

\smallskip
\noindent\emph{The branches \(k\geq6\).}
For \(k=8,7,6\), respectively, the rich-circle bound
\eqref{eq:rich-circle-pencil} gives
\[
 C(P)\geq25,\quad34,\quad28.
\]
Thus none of these branches is compatible with
\eqref{eq:n9-counterassumption}.

\smallskip
\noindent\emph{The branch \(k=5\).}
Fix a five-point circle \(\Gamma\) and put
\(X=P\setminus\Gamma\), so \(|X|=4\).  First,
\cref{lem:rich-line-pencil} excludes every line of size at least five:
for line sizes \(m=5,6,7,8\), its lower bound is, respectively,
\[
 28,\quad36,\quad39,\quad28.
\]
Hence all lines below have size at most four.

For a proper circle other than \(\Gamma\), let \(c_{r,t}\) count those
having \(r\) points on \(\Gamma\) and \(t\) points in \(X\).  Necessarily
\(r\leq2\).  Let \(\lambda_{r,t}\) denote the analogous number of rich
lines and set
\begin{equation}\label{eq:n9-R}
 R:=\sum_t t\lambda_{2,t}.
\end{equation}
At each outsider, the chords of \(\Gamma\) passing through it form a
matching on five vertices.  Therefore
\begin{equation}\label{eq:n9-R-cap}
 R\leq8.
\end{equation}

For \(x\in X\), let \(\mathcal S_x\) consist of the proper circles through
\(x\) and two points of \(\Gamma\).  Exactly the chord-lines through \(x\)
fail to yield such circles, whence
\[
 \sum_{x\in X}|\mathcal S_x|=40-R.
\]
For an outsider pair \(x,y\), the \(\Gamma\)-pairs represented in
\(\mathcal S_x\cap\mathcal S_y\) form a matching, so this intersection has
size at most two.  No circle contains two points of \(\Gamma\) and all four
outsiders, because \(k=5\).  Thus inclusion--exclusion stops after outsider
triples.  If \(c_2\) counts all circles with exactly two
\(\Gamma\)-points, then
\begin{equation}\label{eq:n9-four-outsider-union}
 c_2\geq40-R-12+c_{2,3}=28-R+c_{2,3}.
\end{equation}
Let \(c_{\leq1}\) count circles other than \(\Gamma\) with at most one
\(\Gamma\)-point.  Since \(C(P)=1+c_2+c_{\leq1}\),
\eqref{eq:n9-counterassumption} gives the four-outsider master inequality
\begin{equation}\label{eq:n9-four-outsider-master}
 R\geq5+c_{2,3}+c_{\leq1}.
\end{equation}

We now use only the four possible incidence types of \(X\).

\emph{No three outsiders are collinear and the four are not concyclic.}
The four outsider triples determine four distinct proper circles.  They
are counted by \(c_{2,3}+c_{\leq1}\), so
\eqref{eq:n9-four-outsider-master} gives \(R\geq9\), contradicting
\eqref{eq:n9-R-cap}.

\emph{All four outsiders are collinear.}
Put \(a=c_{1,2}\) and \(b=c_{2,2}\).  Counting the \(30\) triples with one
\(\Gamma\)-point and two outsiders gives
\begin{equation}\label{eq:n9-collinear-four-count}
 a+2b=30.
\end{equation}
There is no line contribution, since every outsider pair already lies on
the four-outsider line; nor can a proper circle contain three of its
collinear points.  The pair-intersection bound above gives \(b\leq12\),
and hence \(a\geq6\).  Moreover, the exact membership count for the
\(\mathcal S_x\) gives \(c_2=40-R-b\).  Therefore
\[
 C(P)=1+(40-R-b)+a
     =26-R+\frac32a\geq27,
\]
again a contradiction.

\emph{Exactly three outsiders are collinear.}
Let \(\lambda\) be their line and let \(s\in\{0,1\}\) record whether
\(\lambda\) contains a point of \(\Gamma\).  Put
\[
 a:=c_{1,3},\qquad h:=c_{2,3},\qquad
 \ell:=\lambda_{1,2}.
\]
The other three outsider triples determine three distinct proper circles.
Thus \(h+c_{\leq1}\geq3\).  Combining this with
\eqref{eq:n9-R-cap}, \eqref{eq:n9-four-outsider-union}, and
\eqref{eq:n9-four-outsider-master} first gives
\(R=8\) and \(h+c_{\leq1}=3\).  If
\[
 I:=\sum_{\{x,y\}\subset X}|\mathcal S_x\cap\mathcal S_y|
   =c_{2,2}+3h,
\]
then \(I\leq12\), while exact inclusion--exclusion gives
\[
 C(P)=1+40-R-I+h+c_{\leq1}=36-I.
\]
Thus \eqref{eq:n9-counterassumption} forces \(I=12\) and equality
throughout:
\begin{equation}\label{eq:n9-three-collinear-equality}
 R=8,\qquad C(P)=24,\qquad
 h+c_{\leq1}=3,\qquad c_{2,2}+3h=12.
\end{equation}
In particular those three outsider-triple circles exhaust
\(c_{\leq1}+h\), so \(c_{1,2}=0\).  The last equality says that each
outsider pair lies on exactly two circles with two \(\Gamma\)-points, and
their two \(\Gamma\)-pairs are disjoint.  There can be no line of type
\((2,2)\): its \(\Gamma\)-pair would have to be disjoint from both circle
pairs, which is impossible on five points.

Counting again the triples with one \(\Gamma\)-point and two outsiders,
and then applying the global line row, gives
\begin{align}
 \ell+3a+3s&=6,                         \label{eq:n9-three-collinear-triples}\\
 3\ell+4s&\leq6.                        \label{eq:n9-three-collinear-lines}
\end{align}
Indeed, the term already covered by the two-\(\Gamma\) circles is
\(2(c_{2,2}+3h)=24\); while the rich lines are the eight chord incidences,
\(\lambda\), and the \(\ell\) lines of type \((1,2)\).
Equations \eqref{eq:n9-three-collinear-triples}--%
\eqref{eq:n9-three-collinear-lines} yield
\[
 \ell=0,\qquad a+s=2.
\]
For each outsider pair, the two disjoint \(\Gamma\)-pairs in
\eqref{eq:n9-three-collinear-equality} leave one \(\Gamma\)-point.  Its
triple with the outsider pair must be owned by one of the \(a\)
three-outsider circles or, when \(s=1\), by \(\lambda\).
The pair-shadows of two distinct three-subsets of a four-set cover only
five of its six pairs.  One pair is therefore left without an owner, a
contradiction.

\emph{The four outsiders are concyclic.}
Let \(\Omega\) be their circle.  It contains \(\delta\in\{0,1\}\) selected
points of \(\Gamma\).  Put
\[
 a:=c_{1,2},\quad b:=c_{2,2},\quad
 u:=\lambda_{2,1},\quad v:=\lambda_{2,2},\quad
 \ell:=\lambda_{1,2}.
\]
Apart from \(\Gamma,\Omega\), every rich circle or line has type
\((2,1),(1,2)\), or \((2,2)\): a distinct carrier meets each of
\(\Gamma,\Omega\) in at most two points.  Consequently
\begin{align}
 R&=u+2v,                                                   \label{eq:n9-concyclic-R}\\
 C(P)&=42-R-b+a.                                           \label{eq:n9-concyclic-C}
\end{align}
Write \(b=12-e\), where \(e\geq0\).  From
\eqref{eq:n9-counterassumption}, \eqref{eq:n9-R-cap}, and
\eqref{eq:n9-concyclic-C},
\begin{equation}\label{eq:n9-concyclic-master}
 R\geq6+e+a,\qquad e+a\leq2.
\end{equation}
The same \(30\)-triple count and the global line row now read
\begin{align}
 \ell+2v&=6-a+2e-6\delta,                 \label{eq:n9-concyclic-triples}\\
 3R+3\ell+v&\leq33.                       \label{eq:n9-concyclic-lines}
\end{align}

If \(\delta=0\), all \(b+v\) generalized blocks of type \((2,2)\)
are cross-blocks between five points of \(\Gamma\) and four points of
\(\Omega\).  By \cref{lem:radical-axis-cross-block},
\(b+v\leq12\), hence \(v\leq e\).  But
\eqref{eq:n9-concyclic-master}--\eqref{eq:n9-concyclic-lines} give
\[
 33\geq3R+3\ell+v\geq36+9e-5v,
\]
or \(5v\geq3+9e\), impossible when \(v\leq e\).

If \(\delta=1\), remove the common selected point of \(\Gamma\) and
\(\Omega\).  A type-\((2,2)\) block cannot use that point: a circle would
share three points with \(\Omega\), while a line would contain three
collinear points of \(\Omega\).  The \(4+4\) case of
\cref{lem:radical-axis-cross-block} therefore gives
\[
 b+v\leq10,\qquad v\leq e-2.
\]
Together with \eqref{eq:n9-concyclic-master}, this forces
\[
 e=2,\qquad a=v=0,\qquad R=8.
\]
Equation~\eqref{eq:n9-concyclic-triples} gives \(\ell=4\), and
\eqref{eq:n9-concyclic-lines} becomes \(36\leq33\), the final
contradiction in the \(k=5\) branch.

\smallskip
\noindent\emph{The branches \(k\leq4\).}
The same rich-line calculation excludes lines of size at least five.
All maximal generalized blocks therefore have size three or four.  Put
\[
 B:=B_4=C_4+L_4,\qquad L:=L_3+L_4.
\]
Triple ownership and the circle objective give
\begin{equation}\label{eq:n9-four-block-count}
 C(P)=84-3B-L.
\end{equation}
The global line row is \(3L_3+7L_4\leq33\), so \(L\leq11\).
Under \eqref{eq:n9-counterassumption},
\eqref{eq:n9-four-block-count} then forces \(B\geq17\).

Fix \(p\in P\) and invert the other eight points at \(p\).  They are
noncollinear by \cref{lem:inversion-dictionary}, and no four are collinear:
such a line would invert to a five-point proper circle or a five-point
line through \(p\).  A three-point image line is exactly a generalized
four-block of \(P\) through \(p\).  Hence
\cref{lem:orchard-eight} gives \(d_4(p)\leq7\).  Summing over all pivots,
\[
 4B=\sum_{p\in P}d_4(p)\leq9\cdot7=63,
\]
so \(B\leq15\), contrary to \(B\geq17\).

Every possible value of \(k\) has now been excluded under
\eqref{eq:n9-counterassumption}.  Thus \(C(P)\geq25\), completing the
proof.
\end{proof}

\begin{theorem}[The small values]\label{thm:small-values}
For the V1 problem,
\[
 \bigl(f(4),f(5),f(6),f(7),f(8),f(9)\bigr)
   =(3,5,8,11,17,25).
\]
\end{theorem}

\begin{proof}
Combine \cref{prop:n4-n5,prop:n6,prop:n7,prop:n8,prop:n9}.
\end{proof}

\section{The finite cases \texorpdfstring{$n=10,11,12$}{n=10,11,12}}
\label{sec:finite-n10-n12}

Throughout this section $C=C(P)$ and
\[
 F(n):=1+\binom{n-1}{2}-\floor{\frac{n-1}{2}}.
\]
Thus $F(10)=33$, $F(11)=41$, and $F(12)=51$.  We prove the
corresponding lower bounds.  The upper bounds are already supplied by
\cref{con:circle-centre}.  All blocks, degrees, and slack variables have
the meaning fixed in \cref{sec:foundations-small}.

\subsection{A deletion-free outer router}

The following lemma removes every largest-circle size except five and six
for eleven and twelve points.  It is useful to record it before the local
geometry, because it also supplies the block-size guards used later.

\begin{lemma}[Outer router]\label{lem:n11-n12-outer-router}
Let $n\in\{11,12\}$ and suppose $C\le F(n)-1$.  If $m$ is the maximum
number of selected points on a proper circle, then $m\in\{5,6\}$.  Moreover,
for $n=11$ every line has at most five selected points, and for $n=12$
every line has at most six.
\end{lemma}

\begin{proof}
The pencil bounds \cref{lem:rich-line-pencil,lem:rich-circle-pencil} give
\begin{equation}
\begin{array}{c|c|c}
n& R_n(m),\ m=7,8,\ldots,n-1&
       \min_{s\ge s_0}L_n(s)\\ \hline
11&55,61,61,41&45\quad(s_0=6),\\
12&61,73,85,76,51&55\quad(s_0=7).
\end{array}\label{eq:outer-rich-values}
\end{equation}
Here $R_n$ and $L_n$ denote the right sides of the rich-circle and
rich-line bounds.  Every entry is at least $F(n)$, proving the asserted
caps and excluding $m\ge7$.

If $m\le3$, the triple partition and the global line row give, respectively,
\[
 C\ge \binom{11}{3}-\frac56\left(\binom{11}{2}-3\right)
 =\frac{365}{3},
\qquad
 C\ge \binom{12}{3}-\frac{10}{9}\left(\binom{12}{2}-3\right)=150.
\]
The constants $5/6$ and $10/9$ are the maxima of
$\binom{s}{3}/(\binom{s}{2}+s-3)$ over the allowed line sizes.

It remains to exclude $m=4$.  Fix a maximal four-circle $\Gamma$ and put
\[
 T'=\sum_{B\ne\Gamma}\binom{|B|}{3}=\binom n3-4.
\]
Let $O$ be the summed Lenchner row $3C_3\ge3n$, let $G$ denote the
left side of \eqref{eq:global-line-row}, and let $K$ denote the summed
Kelly--Moser row.  At $n=11$ it says $3B_3\ge55$.  The fixed combination
\[
 1+\frac14T'+\frac18O-\frac5{24}G+\frac18K
\]
has coefficient one on every proper circle and coefficients
$0,-11/24,0$ on lines of sizes $3,4,5$.  Hence
\[
 C\ge1+\frac{161}{4}+\frac{33}{8}-\frac{5\cdot52}{24}+\frac{55}{8}
 =\frac{497}{12}>41.
\]
At $n=12$ replace $K$ by the summed pivot row
$3B_3-5B_5-12B_6\ge36$.  The combination
\[
 1+\frac14T'+\frac2{15}O-\frac15G+\frac7{60}P
\]
has line coefficients $0,-2/5,-29/60,0$ for sizes $3,4,5,6$ and
coefficient one on every proper circle.  Therefore
\[
 C\ge1+\frac{216}{4}+\frac{72}{15}-\frac{63}{5}+\frac{7\cdot36}{60}
 =\frac{257}{5}>51.
\]
This excludes $m=4$ and completes the router.
\end{proof}

\subsection{Ten points}

\begin{lemma}[The maximum-four branch at ten points]
\label{lem:n10-gamma4}
If $|P|=10$, no proper circle contains more than four selected points, and
$C\le32$, then a contradiction follows.
\end{lemma}

\begin{proof}
A line of size $s\ge6$ gives, by \cref{lem:rich-line-pencil}, respectively
$42,54,52,36$ circles for $s=6,7,8,9$.  Thus a five-block is a line and
all other blocks have size three or four.  Put
$b=B_4$, $g=L_5$, $t=L_3$, and $h=L_4$.  The triple partition, circle count,
and global Melchior row are
\[
 B_3+4b+10g=120,\qquad
 C=120-3b-10g-t-h,\qquad
 3t+7h+12g\le42.
\]
At every pivot, Kelly--Moser and the congruence
$36=d_3+3d_4+6d_5$ give $d_3\ge6$, hence $B_3\ge20$.  It follows first that
$g=0$, then $b=25$ and $t+h\ge13$.  The summed added-centre row
\[
 9L_3+16L_4+25L_5\le90+3B_3-30-5B_5
\]
excludes $t+h=14$ and forces $h=0$ when $t+h=13$.  Therefore the only
possible aggregate profile is
\begin{equation}\label{eq:n10-g4-profile}
(B_3,B_4;L_3,L_4,L_5;C_3,C_4)
 =(20,25;13,0,0;7,25),
\end{equation}
and every pivot has $(d_3,d_4)=(6,10)$.

Invert at any pivot.  The nine-point image satisfies the hypotheses of
\cref{lem:real-nine-link}; its ordinary graph is two triangles and three
isolated points.  Let $\mathcal H$ be the $20$-edge three-uniform
hypergraph of generalized three-blocks.  The link of every vertex of
$\mathcal H$ is two triangles.  If $abc\in\mathcal H$, the edge $bc$ lies
in a unique link triangle, say with fourth vertex $d$; the same link
argument at $b$ forces all four faces of $abcd$.  Thus the $20$ hyperedges
partition into five tetrahedron boundaries.  At most one face of each
boundary is a line, since two faces share two points and $L_4=0$.  Hence
$L_3\le5$, contrary to \eqref{eq:n10-g4-profile}.
\end{proof}

We first isolate the local parity input and the three geometric kernels.
This keeps the final numerical router transparent: it will inspect only
integer moments, while every equality face is discharged by a named lemma.

\begin{lemma}[Ten-point local parity]\label{lem:n10-local-parity}
Assume that $|P|=10$, that no proper circle contains more than five selected
points, and that $C\le32$.  Put $r_p=d_5(p)$ and $q_p=\ell_4(p)$.  Then
\begin{equation}\label{eq:n10-local-pair}
 36=d_3(p)+3d_4(p)+6d_5(p),\qquad d_3(p)\in\{6,9\}.
\end{equation}
If $\ell_5(p)=0$, the largest possible value of $\ell_3(p)$ is
\begin{equation}
\begin{array}{c|c|rrrr}
q_p&\text{type}&r_p=0&r_p=1&r_p=2&r_p=3\\ \hline
0&d_3=6&4&3&2&3\\
0&d_3=9&4&4&3&4\\
1&d_3=6&2&1&2&1\\
1&d_3=9&3&2&3&2\\
2&d_3=6&0&1&0&-\\
2&d_3=9&1&2&1&0\\
3&d_3=6&0&-&-&-\\
3&d_3=9&1&0&-&0.
\end{array}\label{eq:n10-full-local-table}
\end{equation}
A dash means that the indicated local state is impossible.  If
\(\ell_5(p)=1\), the maximum is at least one smaller than the entry in the
\(q_p=0\) row with the same values of \(d_3(p)\) and \(r_p\).
\end{lemma}

\begin{proof}
The rich-line pencil bound excludes lines of size at least six.  Inversion
at $p$ gives a noncollinear nine-point set whose ordinary lines are the
three-blocks through $p$.  Kelly--Moser and the pair partition give
$d_3(p)\ge6$ and $d_3(p)\equiv0\pmod3$.  Deletion of $p$ is nondegenerate:
otherwise the corresponding rich line or rich circle pencil already has at
least $33$ proper circles.  Hence \cref{thm:small-values} gives
$c_3(p)\le C-f(9)\le7$, while the ray bound gives $\ell_3(p)\le4$.
This proves \eqref{eq:n10-local-pair}.

Set $\sigma_p=d_3(p)-3-d_5(p)$.  Added-centre Melchior says that
\begin{equation}\label{eq:n10-kappa-local}
 \kappa_p:=9+\sigma_p-3\ell_3(p)-4\ell_4(p)-5\ell_5(p)\ge0.
\end{equation}
Moreover $\kappa_p\ne1$.  Indeed, after inversion and restoration of the
centre, dualize to an arrangement of ten real projective lines.  The sign
of the product of their defining forms two-colours the cells.  If
$\kappa_+$ and $\kappa_-$ are the sums of $|F|-3$ over the two colours,
edge counting gives $\kappa_+\equiv\kappa_-\pmod3$; total slack one is
therefore impossible.  Substituting $d_3=6,9$ in
\eqref{eq:n10-kappa-local}, and taking the largest integer
$\ell_3\le4$ for which $\kappa_p\ge0$ and $\kappa_p\ne1$, gives every entry
of \eqref{eq:n10-full-local-table} directly.  Setting \(\ell_5=1\) subtracts
five from the numerator in \eqref{eq:n10-kappa-local}; for
\(r=0,1,2,3\), in both the low and high rows, the same integer maximization
lowers \(\ell_3\) by at least one.  This proves the final assertion.
\end{proof}

\begin{lemma}[Punctured-pentagon transversal]\label{lem:n10-punctured-pentagon}
Let five five-subsets of a ten-set have pairwise intersections at most two
and point-degree multiset $(3^5,2^5)$.  At a degree-three point $p$, at most
two four-subsets through $p$ can meet every five-subset in at most two
points and meet one another in at most two points.
\end{lemma}

\begin{proof}
The degree row gives
$\sum_p\binom{d_5(p)}2=5\binom32+5\binom22=20
=2\binom52$.  Since every block pair meets in at most two points, every
pair therefore meets in exactly two.
Every five-subset contains two degree-two and three degree-three points:
if these numbers are $e,t$, then $e+t=5$ and $e+2t=8$.  On the five block
labels, the supports of the degree-two points form a two-regular graph.
It has no double edge, since a doubled $ij$ edge would exhaust the
intersection of blocks $i,j$, whereas their six required occurrences among
the five degree-three supports could then not be supplied.  Thus it is a
five-cycle.  Complements of the degree-three supports form the complementary
five-cycle, because each label pair occurs altogether twice.  This proves
the incidence pattern uniquely, without a support catalogue.

Relabel one degree-three point as $p=124$ and the other nine points as
\[
 A=12,\ B=134,\ C=245,\quad
 D=135,\ Q=15,\ R=23,\ E=235,\ S=34,\ T=45.
\]
The three blocks through $p$ have pairwise vertices $A,B,C$ and private
pairs
\begin{equation}\label{eq:n10-puncture-private}
 \{D,Q\},\qquad\{R,E\},\qquad\{S,T\}.
\end{equation}
A four-set through $p$ cannot use $A,B$, or $C$: such a vertex fills the
available intersection slot on two incident five-blocks and leaves room
for at most one further point.  It must therefore choose one point from
each pair in \eqref{eq:n10-puncture-private}.  Membership in the two
five-blocks $3$ and $5$ is
\[
 D:35,\ Q:5,\ R:3,\ E:35,\ S:3,\ T:5.
\]
Forbidding three chosen points in either omitted block leaves exactly
\begin{equation}\label{eq:n10-puncture-four}
 DRT,\qquad QRS,\qquad QRT,\qquad QES.
\end{equation}
Two residual triples are compatible precisely when they share at most one
point.  The compatible pairs are
\[
 DRT\!-!QRS,\qquad DRT\!-!QES,\qquad QES\!-!QRT,
\]
the edges of a four-vertex path.  Its clique number is two, proving the
claim.
\end{proof}

\begin{lemma}[Four-pentagon obstruction]\label{lem:n10-four-pentagon}
Assume the hypotheses of \cref{lem:n10-local-parity}, and suppose that
exactly four maximal generalized five-blocks occur.  Then the configuration
is impossible.
\end{lemma}

\begin{proof}
Write the blocks as $B_1,\ldots,B_4$, put $r(p)=d_5(p)$, and set
\[
 S=\sum_p\binom{r(p)}2
   =\sum_{i<j}|B_i\cap B_j|.
\]
We have $\sum_pr(p)=20$, $S\le12$, and convexity gives $S\ge10$.
Moreover
\begin{equation}\label{eq:n10-k4-square}
 \sum_p(r(p)-2)^2=2S-20.
\end{equation}
Thus the only degree rows are
\begin{equation}\label{eq:n10-k4-degree-rows}
 \begin{array}{c|c}
 S&\{r(p):p\in P\}\\ \hline
 10&(2^{10})\\
 11&(1,2^8,3)\\
 12&(1^2,2^6,3^2).
 \end{array}
\end{equation}
For each five-block and four-block $D$ introduce centred vectors
\[
 w_i=\mathbf1_{B_i}-\frac12\mathbf1,
 \qquad v_D=\mathbf1_D-\frac25\mathbf1.
\]
Maximality gives the coefficientwise inequalities
\begin{equation}\label{eq:n10-k4-centred}
 \langle v_D,w_i\rangle=|D\cap B_i|-2\le0.
\end{equation}

If $S=10$, put $e_{ij}=2-|B_i\cap B_j|$.  The five actual
pair-incidences of $B_i$, against the maximum six across its three block
partners, give $\sum_{j\ne i}e_{ij}=1$, so the two
deficient pairs form a perfect matching, say $12,34$.  Since every point
has degree two, $\sum_iw_i=0$; equality after summing
\eqref{eq:n10-k4-centred} forces $|D\cap B_i|=2$ for all $i$.  A four-block
through the unique $12$ point must contain the unique $34$ point: the three
remaining membership edges need degree sequence $(1,1,2,2)$, which is
impossible without the edge $34$.  Its last two points form $13|24$ or
$14|23$.  Different four-blocks cannot reuse either of them, whence
$d_4(12)\le4$; but \eqref{eq:n10-local-pair} requires $d_4(12)=5$ or $6$.

If $S=11$, there are a degree-three point $T$, a degree-one point $U$, and
one deficient block pair.  Comparing the incidence and intersection row at
each label shows that, after relabelling, the ten supports are
\begin{equation}\label{eq:n10-k4-S11-supports}
 T=123,\quad U=1,\quad 12,13,23,14,24,24,34,34.
\end{equation}
Here $\sum_iw_i=\mathbf1_T-\mathbf1_U$.  Summing
\eqref{eq:n10-k4-centred} shows that every four-block through $T$ also
contains $U$.  Its last two points must be a matching between the two
$24$ points and the two $34$ points, so $d_4(T)\le2$, whereas
\eqref{eq:n10-local-pair} requires $3$ or $4$.

It remains that $S=12$.  Comparing the label rows now forces, up to
relabelling, the support dictionary
\begin{equation}\label{eq:n10-k4-S12-supports}
 \begin{array}{c|c}
 T,T'&123,124\\
 U_1,U_2&1,2\\
 a,b,c,d&13,14,23,24\\
 e,f&34,34.
 \end{array}
\end{equation}
Indeed, if $t_i,s_i$ count triple and singleton supports at label $i$, the
five-incidence row gives $s_i=t_i-1$; hence the two triples meet in two
labels and the singletons carry those labels.  Also
\begin{equation}\label{eq:n10-k4-centred-sum}
 \sum_iw_i=\mathbf1_T+\mathbf1_{T'}-
            \mathbf1_{U_1}-\mathbf1_{U_2}.
\end{equation}
A four-block through $T$ cannot contain $T'$, and
\eqref{eq:n10-k4-centred-sum} forces at least one singleton.  Splitting
according as it contains both, only $U_1$, or only $U_2$, the four
intersection caps leave precisely
\begin{equation}\label{eq:n10-k4-local-six}
 \begin{array}{lll}
 TU_1U_2e,&TU_1de,&TU_2be,\\
 TU_1U_2f,&TU_1df,&TU_2bf.
 \end{array}
\end{equation}
Within each row the first entry conflicts with the other two, and the three
same-column pairs between rows conflict.  Thus at most three coexist;
equality has one of the two row-switched orientations.  The same holds at
$T'$.  Both points must consequently have $(d_3,d_4,d_5)=(9,3,3)$, and
the orientations must be opposite, since equal orientations would produce
two four-blocks sharing $U_1,U_2$ and one of $e,f$.  Interchanging $e,f$
if necessary, the following six blocks are forced:
\begin{equation}\label{eq:n10-k4-forced-six}
 \begin{array}{lll}
 TU_1U_2e,&TU_1df,&TU_2bf,\\
 T'U_1U_2f,&T'U_1ce,&T'U_2ae.
 \end{array}
\end{equation}

Every remaining four-block avoids $T,T'$.  Its intersections with the four
five-blocks are bounded by
\begin{equation}\label{eq:n10-k4-residual-ineq}
 \begin{array}{ll}
 |D\cap\{U_1,a,b\}|\le2,&|D\cap\{U_2,c,d\}|\le2,\\
 |D\cap\{a,c,e,f\}|\le2,&|D\cap\{b,d,e,f\}|\le2,
 \end{array}
\end{equation}
and compatibility with \eqref{eq:n10-k4-forced-six} forbids the six triples
\begin{equation}\label{eq:n10-k4-forbidden-six}
 U_1U_2e,\ U_1df,\ U_2bf,\ U_1U_2f,\ U_1ce,\ U_2ae.
\end{equation}
According as both, one, or neither singleton is used,
\eqref{eq:n10-k4-residual-ineq}--\eqref{eq:n10-k4-forbidden-six} give
exactly
\begin{equation}\label{eq:n10-k4-residual-list}
\begin{array}{c|l}
U_1,U_2&U_1U_2ac,\ U_1U_2ad,\ U_1U_2bc,\ U_1U_2bd\\
U_1&U_1acd,\ U_1ade,\ U_1bcd,\ U_1bcf\\
U_2&U_2abc,\ U_2abd,\ U_2adf,\ U_2bce\\
\text{neither}&abcd.
\end{array}
\end{equation}
For example, with both singletons the forbidden triples exclude $e,f$ and
the first two inequalities choose one of $a,b$ and one of $c,d$; with only
$U_1$, choosing $a$ leaves $cd$ or $de$, choosing $b$ leaves $cd$ or $cf$;
the other row is symmetric.  With no singleton, the last two inequalities
exclude $e,f$, proving that the list is exhaustive.

Name the thirteen entries of \eqref{eq:n10-k4-residual-list} rowwise
$A,\ldots,M$.  The seven sets
\begin{equation}\label{eq:n10-k4-clique-cover}
 \{A,B\},\{C,D\},\{E,F\},\{G,H\},
 \{I,L\},\{J,K\},\{M\}
\end{equation}
are conflict cliques, so at most seven residual blocks coexist.  Equality
uses $M$ and therefore forces successively $F,H,L,K,A,D$; it is uniquely
\begin{equation}\label{eq:n10-k4-residual-equality}
 U_1U_2ac,\ U_1U_2bd,\ U_1ade,\ U_1bcf,
 \ U_2adf,\ U_2bce,\ abcd.
\end{equation}

Let $\alpha$ be the number of points with $d_3=9$.  Summing $d_3$ and the
triple partition gives $B_3=20+\alpha$ and
$B_4=15-\alpha/4$.  Since $T,T'$ are high, integrality leaves
$\alpha=4$ or $8$.  If $\alpha=4$, then $B_4=14$, but the preceding caps
give at most $6+7=13$.  Hence $\alpha=8$, $B_4=13$, and all blocks in
\eqref{eq:n10-k4-forced-six} and
\eqref{eq:n10-k4-residual-equality} occur.

Invert at $T=123$.  The three incident five-blocks are the sides of a
nondegenerate triangle with vertices
$A=T'$, $B=a$, and $C=c$.  Use directed parameters
\begin{equation}\label{eq:n10-k4-coordinate-dictionary}
\begin{array}{c|c}
AB&A=0,\ B=1,\ U_1=u,\ b=v\\
AC&A=0,\ C=1,\ U_2=s,\ d=t\\
BC&B=0,\ C=1,\ e=x,\ f=y.
\end{array}
\end{equation}
The first row of \eqref{eq:n10-k4-forced-six} gives the collinearities
\begin{equation}\label{eq:n10-k4-collinearities}
 U_1U_2e,\qquad U_1df,\qquad U_2bf,
\end{equation}
while its second row gives the cyclic quadruples
\begin{equation}\label{eq:n10-k4-cyclic-data}
 AU_1U_2f,\qquad AU_1Ce,\qquad AU_2Be.
\end{equation}
The nonincident five-block and the last residual block are the proper
circles
\begin{equation}\label{eq:n10-k4-two-circles}
\{A,b,d,e,f\},\qquad\{B,b,C,d\}.
\end{equation}
Here \(u-s\ne0\): equality would put the finite, distinct point \(e\) at
the pole of the first collinearity \(U_1U_2e\), contrary to the chosen
finite side-parameter chart.  No affine transformation of the ambient
Euclidean plane is being used here; the parameters in
\eqref{eq:n10-k4-coordinate-dictionary} are directed coordinates on the
three actual sides of the inverted triangle.
Menelaus applied to \eqref{eq:n10-k4-collinearities} gives
\begin{equation}\label{eq:n10-k4-menelaus}
 x=\frac{s(u-1)}{u-s},\qquad
 y=\frac{t(u-1)}{u-t}=\frac{s(v-1)}{v-s}.
\end{equation}
Put $R=|AC|^2/|AB|^2>0$.  Subtracting the circle equations in
\eqref{eq:n10-k4-cyclic-data} from the first circle in
\eqref{eq:n10-k4-two-circles}, at $e$ and $f$, yields
\begin{align}
 (v-u)(1-x)+R(t-1)x&=0,\notag\\
 (v-1)(1-x)+R(t-s)x&=0,\label{eq:n10-k4-power-three}\\
 (v-u)(1-y)+R(t-s)y&=0.\notag
\end{align}
The first two equations together with \eqref{eq:n10-k4-menelaus} give
$R=-u/s$; power of $A$ to the second circle in
\eqref{eq:n10-k4-two-circles} then gives $v=Rt=-ut/s$.
Substitution in \eqref{eq:n10-k4-menelaus} and
\eqref{eq:n10-k4-power-three} produces
\begin{align}
 s^2u+st-s-tu&=0,\notag\\
 2stu-st-su+t^2-tu&=0,\label{eq:n10-k4-three-polynomials}\\
 s^2t-s^2+st^2-stu+t^2u-t^2&=0.\notag
\end{align}
All cleared factors are nonzero because the named side-points are distinct.
In particular $t-u\ne0$ follows from the second collinearity,
$s^2+tu=-s(v-s)\ne0$ from the third, and $s^2-t\ne0$ below would otherwise
force $t=1$.  The first equation in
\eqref{eq:n10-k4-three-polynomials} therefore gives
\[
 u=\frac{s(1-t)}{s^2-t}.
\]
The remaining two become
\begin{align}
 (s+t)(s^2t-3st+s+t^2)&=0,\notag\\
 (s+t)(s^3t-s^3+s^2t-2st^2+t^2)&=0.
 \label{eq:n10-k4-final-two}
\end{align}
If $s+t=0$, then $v=u$, contrary to distinctness.  Subtracting
$(1-2s)$ times the first remaining factor from the second gives
$s(s-1)^2(3t-1)=0$.  Hence $t=1/3$, and the first factor becomes
$3s^2+1=0$, impossible over $\R$.  This closes the last row of
\eqref{eq:n10-k4-degree-rows}.
\end{proof}

\begin{lemma}[Three-pentagon rich-line endpoint]
\label{lem:n10-three-pentagon-endpoint}
Assume the hypotheses of \cref{lem:n10-local-parity} and the aggregate row
\[
 C=32,\quad B_5=3,\quad
 |\{p:d_3(p)=9\}|=2,\quad L_5=0,
 \quad L_4=h\in\{1,2,3\},\quad L_3=10-h.
\]
Then the configuration is impossible.
\end{lemma}

\begin{proof}
Put $r(p)=d_5(p)$ and $q(p)=\ell_4(p)$.  On each four-line $D$,
maximality gives
\begin{equation}\label{eq:n10-k3-line-degree}
 \sum_{p\in D}r(p)=\sum_{i=1}^3|D\cap B_i|\le6.
\end{equation}
The degree sum is $15$.  Convexity and the pair-intersection cap give
$5\le\sum_p\binom{r(p)}2\le6$, and the only degree profiles are
\begin{equation}\label{eq:n10-k3-profiles}
 P_1=(1^5,2^5),\qquad
 P_2=(1^6,2^3,3),\qquad
 P_3=(0,1^3,2^6).
\end{equation}
Indeed, the lower moment is the distribution $(1^5,2^5)$; increasing it
by one either replaces two \(1\)'s by \(0,2\), giving \(P_3\), or replaces
two \(2\)'s by \(1,3\), giving \(P_2\).  The $q=0$ rows of
\eqref{eq:n10-full-local-table} give all-low
triple-line capacities
\begin{equation}\label{eq:n10-k3-baselines}
 25,\qquad27,\qquad25,
\end{equation}
and each of the two high designations raises a baseline by at most one.

For $h=1$, \eqref{eq:n10-k3-line-degree} forces at least two degree-one
points on the line in $P_1,P_2$; each loses two units from its $q=0$
capacity, even if high.  The total capacities are therefore at most
$23,25<3L_3=27$.  In $P_3$ the line contains either the degree-zero point
or two degree-one points.  Making the zero point high reduces its line loss
by one but gives no baseline gain, so the total is at most $25<27$.

For $h=2$, two lines need four degree-one incidences in $P_1,P_2$.  If
their intersection is degree one, its $q=2$ row pays two units and the two
remaining incidences pay two each; otherwise the loss is larger.  Starting
from the high-upgraded maxima $27,29$ leaves at most $21,23<3L_3=24$.
For $P_3$ the cheapest possibility is that the lines meet at the
degree-zero point; \eqref{eq:n10-full-local-table} still leaves at most
$23<24$.

Suppose $h=3$.  The twelve line incidences have only three possible
intersection topologies: a common point of all three lines, two distinct
pairwise intersections, or three distinct pairwise intersections.  Indeed,
without a common point their union has size $12-j\le10$, where $j$ is the
number of selected pairwise intersections, so $j=2$ or $3$.
Measure loss from \eqref{eq:n10-k3-baselines}, already crediting the two
possible high-point discounts.  A degree-one point on one or two lines
costs two, a degree-two point costs zero on one line and two on two, and a
degree-zero point costs two on one line and four on two or three.  Directly
from these four costs and \eqref{eq:n10-k3-line-degree}, the losses are
\begin{equation}\label{eq:n10-k3-loss-table}
\begin{array}{c|ccc}
&\text{common point}&\text{two intersections}&\text{triangle}\\ \hline
P_1&10&8&6\\
P_2&14&12&10\\
P_3&8&6&6.
\end{array}
\end{equation}
Here is a derivation rather than a case catalogue.  In the common-point
topology that point cannot have degree two; the first two rows cost
$3+4\cdot2-1=10$ and $3+5\cdot2+2-1=14$, while the third costs at least
$8$.  With two intersections, let $x$ be the number of degree-one double
points and $z$ indicate that the degree-zero point is double.  Before the
two discounts the first two losses are at least
$4+2(5-x)$ and $4+2(6-x)+2$; for $P_3$ they are
$4z+2(2-z)+2(1-z)+2(3-x)-2=10-2x\ge6$.
In the triangle topology let $x$ count degree-one vertices and let
$\epsilon$ indicate whether the off-line point has degree one or zero.
The first row costs at least $6+2(5-x-\epsilon)-2\ge6$; the analogous
$P_2$ expression has the extra degree-three deficit and is at least $10$.
For $P_3$ the exact loss is
\begin{equation}\label{eq:n10-k3-P3-loss}
 12-2x-2\epsilon.
\end{equation}
It is at least six except when the three triangle vertices are exactly the
three degree-one points and the off-line point is the degree-zero point.
In that sole exception each side contains two degree-two points and both
high discounts are attained; the loss is four.  All other entries of
\eqref{eq:n10-k3-loss-table} leave capacity below $3L_3=21$.

Let $p$ be the off-line degree-zero point in the equality pattern.  It is
low, $\ell_3(p)=4$, and \eqref{eq:n10-local-pair} gives
\[
 (d_3(p),d_4(p),d_5(p))=(6,10,0).
\]
After inversion at $p$, the other nine points have line profile
$(t_2,t_3,t_4)=(6,10,0)$.  By \cref{lem:real-nine-link}, their six ordinary
lines form $C_3\sqcup C_3$.  The four original three-lines through $p$
would give four pairwise disjoint ordinary edges, but
$\nu(C_3\sqcup C_3)=2$.  This final contradiction closes the endpoint.
\end{proof}

\begin{lemma}[Disjoint golden-link endpoint]\label{lem:n10-golden-link-endpoint}
Assume the hypotheses of \cref{lem:n10-local-parity} and the aggregate row
\begin{equation}\label{eq:n10-k2-endpoint-row}
 C=32,\quad B_3=B_4=20,\quad B_5=2,\quad
 L_3=10,\quad L_4=L_5=0,
 \quad d_3(p)=6\ \text{for every }p.
\end{equation}
Then the configuration is impossible.
\end{lemma}

\begin{proof}
Let the two five-blocks meet in $m=0,1$, or $2$ selected points.  Their
degree profiles are, respectively,
\begin{equation}\label{eq:n10-k2-degree-profiles}
 (1^{10}),\qquad(0,1^8,2),\qquad(0^2,1^6,2^2).
\end{equation}
The $q=0$, low row of \eqref{eq:n10-full-local-table} gives total
triple-line capacity exactly $30$ in each profile, equal to
$\sum_p\ell_3(p)=3L_3$.  Thus every point attains its local maximum.  If
$m>0$, choose a degree-zero point $p$.  Then
$(d_3,d_4,d_5,\ell_3)(p)=(6,10,0,4)$.  Inversion at $p$ gives the
nine-point line profile $(t_2,t_3,t_4)=(6,10,0)$, whose ordinary graph is
$C_3\sqcup C_3$ by \cref{lem:real-nine-link}.  The four original lines
through $p$ would be a four-edge matching in this graph, contradicting
$\nu(C_3\sqcup C_3)=2$.  Hence $m=0$.

The blocks are therefore disjoint proper circles; call their point sets
$A,B$.  Every four-block has two points on each, since three on a base
circle would repeat an owned triple.  Fixing an $A$-pair, at most two
four-blocks contain it, because their $B$-pairs must be disjoint.  There
are ten $A$-pairs and exactly twenty four-blocks, so equality holds: every
$A$-pair has exactly two $B$ partners.  Symmetrically every $B$-pair has
exactly two $A$ partners.

Use the quadratic lift and \cref{lem:circle-lift-trace}.  No lift of an
\(A\)-point lies in \(\Pi_B\), since \(\Pi_B\) would then cut the quadric
in a generalized block containing the five points of \(B\) and that sixth
selected point; symmetrically no \(B\)-point lies in \(\Pi_A\).
Thus the base planes are distinct and meet in a projective line \(\ell\).
A $2A+2B$ block plane
cuts them in two chords whose chord lines meet $\ell$ at the same centre;
conversely, an $A$-chord and a $B$-chord with the same centre span a block
plane.  That plane cannot acquire a fifth selected point, since it is
different from the two base planes and \eqref{eq:n10-k2-endpoint-row}
contains only the two five-blocks $A,B$.  At one centre the chords on either
base form a matching.  If it carries $\alpha$ $A$-chords and $\beta$
$B$-chords, all $\alpha\beta$ cross-pairs are four-blocks; the exact partner
degree two of every chord gives
\begin{equation}\label{eq:n10-k2-centre-degree}
 \alpha=\beta=2.
\end{equation}
Thus the edges of each $K_5$ split into five two-edge matchings, in common
centre correspondence.  By \cref{lem:pentagon-golden}(P2), after labelling
the $B$-vertices $0,\ldots,4$ their matching rows are
\begin{equation}\label{eq:n10-k2-near-factor}
 M_0=13|24,\ M_1=02|34,\ M_2=03|14,\
 M_3=04|12,\ M_4=01|23.
\end{equation}

Choose $p\in A$.  Row $0$ is the unique $A$-matching omitting $p$.  For
$i=1,\ldots,4$, let $q_i$ be the other endpoint of the $A$-chord through
$p$ whose centre corresponds to $M_i$.  After inversion at $p$, the four
$q_i$ lie on the image line of $A$, and $q_i$ is the intersection of the
two $B$-chords in $M_i$.  Normalize the five $B$-points as
\begin{equation}\label{eq:n10-k2-conic-frame}
 b_0=(1,0,0),\ b_1=(0,1,0),\ b_2=(0,0,1),\
 b_3=(1,1,1),\ b_4=(a,b,1).
\end{equation}
The four relevant centres are
\begin{equation}\label{eq:n10-k2-centres}
 q_1=(a-b,0,1-b),\quad q_2=(a,1,1),\quad
 q_3=(0,b,1),\quad q_4=(-1,-1,0).
\end{equation}
Their collinearity gives, exactly as in
\cref{lem:pentagon-golden}(P3),
\begin{equation}\label{eq:n10-k2-golden-wall}
 a=1-b=b^2,\qquad b^2+b-1=0.
\end{equation}

The eight four-blocks through $p$ give, at $q_i$, the two three-lines
prescribed by $M_i$.  These cover every $q_iB$ pair except $q_ib_i$ and
every internal $B$-pair except the two edges of $M_0$.  Hence the six
ordinary edges in the inverted link are exactly
\begin{equation}\label{eq:n10-k2-two-paths}
 q_ib_i\ (1\le i\le4),\qquad b_1b_3,\qquad b_2b_4,
\end{equation}
namely the two paths
\[
 q_1-b_1-b_3-q_3,\qquad q_2-b_2-b_4-q_4.
\]
In the disjoint profile every point has $r=1$; the local table and
$\sum_p\ell_3(p)=30$ force $\ell_3(p)=3$ everywhere.  Thus the three
original lines through $p$ give a concurrent three-edge matching in these
two paths.  Such a matching takes both end edges of one path and one edge
of the other.  The automorphism $\rho=(1234)$ of
\eqref{eq:n10-k2-near-factor} interchanges the paths, while $\rho^2$
reverses both.  Consequently there are only two orbits, represented by
\begin{equation}\label{eq:n10-k2-two-matchings}
 \{q_1b_1,q_3b_3,q_2b_2\},\qquad
 \{q_1b_1,q_3b_3,b_2b_4\}.
\end{equation}
Writing a line as the cross product of its endpoint vectors, their
concurrency determinants are
\begin{align}
 \det(q_1\!\times b_1,q_3\!\times b_3,q_2\!\times b_2)
   &=-(a-1)(ab-a+b),\notag\\
 \det(q_1\!\times b_1,q_3\!\times b_3,b_2\!\times b_4)
   &=a^2b-a^2+ab-b^2.\label{eq:n10-k2-determinants}
\end{align}
On the golden wall \eqref{eq:n10-k2-golden-wall} these reduce to
$4-6b$ and $11b-7$.  Their possible roots $2/3$ and $7/11$ do not satisfy
$b^2+b-1=0$.  Neither matching in \eqref{eq:n10-k2-two-matchings} is
concurrent, the required contradiction.
\end{proof}

\begin{lemma}[Ten-point five-block kernel]\label{lem:n10-five-block-kernel}
Assume $|P|=10$, the largest proper circle has size five, and $C\le32$.
Put
\[
x=B_3,\quad b=B_4,\quad k=B_5,\quad
L=L_3+L_4+L_5,\quad h=L_4,\quad g=L_5.
\]
Then no such configuration exists.
\end{lemma}

\begin{proof}
The rich-line bound excludes lines of size at least six.  Triple ownership,
the circle objective, and the two Melchior rows give
\begin{align}
x+4b+10k&=120,\label{eq:n10-g5-T}\\
C&=30+\frac34x-\frac32k-L,\label{eq:n10-g5-C}\\
15x-34k+28h+64g&\le36C-840.\label{eq:n10-g5-AC}
\end{align}
By \cref{lem:n10-local-parity}, $d_3(p)\in\{6,9\}$.  Let $\alpha$ be
the number of high points, those with $d_3=9$.  Summing $d_3$ gives
\begin{equation}\label{eq:n10-high-count}
 x=20+\alpha.
\end{equation}
For $r_p=d_5(p)$ we also have
\begin{equation}\label{eq:n10-five-moment}
 \sum_pr_p=5k,\qquad
 \sum_p\binom{r_p}{2}\le2\binom k2.
\end{equation}

We first remove the large values of $k$.  The centred-vector argument and
the real triangle-power obstruction in \cref{lem:pentagon-golden}(P1) give
$k\le5$.  If $k=5$, convexity and \eqref{eq:n10-five-moment} force
$(r_p)=(3^5,2^5)$ and equality in every pair intersection.  At a
degree-three point, \cref{lem:n10-punctured-pentagon} gives $d_4\le2$,
whereas \eqref{eq:n10-local-pair} gives $d_4=3$ or $4$.  Thus $k\ne5$;
\cref{lem:n10-four-pentagon} removes $k=4$.  We have proved
\begin{equation}\label{eq:n10-kcap}
 1\le k\le3.
\end{equation}

Values $C\le27$ disappear before any further routing.  Indeed,
$d_3(p)=c_3(p)+\ell_3(p)\le(C-25)+4=C-21$, while $d_3(p)\ge6$.
At $C=27$ equality would force $\ell_3(p)=4$ at all ten points, so
$3L_3=40$, impossible.

For $C\ge28$, equations \eqref{eq:n10-g5-T}--\eqref{eq:n10-g5-AC} and
\eqref{eq:n10-high-count} become the exact scalar router
\begin{align}
 \alpha+2k&\equiv0\pmod4,\notag\\
 L&=45-C+\frac{3\alpha-6k}{4},\label{eq:n10-scalar-router}\\
 15\alpha-34k+28h+64g&\le36C-1140.\notag
\end{align}
For $k=1,2,3$, the least nonnegative $\alpha$ allowed by the congruence is
$2,0,2$, respectively; increasing it by four raises the last left side by
$60$.  Comparing with the five right sides
$-132,-96,-60,-24,12$ for $C=28,\ldots,32$ therefore gives, without a
configuration search, the complete list
\begin{equation}\label{eq:n10-scalar-faces}
\begin{array}{c|c|c}
C&(k,\alpha,L)&(h,g)\\ \hline
28,29&\text{none}&\text{--}\\
30&(2,0,12),(3,2,12)&(0,0)\\
31&(2,0,11),(3,2,11)&(0,0),(1,0)\\
32&(1,2,13),(2,4,13),(3,6,13)&(0,0)\\
  &(2,0,10)&(0,0),(1,0),(2,0),(0,1)\\
  &(3,2,10)&(0,0),(1,0),(2,0),(3,0),(0,1).
\end{array}
\end{equation}
For example, at $C=30$ the two base values in the last inequality are
$-68,-72$, leaving only $8,12$ units and hence $h=g=0$; the other
line-state columns follow by the same displayed costs $28,64$.

It remains to apply the local table, and we record every capacity comparison.
For two five-blocks the three degree profiles
\eqref{eq:n10-k2-degree-profiles} have all-low capacity $30$; for three
five-blocks the profiles \eqref{eq:n10-k3-profiles} have all-low capacities
$25,27,25$.  Each high point adds at most one.  Thus at $C=30$ the two
rows in \eqref{eq:n10-scalar-faces} have capacities at most $30,29$, below
$3L_3=36$.  At $C=31$, the three-block row has capacity at most $29$,
below both required totals $33,30$.  In the two-block row, $h=0$ gives
$30<33$; on a four-line $D$ one has $\sum_{p\in D}r_p\le4$, and the weak
loss $2-r_p$ sums to at least four, leaving at most $26<30$ when $h=1$.

At $C=32$, the three $L=13$ rows have capacities at most
\begin{equation}\label{eq:n10-C32-easy-capacities}
 37,\qquad34,\qquad33,
\end{equation}
against the required $39$.  For the first number, the five points of the
unique five-block contribute $3$ each, its outsiders $4$ each, and at most
two high points on the block add one; the other two numbers are
$30+4$ and $27+6$.

Consider $(k,\alpha,L)=(2,0,10)$.  A five-line loses at least one unit at
each of its points, leaving at most $25<27$.  For one or two four-lines,
sum the weak loss $2-r_p$ on each line; it is at least four because
$\sum_Dr_p\le4$.  This addition does not overcount at an intersection:
for $r=2,1,0$ the two-line losses from
\eqref{eq:n10-full-local-table} are respectively $2,2,4$, at least the
corresponding doubled weak weights $0,2,4$.  Hence capacity is at most
$30-4h<30-3h=3L_3$.  Only $h=g=0$ survives, and it is excluded by
\cref{lem:n10-golden-link-endpoint}.

Finally take $(k,\alpha,L)=(3,2,10)$.  The no-rich-line state has capacity
at most $29<30$, and a five-line leaves at most $29-5<27$.  The three
remaining states $h=1,2,3$, $g=0$, are precisely the hypotheses of
\cref{lem:n10-three-pentagon-endpoint}.  This exhausts
\eqref{eq:n10-scalar-faces} and completes the five-circle branch.
\end{proof}

\begin{lemma}[The ten-point six-circle certificate]
\label{lem:n10-gamma6}
If ten admissible points contain a proper circle with at least six selected
points, then $C\ge33$.
\end{lemma}

\begin{proof}
The rich-circle bound gives $46,45,33$ for occupancies $7,8,9$; the
rich-line bound gives $54,52,36$ for line occupancies $7,8,9$.  Thus a
counterexample has a maximal six-circle $\Gamma$, four outsiders, and no
block larger than six.

Let $L_{ij},C_{ij}$ count line and circle blocks containing $i$ points of
$\Gamma$ and $j$ outsiders.  We use
\begin{equation}
 \sum ijL_{ij}\le24,\qquad \sum C_{ij}\le32,\qquad
 \sum\left(\binom{i+j}{2}+i+j-3\right)L_{ij}\le42.\label{eq:n10-g6-rows}
\end{equation}
the U17 row
\[
 \sum_j\binom j2C_{2j}\le17,
\]
and, for $q=0,1,2,3$, the four exact triple rows
\[
 \sum_{i,j}\binom iq\binom j{3-q}(L_{ij}+C_{ij})
 =\binom6q\binom4{3-q}.
\]
Also $C_{60}=1$.  Multiply the three inequalities in
\eqref{eq:n10-g6-rows}, the U17 row, and the four equalities by
\[
 24,84,12,82;\qquad -21,1,-84,684,
\]
and add $-13764C_{60}$.  A block distinct from $\Gamma$ meets it in at most
two points, so there are exactly nine types.  Their coefficients are
\[
\begin{array}{c|rrrrrrrrr}
(i,j)&(0,3)&(0,4)&(1,2)&(1,3)&(1,4)&(2,1)&(2,2)&(2,3)&(2,4)\\ \hline
L_{ij}&15&0&85&138&162&0&14&21&0\\
C_{ij}&63&0&85&66&6&0&0&63&168.
\end{array}
\]
The selected circle has coefficient zero, so the combined left side is
nonnegative.  The combined right side is
\[
24\cdot24+84\cdot32+12\cdot42+82\cdot17
-21\cdot4+36-84\cdot60+684\cdot20-13764=-10,
\]
a contradiction.
\end{proof}

\begin{proposition}\label{prop:n10-value}
One has $f(10)=33$.
\end{proposition}

\begin{proof}
Assume $C\le32$ and split by the maximum proper-circle size.  The three
possibilities at most four, exactly five, and at least six are excluded by
\cref{lem:n10-gamma4,lem:n10-five-block-kernel,lem:n10-gamma6}.
Construction \ref{con:circle-centre} has $33$ circles.
\end{proof}

\subsection{Real link and circle-lift modules}

We next isolate the real geometry used more than once.  These statements
are independent of the endpoint arithmetic.

\begin{lemma}[The real nine-point link]\label{lem:real-nine-link}
Let $Q$ be nine distinct real projective points, no four collinear, with
exactly ten three-point lines.  Then it has six ordinary lines, whose graph
is
\[
 C_3\sqcup C_3\sqcup3K_1.
\]
Consequently ten such points, no four collinear, have at most twelve
three-point lines, and eleven such points have at most sixteen.
\end{lemma}

\begin{proof}
We first need the eight-point Orchard fact: eight noncollinear real
projective points, no four collinear, have at most seven three-point lines.
Indeed, eight such lines would leave four ordinary lines.  Local pair
counting makes the ordinary lines a perfect matching and every point lie on
three three-lines.  Label its four pairs $A,B,C,D$.  Pair ownership forces,
up to relabelling, the eight triples
\[
\begin{gathered}
A_0B_0C_0,\ A_1B_1C_1,\ A_0B_1D_0,\ A_1B_0D_1,\\
A_0C_1D_1,\ A_1C_0D_0,\ B_0C_1D_0,\ B_1C_0D_1.
\end{gathered}
\]
Normalize $A_0,A_1,B_0,B_1$ to a projective frame.  Writing the remaining
four points with parameters $u,v,w,z$, four of the incidences give
$vz=uw=vw=1$ and $u+u^{-1}=1$.  Thus $u^2-u+1=0$, impossible over $\R$.

For the nine-point set, pair counting gives $t_2=6$.  If $k_v$ and $e_v$
are the three-line and ordinary degrees at $v$, deleting $v$ and applying
the preceding fact gives $k_v\ge3$, while
\[
 8=2k_v+e_v.
\]
Hence three points have degree pair $(4,0)$ and six have $(3,2)$.  The
ordinary graph is either $C_6\sqcup3K_1$ or the graph in \eqref{eq:real9}.
We reject the first alternative directly.  Calling the isolated vertices
$U_6,U_7,U_8$, the ten three-lines have one of the two distributions
\[
 (a_3,a_2,a_1,a_0)=(1,0,9,0),\qquad(0,3,6,1),
\]
according to the number of $U$-vertices on a line.  In the first case the
nine $UVV$ lines force the three matchings
\[
 02,14,35;\qquad03,15,24;\qquad04,13,25.
\]
Sending the $U$-line to infinity and using two independent directions
reduces the last required parallelism to $2/\lambda=0$.  In the second
case a projective frame at $U_6,U_7,U_8$ reduces the last three incidences
to
\[
 \alpha\gamma=1,\qquad\beta\gamma=-1,\qquad
 \alpha\beta=1,
\]
and hence $\beta^2=-1$.  Both are impossible over $\R$, proving
\begin{equation}\label{eq:real9}
\operatorname{Ord}(Q)=C_3\sqcup C_3\sqcup3K_1.
\end{equation}

If ten points had thirteen three-lines, their ordinary-degree multiset
would be $(3,1^9)$, so the ordinary graph would be
$K_{1,3}\sqcup3K_2$.  Delete its degree-three vertex.  The remaining
ordinary graph is $C_6\sqcup3K_1$ and there remain ten three-lines,
contrary to \eqref{eq:real9}.  Fourteen or more three-lines leave fewer
than five ordinary pairs, contradicting the positive odd local ordinary
degrees.  For eleven points, seventeen three-lines leave an ordinary
$C_4$; deleting a cycle vertex gives the forbidden ten-point case.
Eighteen leave one ordinary edge although every ordinary degree is even.
This proves both ceilings.
\end{proof}

Use the circular lift
\[
 v(x,y)=[x^2+y^2:x:y:1]\in
 \mathcal Q:\ XW-Y^2-Z^2=0\subset\PP^3(\R).
\]
The real quadric $\mathcal Q$ contains no real projective line.  A line or
proper circle in the original plane is a plane section of $\mathcal Q$.

\begin{lemma}[Trace dictionary]\label{lem:circle-lift-trace}
Let a selected generalized block have lift plane $H_B$, and let $H$ be the
plane of a selected proper circle $\Gamma$.  Then:
\begin{enumerate}[label=\textup{(T\arabic*)}]
\item $H_B\cap H$ is one trace line, containing every $\Gamma$-chord and
every outsider-secant direction belonging to $B$.
\item For a fixed outsider pair, its host chords on $\Gamma$ form a
matching.  Two hosts determine their diagonal centre on the trace.
\item Equal direction points cannot occur on adjacent outsider edges; a
direction fibre is a matching.
\item The three diagonal points of a complete quadrangle on a nonsingular
real conic are distinct and noncollinear (indeed, self-polar).
\end{enumerate}
In particular, a block containing one point $\gamma_i\in\Gamma$ and three
outsiders cannot have all three outsider edges double-hosted by chords
avoiding $\gamma_i$.
\end{lemma}

\begin{proof}
Plane sections give (T1).  If two hosts shared $\gamma$, they would both
own the selected triple consisting of $\gamma$ and the outsider pair,
proving (T2).  Equal directions on $xy,xz$ would put the three real quadric
points $v_x,v_y,v_z$ on one real line, proving (T3).  In a conic frame
$[t^2:t:1]$, the three diagonal points are the poles of the three diagonal
sides, which proves (T4).  The final assertion would put all three diagonal
points and $\gamma_i$ on the trace of one block.
\end{proof}

\begin{lemma}[Pentagon, golden axis, and polarity]\label{lem:pentagon-golden}
Let $\mathcal F$ be a family of five-subsets of a ten-set, any two meeting
in at most two points.
\begin{enumerate}[label=\textup{(P\arabic*)}]
\item $|\mathcal F|\le6$.  At equality every point has degree three and
every pair of blocks meets in two points.  Up to relabelling the incidence
design is
\begin{equation}
 123,124,135,236,146,156,245,256,345,346.\label{eq:pentagon-design}
\end{equation}
Six such blocks cannot all be realized as generalized blocks of ten real
points.
\item Five two-edge matchings which partition $E(K_5)$ form, up to
relabelling, the near-one-factorization
\begin{equation}
 13|24,\ 02|34,\ 03|14,\ 04|12,\ 01|23.\label{eq:near-factor}
\end{equation}
\item Let five points lie on a real conic and let the centres in
\eqref{eq:near-factor} be collinear.  After a projective normalization the
modulus satisfies
\begin{equation}
 a=1-b=b^2,\qquad b^2+b-1=0.\label{eq:golden}
\end{equation}
The fifth centre is then on the same line.  The five opposite diagonal
lines concur at its pole; no selected chord passes through the pole, and
the tangent at an omitted point contains its opposite centre but no second
centre omitting that point.
\end{enumerate}
\end{lemma}

\begin{proof}
For the incidence vectors $x_F$ put $y_F=x_F-\frac12\mathbf1$.  Then
$\|y_F\|^2=5/2$ and $\langle y_F,y_G\rangle\le-1/2$.  The squared norm of
$\sum_Fy_F$ gives $|\mathcal F|\le6$; equality forces all intersections and
degrees stated in (P1).  Deleting one block, the degree-two points and the
complements of the degree-three points are complementary five-cycles,
which forces \eqref{eq:pentagon-design}.

We make the geometric part of (P1) explicit.  Invert a putative real
realization at the point labelled $123$.  Since equality makes every pair
of blocks meet in two selected points, at most one of the six generalized
blocks can be an original line.  The three blocks through $123$ have the
distinct second intersections
\[
 A=124,\qquad B=135,\qquad C=236.
\]
Their images are therefore three pairwise nonparallel lines.  They cannot
be concurrent, since their pairwise intersections $A,B,C$ are distinct;
the possible original line through the inversion centre causes no
exception.  Thus these lines form a nondegenerate triangle $ABC$.  Every
nonincident block becomes a proper circle: a proper circle missing the
pivot remains proper, while the sole possible original line missing the
pivot becomes a proper circle through the inversion centre.  Inversion is
injective away from the pivot, so the remaining nine selected points stay
distinct.  Write
\[
 X=146,\quad Y=156,\quad U=245,\quad V=256,
 \quad W=345,\quad Z=346.
\]
The three nonincident circles have supports
\[
 A,X,U,W,Z;\qquad B,Y,U,V,W;\qquad C,X,Y,V,Z.
\]

Use directed affine parameters on the triangle sides
\[
 \begin{array}{c|cccc}
 AB&A=0&B=1&X=x&Y=y\\
 AC&A=0&C=1&U=u&V=v\\
 BC&B=0&C=1&W=w&Z=z
 \end{array}
\]
and put $c^2=|AB|^2$, $b^2=|AC|^2$, and $a^2=|BC|^2$.
Directed power at the three vertices, applied to the displayed circles,
gives all six identities
\begin{align*}
 c^2(1-x)&=a^2wz,&
 b^2(1-u)&=a^2(1-w)(1-z),\\
 c^2y&=b^2uv,&
 a^2(1-w)&=b^2(1-u)(1-v),\\
 c^2xy&=b^2v,&
 c^2(1-x)(1-y)&=a^2z.
\end{align*}
The side lengths are positive, and distinctness from the side vertices says
that $x,y,u,v,w,z$ and all six complements
$1-x,1-y,1-u,1-v,1-w,1-z$ are nonzero.  Hence the indicated divisions are
legal.  Dividing the last-row identities by the corresponding identities
above, and comparing the second and fourth identities, gives
\[
 x=u^{-1},\qquad 1-y=w^{-1},\qquad 1-z=(1-v)^{-1}.
\]
Set $R=b^2/c^2$ and $S=a^2/c^2$.  The first and third identities now give
\[
 R=\frac{w-1}{wuv},\qquad
 S=\frac{(u-1)(v-1)}{uwv}.
\]
After division by $c^2$, the fourth identity is
\[
 S(1-w)=R(1-u)(1-v).
\]
Its two sides are respectively $-K$ and $K$, where
\[
 K=\frac{(u-1)(v-1)(w-1)}{uvw}\ne0.
\]
Thus $-K=K$ over the reals, a contradiction.  This proves the geometric
nonrealizability in (P1).

For (P2), fix the matching missing vertex $0$; its two edges force the next
unused edge at each of the other four vertices, and four successive unused
edge rows give \eqref{eq:near-factor}.  For (P3), put four conic points at
$[1:0:0],[0:1:0],[0:0:1],[1:1:1]$ and the fifth at $[a:b:1]$.
Collinearity of the first four matching centres gives
$a=1-b=b^2$.  Direct substitution places the fifth centre on their line.
Polarity sends that line to the common point of the five opposite sides.
La Hire's theorem then gives the tangent statement; a second centre on the
tangent would make one self-polar diagonal triangle collinear.
\end{proof}

\begin{lemma}[Five-conic trace rigidity]\label{lem:five-conic-rigidity}
Let $\Gamma=\{\gamma_0,\ldots,\gamma_4\}$ be five cyclically ordered points
on a nonsingular real conic.  Work in the trace-dictionary setting above
and assume that $\Gamma$ is a maximum selected generalized five-block, with
lift plane $\Pi$.  Thus every outsider lift lies outside $\Pi$, and no plane
section of $\mathcal Q$ contains more than five selected lifts.  For each
$i$, let $D_i$ be the diagonal triangle of
$\Gamma-\{\gamma_i\}$, and let $\delta_i$ be the vertex opposite the unique
side whose line meets the missing real arc at $\gamma_i$.
Then the following hold.
\begin{enumerate}
\item[\textup{(K2.1)}] A one-$\Gamma$ block cannot contain an all-double outsider triangle.
If exactly one edge is single and the other two are double, its host chord
is one of
\begin{equation}
 \gamma_{i-1}\gamma_{i+1},\qquad
 \gamma_{i-2}\gamma_{i+2}.\label{eq:cyclic-host}
\end{equation}
These five two-element chord classes partition the ten chords.  Consequently
a fixed single-hosted outsider pair occurs in at most one such
one-$\Gamma$ block.
\item[\textup{(K2.4)}] An all-double outsider $K_5$ has exactly five centres, one omitting
each colour; its trace is the golden axis.  The tangent separation in
\cref{lem:pentagon-golden} forbids extending it by a one-colour page with
two adjacent double edges omitting that colour.
\end{enumerate}
\end{lemma}

\begin{proof}
The all-double assertion in (K2.1) is the final clause of
\cref{lem:circle-lift-trace}.  For the remaining assertion, in the conic frame
$[t^2:t:1]$, the three diagonal sides of four cyclic points meet the missing
arc in exactly one side, and its opposite matching is \eqref{eq:cyclic-host}.
The five resulting two-chord classes are disjoint and exhaust the ten
chords.  A fixed single outsider edge has one host chord, which belongs to
exactly one of these classes; once its colour $i$ is fixed, triple ownership
forbids two blocks containing the same triple $\gamma_i xy$.  This proves
the fixed-single consequence.

For (K2.4), the ten edges of an all-double outsider $K_5$ have matching
direction fibres, so its trace contains at least five diagonal centres.
The trace is not a $\Gamma$-chord: otherwise its block plane would contain
the five outsiders and the two selected endpoints of that chord, contrary
to maximality of $\Gamma$.  A nonchord trace meets each of the ten
$\Gamma$-chords once.  Each diagonal centre uses two such chords, and two
distinct centres on the trace cannot reuse one, so there are at most five.
Equality partitions the ten chords into five two-edge matchings.
By \cref{lem:pentagon-golden}(P2), these form the near-one-factorization
\eqref{eq:near-factor}, with exactly one matching omitting each colour; by
(P3), their trace is the golden axis and $\gamma_i\delta_i$ is its polar
tangent.

Now consider a one-colour page of colour $i$ with two adjacent double
outsider edges.  Their direction points are distinct by the no-real-line
clause of the trace dictionary.  Triple ownership makes both host matchings
omit $i$, so both directions are vertices of $D_i$.  The direction of the
edge lying in the outsider $K_5$ is its unique centre $\delta_i$ omitting
$i$; the page trace is therefore the tangent $\gamma_i\delta_i$.  The
second direction would be a second vertex of $D_i$ on that tangent,
contrary to \cref{lem:pentagon-golden}(P3).
\end{proof}

\begin{lemma}[Six-marked conic and repetition events]
\label{lem:six-conic-events}
Let six cyclically ordered selected points lie on a proper circle $\Gamma$.
For an outsider edge $e=xy$, let $q_e\le3$ be the number of circles through
$x,y$ and a pair of points of $\Gamma$.
\begin{enumerate}[label=\textup{(S\arabic*)}]
\item There are at most four full signatures $q_e=3$; the support of each
is a matching.  For any four outsiders $Y$,
\begin{equation}\label{eq:U17}
 \sum_{e\in\binom Y2}q_e\le17.
\end{equation}
On any fixed trace line lie the centres of at most three full signatures.
\item Two disjoint edges carrying the same signature determine a unique
maximal generalized host of their four endpoints.  If such a host contains
$b$ outsiders and $R(H)$ repetition events, then
\begin{equation}\label{eq:conic-event-cap}
 R(H)\le \min\left\{2\binom b4,
 3\binom{\floor{b/2}}2\right\}.
\end{equation}
\item For five outsiders put $W=\sum q_e$.  A unique four-outsider host, or
a five-outsider host disjoint from $\Gamma$, gives $W\le26$; equality has
signature multiplicities $2,2,1,1$ or $2,2,2$, respectively.  A
five-outsider host meeting $\Gamma$ gives $W\le20$.
\item If $J$ counts incidences of outsiders with original lines carrying a
pair of $\Gamma$ points, then $J\le14$; on one fixed line or circle,
$J\le13$.
\end{enumerate}
\end{lemma}

\begin{proof}
A real projectivity of the conic preserves or reverses cyclic order.  The
only fixed-point-free involutions of six cyclic labels are
$R_3,S_1,S_3,S_5$, proving the cap four.  If all four occur, they generate
the rotation $R_1$; hence the six points form one orbit of a real
projectivity of order six.  A projective coordinate on the conic may
therefore normalize that orbit to $(0,1,3,\infty,-3,-1)$, which gives centres
\[
 (3,0,1),(3,3,-1),(-3,0,1),(-3,3,1),
\]
exactly three of which are collinear.  This also proves the trace-line cap.
Matching support follows from \cref{lem:circle-lift-trace}.  If all six
edges of a four-set were full, a repeated signature would put the three
perfect-matching centres on one trace, contradicting the noncollinearity of
the diagonal triangle.  Hence at least one edge has weight at most two,
which is \eqref{eq:U17}.

A repeated signature on $xy,zw$ puts their two lift secants and hence all
four lifts in one plane; its pullback is the unique line or circle through
the four endpoints.  Each four-set supports at most two events by U17, and
on one trace there are at most three signatures, each supported on a
matching of size at most $\floor{b/2}$.  These are the two bounds in
\eqref{eq:conic-event-cap}.

For five outsiders write $F$ for the number of full edges.  Then
$W\le20+F$.  In a unique four-host, at most two signatures repeat, so
$F\le2+2+1+1=6$.  On a five-host, at most three trace centres occur and
each supports two edges, again $F\le6$.  If the host contains a selected
$\gamma_i$, no full matching is possible because it would repeat a triple
$\gamma_i xy$, giving $W\le20$.  Finally, every outsider carries at most
three original chord lines.  Five degree-three outsiders would require five
signatures, and on a fixed host at most three may have degree three.  This
gives $J\le4\cdot3+2=14$ and $J\le3\cdot3+2\cdot2=13$.
\end{proof}

\subsection{The real grid and gallery modules}

\begin{lemma}[Real $3\times3$ grid]\label{lem:real-grid}
Let two triples of real projective lines form a $3\times3$ grid.  There are
at most two further lines meeting three grid points with one point in each
row and column.  If two exist, the grid is projectively the affine grid
$\{-1,0,1\}^2$.  No point outside this grid lies on three ordinary grid
secants.
\end{lemma}

\begin{proof}
Send the line through the two pencil centres to infinity.  A transversal is
the graph of an affine bijection between two real three-sets.  After one is
made the identity, every other is an affine automorphism of a real
three-set; a nonconstant affine map is determined by two values, and such a
set has at most the identity and its reflection.  Equality normalizes both
coordinate sets to $\{-1,0,1\}$.  In this frame the twelve ordinary grid
secants are
\[
 x\pm y=\pm1,\qquad x\pm2y=\pm1,\qquad 2x\pm y=\pm1.
\]
At infinity their six directions are distinct and each occurs exactly
twice, so no ideal point lies on three of them.  Now let a finite point
$(u,v)$ lie on three.  If $uv\ne0$, it can lie on at most one line from
each of the pairs with left sides $u\pm v$, $u\pm2v$, and $2u\pm v$.
It must therefore use one from every pair.  The first two equations imply
$|v|\in\{2,2/3\}$, while the first and third imply
$|v|\in\{1,3,1/3\}$, a contradiction.  If $uv=0$, the displayed equations
show directly that three incidences are possible only at a point of
$\{-1,0,1\}^2$.  Thus an external point, finite or ideal, lies on at most
two ordinary grid secants.
\end{proof}

\begin{lemma}[Gallery skeleton and its two completions]
\label{lem:gallery-completions}
There is no arrangement of twelve distinct real projective lines with
\begin{equation}
 (t_2,t_3,t_4)=(6,14,3)\label{eq:typeA-vector}
\end{equation}
and a distinguished line with complete word $3D+4T$.  Nor is there one
with
\begin{equation}
 (t_2,t_3,t_5)=(7,13,2)\label{eq:M4-vector}
\end{equation}
and the same distinguished word.
\end{lemma}

\begin{proof}
Melchior is sharp in either case, so the arrangement is simplicial and not
a near-pencil.  We first make the gallery entrance explicit.  Two ordinary
vertices cannot be adjacent on a line $\ell$: if they are
$\ell\cap A$ and $\ell\cap B$, send $\ell$ to infinity.  The two triangular
chambers along the intervening edge become opposite sectors bounded by
$A,B$ at $A\cap B$.  Any further arrangement line not through $A\cap B$
crosses one of those sectors (its direction cannot subdivide the chosen
edge), hence cuts a chamber.  Thus every further line passes through
$A\cap B$, the excluded near-pencil.  The three double vertices therefore
separate the four triple vertices into positive cyclic gaps of sizes
$2,1,1$, so the word is $TTDTDTD$.

Send the distinguished line to infinity.  Its four triple directions give
pairs $(a_i,b_i)$, $0\le i\le3$, and its three ordinary directions give
$c_1,c_2,c_3$.  At a consecutive flag
$T_{i-1}-D_i-T_i$, let $P^+,P^-$ be the other endpoints of the two edges of
$c_i$ incident with $D_i$.  The four triangular chambers at this flag show
that each $P^\pm$ lies on $c_i$, on one of the two lines through
$T_{i-1}$, and on one of the two lines through $T_i$.  On opposite sides
of a line germ at a transverse triple vertex the adjacent germs are the two
different remaining lines.  Hence $P^+,P^-$ define a perfect matching
between the two pairs.  The three flags form a path, so labels propagate
without a monodromy choice and give
\begin{equation}
 A_i=\{a_{i-1},a_i,c_i\},\qquad
 B_i=\{b_{i-1},b_i,c_i\}\quad(1\le i\le3),\label{eq:gallery-bases}
\end{equation}
together with the four rungs $\{a_i,b_i\}$.

After deleting their used pairs, split the eleven labels as
\[
 L=\{a_0,a_1,b_0,b_1,c_1,c_2\},\qquad
 R=\{a_2,a_3,b_2,b_3,c_3\}.
\]
The only residual within-side edges are
\begin{equation}\label{eq:gallery-H}
H=\{a_0b_1,a_0c_2,a_1b_0,b_0c_2,c_1c_2,
       a_2b_3,a_3b_2\}.
\end{equation}
The left part is a tree and the right part a matching; hence every residual
triangle uses exactly one edge of $H$.  If $e$ bases in \eqref{eq:gallery-bases}
are upgraded by high vertices and those vertices consume $h$ edges of $H$,
pair ownership gives the common budget
\begin{equation}\label{eq:gallery-budget}
 e+h\le13-t'_3,
\end{equation}
where $t'_3$ is the required number of triples after deleting the
distinguished line.

For \eqref{eq:typeA-vector}, deleting the distinguished line leaves
$(t'_2,t'_3,t'_4)=(7,10,3)$.  A nonupgrading quadruple is a four-clique of
the residual graph.  It has two labels on each side of the cut.  For the
right edge $a_2b_3$ its common left neighbours are
$\{a_0,b_0,b_1,c_1\}$, containing only the $H$-edge $a_0b_1$; for
$a_3b_2$ they are $\{a_0,a_1,b_0,c_1\}$, containing only $a_1b_0$.
Thus the only two nonupgrading quadruples are
\[
 \{a_0,a_2,b_1,b_3\},\qquad
 \{a_1,a_3,b_0,b_2\},
\]
and each consumes two $H$-edges.  If $e$ of the three quadruples upgrade
bases, the other $3-e$ consume at least $2(3-e)$ edges of $H$.
Equation \eqref{eq:gallery-budget} then gives
$e+2(3-e)\le3$, whence $e=3$ and no $H$-edge is consumed.

Only $A_1,B_1,A_3,B_3$ can be upgraded without using $H$.  Track exchange
$a\leftrightarrow b$ and gallery reversal act transitively on these four
bases, so assume $B_3$ is omitted and write the upgrades as
$A_1+x,B_1+y,A_3+z$.  Pair ownership gives
\[
 \begin{gathered}
 x\in\{a_3,b_2,b_3,c_3\},\qquad
 y\in\{a_2,a_3,b_3,c_3\},\\
 z\in\{a_0,b_0,b_1,c_1\},\qquad x\ne y,\\
 \neg(x\in\{a_3,c_3\}\land z\in\{a_0,c_1\}),\\
 \neg(y\in\{a_2,a_3,c_3\}\land z\in\{b_0,b_1,c_1\}).
 \end{gathered}
\]
Requiring only the two corner bundles $a_0b_1$ and $a_3b_2$ to remain
nonempty gives the complete four-state propagation certificate
\[
\begin{array}{c|c}
z& (x,y)\text{ still possible}\\ \hline
a_0&(b_2,a_2),(b_2,c_3)\\
b_0&(c_3,b_3)\\
b_1&\varnothing\\
c_1&(b_2,b_3).
\end{array}
\]
The three states
\[
 (b_2,a_2,a_0),\qquad (b_2,b_3,c_1),\qquad(c_3,b_3,b_0)
\]
empty, respectively, the bundles at $a_2b_3,c_1c_2,b_0c_2$.  Hence only
$(x,y,z)=(b_2,c_3,a_0)$ survives.  There all five left $H$-edges can be
completed only through the two poles $a_3,b_3$.  After the mandatory
triangles on the two right $H$-edges, their residual capacities are
\[
 (\floor{6/2}-1)+(\floor{7/2}-1)=4<5.
\]
This proves the first assertion by a fixed corner certificate.

For \eqref{eq:M4-vector}, deletion leaves
$(t'_2,t'_3,t'_5)=(8,9,2)$.  For a fivefold vertex $F$, let $e(F)$ be the
number of bases it upgrades and let $h(F)$ be the number of edges of $H$
among its pairs.  If $e(F)=0$, all ten pairs of $F$ lie in the residual
graph.  Each side of the cut in \eqref{eq:gallery-H} is triangle-free, so
$F$ has at most two labels on either side, contrary to $|F|=5$.  Hence
$e(F)\ge1$.

Two upgraded bases must be consecutive on the same rail:
\[
 A_1A_2,\quad A_2A_3,\quad B_1B_2,\quad B_2B_3.
\]
Indeed, these are the only base pairs whose union has size five; the
apparent pair $A_i,B_i$ repeats two rungs, every other pair has union at
least six, and three bases do not fit in five labels.  Thus $e(F)\le2$.

Any five-set has at least
$\binom22+\binom32=4$ within-side pairs.  A single middle base $A_2$ or
$B_2$ absorbs only one of them; every other within-side pair is residual by
pair ownership and hence belongs to $H$, so $h(F)\ge3$.  An end base absorbs
three.  Equality $e(F)+h(F)=2$ is possible only when its two extra labels
lie on the other side and form the unique available edge of $H$ in their
common-neighbour set.  For example, the common neighbours of $A_1$ in $R$
are $\{a_3,b_2,b_3,c_3\}$, whose unique $H$-edge is $a_3b_2$; rail exchange
and end reflection give the four end stars
\[
\begin{aligned}
 S_A^-&=A_1\cup\{a_3,b_2\},&
 S_B^-&=B_1\cup\{a_2,b_3\},\\
 S_A^+&=A_3\cup\{a_0,b_1\},&
 S_B^+&=B_3\cup\{a_1,b_0\}.
\end{aligned}
\]
If $F$ upgrades two bases, direct within-side pair counting gives
$h(A_1\cup A_2)=h(B_1\cup B_2)=2$ and
$h(A_2\cup A_3)=h(B_2\cup B_3)=0$.  Thus the first two unions have cost
four and the equality unions are
\[
 P_A=A_2\cup A_3,\qquad P_B=B_2\cup B_3.
\]
Consequently every fivefold vertex satisfies
\begin{equation}\label{eq:M4-cut-cost}
 e(F)+h(F)\ge2,
\end{equation}
with equality exactly for the four end stars and $P_A,P_B$.

For the two fivefold vertices, pair ownership makes their upgraded bases
and consumed $H$-edges disjoint.  Since $9-(6-e)=3+e$ new triples are
required and only $7-h$ edges of $H$ remain, \eqref{eq:gallery-budget} and
\eqref{eq:M4-cut-cost} force equality for both vertices.  Directly comparing
the six structural equality supports leaves the four-edge compatibility
graph
\[
 (P_A,S_B^-),\quad(P_B,S_A^-),\quad
 (S_A^-,S_B^-),\quad(S_A^+,S_B^+).
\]
For $(S_A^-,P_B)$, the $a_0c_2,c_1c_2$ bundles are empty and
$a_3b_2$ is consumed, leaving four bundles for six triples; the other mixed
pair is its reflection.  The two end-star pairs are mirror images, so take
$(S_A^-,S_B^-)$.  All five residual triangles are forced:
\[
 \{a_0,b_1,c_3\},\ \{a_0,b_3,c_2\},\
 \{a_1,b_0,c_3\},\ \{a_3,b_0,c_2\},\ \{c_1,c_2,c_3\}.
\]
The five lines of $S_A^-$ form a pencil $P$, those of $S_B^-$ a pencil
$Q$, and their common line is $c_1=PQ$.  Send it to infinity.  Then
$a_0,a_1,a_3,b_2$ are one parallel family, say $x=\text{constant}$, while
$b_0,b_1,a_2,b_3$ are a second, say $y=\text{constant}$.  The final forced
triple makes
$c_2:y=mx+\beta$ and $c_3:y=mx+\beta'$ parallel.  Writing $x_1,x_3$ for
the constants of $a_1,a_3$ and $y_2,y_0$ for those of $a_2,b_0$, the bases
and forced triples give
\[
 y_2=mx_1+\beta,\quad y_0=mx_3+\beta,\quad
 y_0=mx_1+\beta',\quad y_2=mx_3+\beta'.
\]
Hence $\beta'-\beta=y_0-y_2=y_2-y_0$, so $y_0=y_2$, contradicting the
distinctness of $b_0,a_2$.  This proves the second assertion.
\end{proof}

\subsection{Eleven points: the local K1--K4 package}

For a pivot $p$ in an eleven-point configuration whose blocks have size at
most five, inversion gives the exact link equation
\begin{equation}\label{eq:n11-link}
 d_3(p)+3d_4(p)+6d_5(p)=45.
\end{equation}
The next lemma records the real part of this link geometry.

\begin{lemma}[K3: inversion-link rigidity]\label{lem:n11-K3}
Let $P$ be a real eleven-point configuration in which every maximal
generalized block has size at most five, and let $p\in P$.  After
inversion at $p$, the other ten marked points form a real linear space with
line counts $(t_2,t_3,t_4)=(d_3,d_4,d_5)$, unique pair ownership, and two
distinct link lines meeting in at most one marked point.  Under
\eqref{eq:n11-link}, the following assertions hold.
\begin{enumerate}[label=\textup{(K3.\arabic*)}]
\item If $d_5=3$ and $d_3\in\{6,9\}$, two of the three four-point link
lines cannot be disjoint.
\item If $(d_3,d_4,d_5)=(9,4,4)$, the four four-point link lines form a
complete quadrangle.  Their four residual three-lines have, up to
relabelling, the private-pair graph ``triangle plus pendant.''  The four
points left on each corresponding five-block form a harmonic quadruple.
\item A five-block contains at most two pivots of type $(9,4,4)$.
\end{enumerate}
\end{lemma}

\begin{proof}
For (K3.1), three four-point link lines with one disjoint pair have
intersection graph a path.  Indeed, the disjoint pair uses eight of the ten
marked points.  The third four-point line contains at least two of those
eight and meets either base in at most one marked point, so it meets both
 exactly once.  Dualize.  Their private line families have
 sizes $3,2,3$, so the three cross-edge classes have sizes $6,9,6$.
 Denote the two connector lines by $x,y$ and the private families by
 $A,B,D$, of sizes $3,2,3$.  The three intersections of $x$ with the lines
 of $D$, and the three intersections of $y$ with the lines of $A$, are six
 forced ordinary vertices.  Every other intersection is represented by an
 edge of the complete tripartite graph on $A,B,D$.  The link equation gives
 $(t_2,t_3,t_4)=(6,7,3)$ or $(9,6,3)$.
 When $d_3=6$, seven triple vertices would have to decompose all three
classes equally, impossible.  When $d_3=9$, six triples exhaust the two
six-edge classes and leave three ordinary edges in the middle class.  The
two middle lines then give two edge-disjoint transversals of a real
$3\times3$ grid.  After one is diagonal and the coordinate set is
$(0,1,t)$, the other has determinant
$t^2-t+1>0$, a contradiction.

For (K3.2), four base lines have sixteen incidences on ten points and at
most six pairwise intersections.  Convexity forces degrees $2^6 1^4$:
the complete quadrangle with private points $u_i$ and pair-points $v_{ij}$.
A residual triple is either
\[
T_{ijk}=\{u_i,u_j,u_k\},\qquad
S_{ij}=\{u_i,u_j,v_{kl}\},\quad\{i,j,k,l\}=[4].
\]
No two residual lines repeat a private pair.  Hence at most one $T$-line
occurs.  If $T_{123}$ occurs, the other three lines are
$S_{14},S_{24},S_{34}$.  With no $T$-line, the four distinct $S$-indices
form a four-edge simple graph on four vertices, and pair ownership leaves
only a four-cycle or a triangle with a pendant edge.

Normalize the bases to $x=0,y=0,z=0,x+y+z=0$ and put
\[
u_1=(0,1,a),\ u_2=(1,0,b),\ u_3=(1,c,0),\\
u_4=(1,d,-1-d).
\]
Since every $u_i$ is private, distinct from all three pair-intersections on
its base, one has
\[
 a,b,c,d\notin\{0,-1\}.
\]
The seven relevant determinants are
\[
b+ac;\quad b-a,\ 1-ac,\ 1+d(1+a),\ b-c,\ 1+d+b,\ d-c.
\]
For the $T$-star, the three $S$-rows give
$d=c$, $b=-1-c$, and $ac=b$; the $T$-row then becomes $2b=0$, contrary to
nondegeneracy.  For a four-cycle the rows give $a=b=c=d$ and
$a^2+a+1=0$, which has no real root.  For the triangle-pendant
representative $S_{12},S_{13},S_{23},S_{14}$, they give $a=b=c$ and
 $a^2=1$.  The value $a=-1$ is incompatible with the $S_{14}$ row, while
 $a=1$ forces $d=-1/2$; substitution makes the other three candidate
 determinants nonzero.  Explicitly, at
 $(a,b,c,d)=(1,1,1,-1/2)$ the intended rows are
 \[
 S_{12}=S_{13}=S_{23}=S_{14}=0,
 \]
 while
 \[
 T_{123}=2,\quad T_{124}=1,\quad T_{134}=-1,\quad T_{234}=-2,\qquad
 S_{24}=\frac32,\quad S_{34}=-\frac32.
 \]
 Thus this is the unique real survivor.  Its
cross-ratio on every base is $-1$.  Inversion restricted to a generalized
block through $p$ is a projectivity onto its image line, so this harmonic
cross-ratio pulls back to the four points of the five-block other than $p$.

For (K3.3), identify the carrier line or circle of a five-block with
$\mathbb P^1(\mathbb R)$.  If three of its points $p,q,r$ had harmonic
complements, send the two remaining points to $0,\infty$ and scale so that
$p=1$.  The three harmonicity conditions say
\[
 p/q,\quad p/r,\quad q/r\in S:=\{-1,2,1/2\}.
\]
Since $S$ is closed under reciprocals, this puts the two distinct numbers
$q,r$ in $S$.  But
\[
 ((S/S)\setminus\{1\})=\{-2,-1/2,4,1/4\}
\]
is disjoint from $S$.  This proves the cap.
\end{proof}

\begin{lemma}[K1: rebased five-block moment]\label{lem:n11-K1}
Let $m=B_5$ and let $\Delta$ be a selected five-block.  Write $n_{ij}(\Delta)$
for the number of other $i$-blocks meeting $\Delta$ in $j$ points and put
\[
H_\Delta=n_{42}+3n_{52},\qquad
q_\Delta=n_{40}+3n_{50}.
\]
If $d_3(p)=6+3\omega_p$, $\Omega_\Delta=\sum_{p\in\Delta}\omega_p$,
$S_{55}=\sum_p\binom{d_5(p)}2$, and
$I_\omega=\sum_p\omega_pd_5(p)$, then
\begin{align}
H_\Delta&=57-N_4-2m-\Omega_\Delta+n_{52}+n_{40}+2n_{50},\label{eq:K1-local}\\
\sum_\Delta H_\Delta
&=(58-N_4-3m)m+2S_{55}-I_\omega+\sum_\Delta q_\Delta.\label{eq:K1-global}
\end{align}
For every $\Delta$, $H_\Delta\le30$.
\end{lemma}

\begin{proof}
Equation \eqref{eq:n11-link} is equivalent to
$d_4+2d_5=13-\omega$.  Sum it over the five points of $\Delta$ and group
blocks by intersection size; eliminating $n_{41}$ and $n_{51}$ gives
\eqref{eq:K1-local}.  Summing over $\Delta$, a pair of five-blocks meeting
in $j=0,1,2$ points contributes through the identities
$S_{55}=A_1+2A_2$ and $A_0+A_1+A_2=\binom m2$, which gives
\eqref{eq:K1-global}.  Finally, $H_\Delta$ counts hosts over the fifteen
outsider pairs.  Host chords form a matching on five points, so every pair
has at most two hosts and $H_\Delta\le30$.
\end{proof}

\begin{lemma}[K4: rich lines and the added centre]\label{lem:n11-K4}
In the $n=11$, maximum-circle-five branch with $C=40,L=14$, put $b=B_5$.
Then
\begin{equation}\label{eq:n11-K4-cut}
 7L_4+16L_5\le b+2.
\end{equation}
Consequently $L_5=0$, $L_4\in\{0,1\}$, and every row except
$b=7,8$ is impossible by local moments and Kelly--Moser.
\end{lemma}

\begin{proof}
For a maximal $m$-line with $q=11-m$ outsiders, let $x_s$ count proper
circle owners whose outsider support has size $s=1,2,3$.  For a fixed
outsider pair, its pairs on the $m$-line form a matching, and therefore
\[
x_1+2x_2+3x_3=\binom m2q,\qquad
Q:=x_2+3x_3\le\floor{m/2}\binom q2.
\]
The $D=x_1+x_2=\binom m2q-Q$ owners are distinct, since a proper circle
cannot contain two pairs from the same line.  Hence
\[
C\ge D\ge q\binom m2-\binom q2\floor{m/2}.
\]
For $m=6,7,8,9,10$ this is $45,66,72,68,45$, so all blocks have size at
most five.

We also print the five-block energy.  Deletion to $f(10)=33$ and the five
disjoint ray pairs give $d_3(p)\le7+5=12$.  If five four-point link lines
existed at a pivot, their twenty incidences on ten marked points would force
the complete-pentagon incidence skeleton.  Every further line contains at
most two of its pair-labelled points, whereas the link equation
\(d_3+3d_4+6d_5=45\) and $d_3\le12$ require a three-point link line.
Six four-point lines are excluded by convexity, since their twenty-four
incidences give at least eighteen pairwise meetings, more than
\(\binom62\).  Thus $d_5(p)\le4$ at every pivot, and
\[
5b=\sum_pd_5(p)\le44,
\]
so $b\le8$.

The restored eleven-point set at a pivot is noncollinear: a carrier line
through the restored centre would invert to an original line through the
pivot, while one avoiding it would invert to a generalized circle through
the pivot; either would contain all original points.  Thus restoration of
the inversion centre and Melchior give
\begin{equation}\label{eq:n11-K4-point}
3\ell_3(p)+\sum_{s\ge4}s\ell_s(p)\le10+\sigma_p.
\end{equation}
Here the block equations are
\[
N_3=17+2b,\quad N_4=37-3b,\quad
\sum_p\sigma_p=18+b.
\]
Summing \eqref{eq:n11-K4-point} gives \eqref{eq:n11-K4-cut}.  Since
$b\le8$, it gives $L_5=0$ and $L_4\le1$.

If $L_4=1$, let $\kappa_p$ be the local slack in
\eqref{eq:n11-K4-point}; then $\sum\kappa_p=b-5$.  The congruence
$d_5\equiv1-\ell_4-\kappa\pmod3$ now excludes $b=5,6$ as follows.

For $b=5$, every slack vanishes.  Points on the four-line have
$d_5\in\{0,3\}$, and points off it have $d_5\in\{1,4\}$.  If $a$ counts
off-line degree-four points and $c$ on-line degree-three points, then
\[
 7+3a+3c=25,\qquad a+c=6,\qquad a\ge2.
\]
Consequently
\[
 \sum_p\binom{d_5(p)}2=6a+3c=18+3a\ge24
 >20=2\binom52,
\]
contrary to the two-point intersection bound for five-blocks.

For $b=6$, exactly one pivot has slack one.  If it lies on the four-line,
its five-degree is two; with the same notation,
\[
 a+c=7,\qquad c\le3,\qquad
 \sum_p\binom{d_5(p)}2=1+6a+3c=22+3a>30.
\]
If the exceptional pivot lies off the line, its degree is
$e\in\{0,3\}$, and
\[
 3(a+c)+e=24,\qquad c\le4,\qquad a\ge3,\qquad
 \sum_p\binom{d_5(p)}2=24+3a>30.
\]
Both contradict $2\binom62=30$.

If $L_4=0$, the restored eleven-point set has
$t_2=d_3+10-3\ell_3$.  Kelly--Moser gives $t_2\ge5$, and summing the
resulting bounds over the fourteen triple lines gives
$42\le28+2b$, hence $b\ge7$.  Thus only $b=7,8$ remain.
\end{proof}

\subsection{The eleven-point five-circle branch}

\begin{lemma}[The low layers]\label{lem:n11-g5-low}
In the $n=11$ maximum-circle-five branch one has $C\ge39$.  In particular
$C=36,37,38$ are impossible.
\end{lemma}

\begin{proof}
Deleting a point preserves admissibility unless the other ten points are
collinear or concyclic.  In the first case the outsider pencil gives
$\binom{10}{2}>40$ circles; in the second it gives
$1+\binom{10}{2}-5=41$.  Thus deletion is legitimate below $41$.
 By \cref{cor:pivot-lenchner}, Lenchner gives at least three ordinary proper
 circles through every pivot, and \cref{prop:n10-value} gives $C\ge36$.

For $C=36$, the global rows force
\[
B_3=23,\quad B_5=7,\quad L=12,\quad L_4=L_5=0,
\quad C_3=11.
\]
Equality in the deletion sum gives $c_3(p)=3$ everywhere.  The congruence
in \eqref{eq:n11-link}, Kelly--Moser, and the ray bound give
$\ell_3(p)=3$ everywhere, so $3L_3=33$ although $L_3=12$.

For $C=37$, the only aggregate rows are
\[
(B_5,B_3,L,K)=(7,23,11,1),\qquad(8,25,11,2).
\]
They have respectively one and three pivots with $d_3=9$.  Deletion forces
$c_3(p)\le C-f(10)=4$ at every pivot.  Hence at a pivot with $d_3=9$,
\[
 \ell_3(p)=d_3(p)-c_3(p)\ge9-4=5.
\]
The local added-centre row
$3\ell_3+4\ell_4+5\ell_5\le10+\sigma$ then has left side at least
$15$, whereas $\sigma(p)\le K\le2$ makes its right side at most $12$.

At $C=38$, put $k=B_5$ and $h=L_4$.  The exact reduction is
\begin{equation}\label{eq:n11-C38-front}
B_3=2k+13,\ B_4=38-3k,\ L_3=13-h,\ L_5=0,
\quad k=5,6,7,8,\quad h=0\text{ unless }k=8.
\end{equation}
Every point has $d_3=6$ or $9$.  With
$\kappa=10+\sigma-3\ell_3-4\ell_4$ one has
$\sum\kappa=k-1-7h$.  When $h=0$, the least possible $\kappa$ at a high
point ($d_3=9$) and a low point ($d_3=6$) is
\[
\begin{array}{c|ccccc}
d_5&0&1&2&3&4\\ \hline
\text{high}&1&0&2&1&0\\
\text{low}&1&0&2&1&\text{forbidden}.
\end{array}
\]
For $k=5$ there is one high point and ten low points.  The zero-cost
baseline has degrees $(4,1^{10})$, of sum fourteen.  Supplying the remaining
eleven incidences costs at least five upgrades $1\to3$ and one upgrade
$1\to2$, hence at least $5+2=7>k-1=4$.  Thus $k=5$ is impossible.
At $k=6,7$ equality in the same marginal-cost calculation forces,
respectively,
$(d_5)=(4^3,3^5,1^3)$ and $(4^5,3^4,2,1)$; the pair moments give
\[
\sum\binom{d_5}2=33>2\binom62=30,
\qquad
\sum\binom{d_5}2=43>2\binom72=42.
\]
At $k=8$, both $h=0,1$ force degrees $(4^7,3^4)$.  Rebase at a five-circle
$\Gamma$.  Equations \eqref{eq:K1-local}--\eqref{eq:K1-global} force every
outsider edge to be double-hosted.  A one-$\Gamma$ block with three
outsiders is then forbidden by \cref{lem:circle-lift-trace}; hence no
four- or five-block meets $\Gamma$ once.  The seven other five-blocks and
all fourteen four-blocks therefore meet $\Gamma$ twice, so
\[
H_\Gamma=n_{42}+3n_{52}=14+3\cdot7=35,
\]
whereas the equality face of \eqref{eq:K1-global} gives
$H_\Gamma=30$.  Hence $C=38$ is impossible.
\end{proof}

\begin{lemma}[The \(C=39\) arithmetic front]\label{lem:n11-C39-front-v2}
Assume that eleven admissible points have \(39\) proper circles, contain a
proper five-circle, and have no proper circle larger than five.  Every
generalized block has size at most five, and the number \(L\) of maximal
lines belongs to \(\{12,15\}\).

Fix a proper five-circle \(\Gamma\), with outsider set \(X\).  Let \(A_{ij}\)
count nonselected blocks containing \(i\) points of \(\Gamma\) and \(j\)
outsiders, and put
\[
 H=A_{22}+3A_{23},\quad x=A_{23},\quad
 r=A_{13},\quad s=A_{14},\quad
 q=A_{04}+3A_{05},\quad y=A_{05}.
\]
If \(L=3k\), and if
\[
 B_0=A_{03}+A_{04}+A_{05},\qquad
 T_0=A_{03}+4A_{04}+10A_{05},
\]
then
\begin{align}
A_{22}&=H-3x,&A_{21}&=60-2H+3x,\notag\\
A_{12}&=75-2H-3r-6s,&
B_0&=3H+L-97-x+2r+5s,\notag\\
T_0&=20-x-r-4s,&
q&=39-H-k-r-3s. \label{eq:C39-full-front-v2}
\end{align}
Moreover
\begin{equation}\label{eq:C39-rich-cap-v2}
 0\le q\le3,\qquad y=1\Longrightarrow q=3,
\end{equation}
and, if \(e_\Gamma,e_X\) count points with \(d_3=9\) on the two sides,
\begin{equation}\label{eq:C39-high-split-v2}
 e_\Gamma=55-2H+2x-r-2s,\qquad
 e_X=2k-r-2s-2q+2y.
\end{equation}
Every point has \(d_3\in\{6,9\}\) and \(d_5\le4\).
\end{lemma}

\begin{proof}
Let \(A\) be a maximal \(a\)-line and put \(b=11-a\).  For each pair of
outsiders, its pair-hosts on \(A\) form a matching, so the line pencil gives
the uniform lower bound
\[
 C\ge b\binom a2-\binom b2\floor{a/2}.
\]
For \(a=6,7,8,9,10\) its values are respectively
\(45,66,72,68,45\), so a line of size at least six is impossible.  Let
\(N_j\) be the number of generalized \(j\)-blocks.  Counting blocks and
triples gives
\[
 N_3+N_4+N_5=39+L,\qquad N_3+4N_4+10N_5=165.
\]
Consequently
\begin{equation}\label{eq:C39-global-blocks-v2}
 L=3k,\qquad N_4+3N_5=42-k,\qquad
 N_3=-3+4k+2N_5.
\end{equation}

At a pivot, inversion gives
\[
 d_3+3d_4+6d_5=45.
\]
Kelly--Moser gives \(d_3\ge5\).  Deleting the pivot destroys at most
\(39-f(10)=6\) three-circles, and its three-lines use disjoint pairs of the
other ten points, so \(d_3\le11\).  The displayed congruence now yields
\(d_3\in\{6,9\}\).

If five four-point image lines existed at a pivot, their twenty incidences
on ten points and
\(\sum_u\binom{r_u}{2}\le\binom52\) would force every image point to be
the intersection of exactly two of the lines.  A further image line distinct
from those five contains at most two such pair-labelled points: three labels
would either share an index, giving one of the five lines, or be three
disjoint edges of \(K_5\), which is impossible.  The link equation nevertheless
requires a further three-point line.  Six four-point image lines are excluded
already by \(\sum_u\binom{r_u}{2}\ge18>\binom62\).  Hence \(d_5\le4\).

Melchior on the original lines gives
\[
 3\ell_3+7\ell_4+12\ell_5\le52,
\]
so \(L\le15\).  Summing the pivot slack and the deletion bound gives
\[
 \sum_p(d_3(p)-3-d_5(p))=4L+N_5-42\ge0,
 \qquad 2N_5\le25-k.
\]
These inequalities exclude \(L=0,3,6\).  For \(L=9\), put
\(E=\#\{p:d_3(p)=9\}=2N_5-13\).  Then \(7\le N_5\le11\).
Splitting the pivot slack between \(\Gamma\) and \(X\) gives
\[
 45+4y-4N_5\le3r+7s+6q\le39+4y-3N_5,
 \qquad r+3s+q\ge6.
\]
For \(q=0,1,2,3\), these intervals leave only
\[
 (N_5,r,s)=(7,6,0),(7,1,2),(8,0,2).
\]
The first and third rows have \(H=30\), hence no deficient outsider pair,
although respectively an \(A_{13}\) or an \(A_{14}\) block would require
one.  The middle row has \(H=29\), whose sole deficient edge cannot meet all
four outsider triples of an \(A_{14}\) block.  Thus \(L\ne9\), proving
\(L\in\{12,15\}\).

The three triple partitions relative to \(\Gamma\) give
\eqref{eq:C39-full-front-v2}.  Distinct outsider \(K_4\)'s have disjoint
complement pairs, while an outsider \(K_5\) excludes every \(K_4\).
Therefore either \(A_{05}=1,A_{04}=0\), or \(A_{05}=0,A_{04}\le3\), which
is \eqref{eq:C39-rich-cap-v2}.  Finally, summing three-block incidences on
\(\Gamma\) and on \(X\) gives \eqref{eq:C39-high-split-v2}.
\end{proof}

\begin{lemma}[Added-centre and singleton interfaces]
\label{lem:n11-C39-interfaces-v2}
Under the hypotheses of \cref{lem:n11-C39-front-v2}, put \(m=N_5\).  Then
\begin{equation}\label{eq:C39-line-cut-v2}
 L=12:\quad7\ell_4+16\ell_5\le m+8,
 \qquad
 L=15:\quad7\ell_4+16\ell_5\le m-7.
\end{equation}
If two generalized five-blocks meet in the single point \(p\), then
\begin{equation}\label{eq:C39-singleton-v2}
 (d_3(p),d_5(p))=(9,2).
\end{equation}
Moreover, every generalized five-block contains at most two points of type
\((d_3,d_4,d_5)=(9,4,4)\).
\end{lemma}

\begin{proof}
Write \(\sigma_p=d_3(p)-3-d_5(p)\).  The inversion images are noncollinear:
otherwise inversion back would put all eleven original points on one
generalized block.  Thus \cref{prop:universal-melchior-rows} applies,
gives \(\sigma_p\ge0\), and its restored-centre row is
\[
 3\ell_3(p)+\sum_{j\ge4}j\ell_j(p)\le10+\sigma_p.
\]
After summing over the eleven pivots and using
\eqref{eq:C39-global-blocks-v2},
\[
 9\ell_3+16\ell_4+25\ell_5
 \le110+\sum_p\sigma_p,
 \qquad \sum_p\sigma_p=4L+m-42.
\]
Since \(L=\ell_3+\ell_4+\ell_5\), this is exactly
\eqref{eq:C39-line-cut-v2}.

For the singleton assertion, invert at its carrier \(p\).  The two
five-blocks become disjoint four-point link lines, so \(z=d_5(p)\ge2\).
The cap \(z\le4\), pointwise \(\sigma_p\ge0\), and
\(d_3\in\{6,9\}\) leave only
\[
 (d_3,z)=(6,2),(6,3),(9,2),(9,3),(9,4).
\]
The two cases \(z=3\) are excluded by \cref{lem:n11-K3}(K3.1).  In the
case \((d_3,z)=(9,4)\), the link equation gives \(d_4=4\), so the
complete-quadrangle conclusion (K3.2) says that every pair of base lines
meets at a selected link point.

It remains to exclude \((d_3,z)=(6,2)\).  Dualizing the ten-point link turns
the two disjoint four-point lines into two fourfold points.  The ten dual
lines split into families of sizes \(4,4,2\); after removing pairs already
meeting at the two fourfold points, the four pair classes have sizes
\[
 16,\qquad8,\qquad8,\qquad1.
\]
The link equation requires nine triple vertices.  If the sole pair in the
last class is unused, all nine triples require one pair from each eight-edge
side class.  If it is used, the other eight triples exhaust both side
classes before the exceptional triple uses two further side edges.  Both
alternatives are impossible.  This proves \eqref{eq:C39-singleton-v2} for
circles and for the possible five-line alike.

Finally, a point of type \((9,4,4)\) supplies a harmonic complement on every
incident five-block by (K3.2).  The real cross-ratio argument (K3.3) allows
at most two such complements on five points, proving the last assertion.
\end{proof}

\begin{lemma}[The \(L=15\) harmonic collapse]
\label{lem:n11-C39-L15-v2}
The \(C=39,L=15\) branch is impossible.
\end{lemma}

\begin{proof}
Let \(z_p=d_5(p)\), and call \(p\) high when \(d_3(p)=9\).  From
\eqref{eq:C39-global-blocks-v2},
\[
 E:=\#\{p:p\text{ is high}\}=2m-5,
 \qquad \sum_pz_p=5m.
\]
Thus \(3\le m\le8\).  The second cut in
\eqref{eq:C39-line-cut-v2} forces \(m\in\{7,8\}\) and
\(\ell_4=\ell_5=0\); in particular all fifteen maximal lines are triples
and all generalized five-blocks are proper circles.

Define the nonnegative integral restored-centre slack
\[
 \kappa_p=10+\sigma_p-3\ell_3(p).
\]
Summing gives \(\sum_p\kappa_p=m-7\), while pointwise
\begin{equation}\label{eq:C39-kappa-point-v2}
 3\ell_3(p)=
 \begin{cases}
 13-z_p-\kappa_p,&p\text{ low},\\
 16-z_p-\kappa_p,&p\text{ high}.
 \end{cases}
\end{equation}
For \(m=7\), every \(\kappa_p\) vanishes.  Pointwise
\(\sigma_p\ge0\) and \eqref{eq:C39-kappa-point-v2} make the two low points
have \(z=1\), while
the nine high points have \(z\in\{1,4\}\).  Since \(\sum z=35\), exactly
eight high points have \(z=4\), and the link equation makes their type
\((9,4,4)\).  Counting their incidences with five-blocks gives
\(8\cdot4\le2m=14\), contrary to
\cref{lem:n11-C39-interfaces-v2}.

For \(m=8\), all eleven points are high and precisely one has
\(\kappa=1\).  The ten zero-slack points have \(z\in\{1,4\}\), while the
exceptional point has \(z\in\{0,3\}\).  The sum \(\sum z=40\) therefore
forces respectively ten or nine points with \(z=4\).  Their incidences give
\(40\le16\) or \(36\le16\), again contradicting the harmonic cap.
\end{proof}

\begin{lemma}[The \(L=12\) low-host moment]
\label{lem:n11-C39-low-L12-v2}
In the \(C=39,L=12\) branch, some proper five-circle has host load
\(H\ge26\).
\end{lemma}

\begin{proof}
Suppose every proper five-circle has host load at most \(25\), and put
\[
 E=2m-9,\qquad z_p=d_5(p),\qquad
 G=2\sum_p\binom{z_p}{2}-\sum_{p\text{ high}}z_p.
\]
Here \(m\ge5\), \(\sum_pz_p=5m\), and the K1 moment is
\begin{equation}\label{eq:C39-global-moment-v2}
 \sum_\Delta H_\Delta=20m+G+\sum_\Delta q_\Delta.
\end{equation}
The successive costs of increasing \(z_p\) from zero to four are
\[
 \begin{array}{c|rrrr}
 &1&2&3&4\\ \hline
 p\text{ low}&0&2&4&6\\
 p\text{ high}&-1&1&3&5.
 \end{array}
\]
Filling these increasing marginal costs gives \(G\ge31\) at \(m=5\).
For \(m\ge6\), Cauchy and
\(\sum_{p\text{ high}}z_p\le4(2m-9)\) give
\[
 G\ge\frac{25m^2}{11}-13m+36>5m+5.
\]

The original-line Melchior row is
\(36+4\ell_4+9\ell_5\le52\), so there is at most one five-line.  A proper
five-circle contributes at most \(25\) to the left side of
\eqref{eq:C39-global-moment-v2}, and a possible five-line at most the
universal host cap \(30\).  Hence \(\sum H_\Delta\le25m+5\), whereas the
right side is at least \(131\) for \(m=5\) and is strictly larger than
\(25m+5\) for \(m\ge6\).  This contradiction proves the lemma.
\end{proof}

\begin{lemma}[The \(L=12\) maximum-host router]
\label{lem:n11-C39-maxhost-v2}
In the \(C=39,L=12\) branch, no proper five-circle of maximum host load has
\(26\le H\le30\).
\end{lemma}

\begin{proof}
Let \(h\) be the maximum proper-circle host load, let
\(t=\ell_5\in\{0,1\}\), and use \(m,E,G\) as in the preceding proof.
Besides \(G\ge31\) for \(m=5\), marginal filling gives
\begin{equation}\label{eq:C39-G-small-v2}
 G\ge45\quad(m=6),\qquad G\ge61\quad(m=7).
\end{equation}
and the Cauchy bound above is greater than \(8m+2\) for \(m\ge8\).
Since the possible five-line has host at most thirty, K1 gives
\begin{equation}\label{eq:C39-maxhost-moment-v2}
 20m+G\le h(m-t)+30t.
\end{equation}

For \(h=26\), \eqref{eq:C39-G-small-v2} and
\eqref{eq:C39-maxhost-moment-v2} leave
only \((m,t)=(5,1)\).  For \(h=27\), they
leave \((5,0),(5,1),(6,1)\).  But
\eqref{eq:C39-line-cut-v2} gives \(16t\le m+8\), removing every listed row
with \(t=1\).  Thus \(h\ne26\), and at \(h=27\) necessarily
\begin{equation}\label{eq:C39-H27-face-v2}
 m=5,\qquad t=0.
\end{equation}

We record a consequence used twice.  Rebase at a proper five-circle
\(\Delta\), put \(R=30-H_\Delta\), and let \(r,s,q,y\) have the front
meaning.  If \(R\ge3\) and \(s=0\), then
\begin{equation}\label{eq:C39-no-singleton-front-v2}
 A_{12}=3q-R,\qquad e_X=3-R-q+2y.
\end{equation}
If \(y=1\), then \(q=3\) and the second quantity is \(2-R<0\).  If
\(y=0\), nonnegativity in \eqref{eq:C39-no-singleton-front-v2} gives
simultaneously \(3q\ge R\) and
\(q\le3-R\), again impossible.  Hence every circle with \(H\le27\) has
\begin{equation}\label{eq:C39-singleton-degree-v2}
 s_\Delta\ge1.
\end{equation}

In the face \eqref{eq:C39-H27-face-v2}, form the graph whose five vertices
are the five-blocks and whose edges are their singleton intersections.  By
\eqref{eq:C39-singleton-degree-v2} it has minimum
degree one and hence at least three edges.  But
\eqref{eq:C39-singleton-v2} injects its edges into the \(E=1\) high points:
a carrier has \(d_5=2\) and can carry only that one edge.  This excludes
\(h=27\).

Suppose next that \(h=28\).
Equations \eqref{eq:C39-G-small-v2} and
\eqref{eq:C39-maxhost-moment-v2} leave
\(m\in\{5,6\}\),
and the line cut gives \(t=0\).  We claim that
\eqref{eq:C39-singleton-degree-v2} still holds for a circle
with \(H=28\).  Otherwise its outsider deficiency
\(\delta(e)=2-h_e\) has total weight two.  Its support is one zero edge
\(Z\), two adjacent single edges \(S_a\), or two disjoint single edges
\(S_d\).  Every \(A_{13}\) triangle contains a deficient edge by
\cref{lem:five-conic-rigidity}(K2.1).  A single edge has one host chord;
(K2.1) places that chord in a unique page-colour class, and triple ownership
allows at most one page of that colour through the edge.  Thus a zero edge
lies in four outsider triangles, while each single edge supports at most one
clean page, with the one
triangle on two adjacent singles as the only overlap.  Therefore
\begin{equation}\label{eq:C39-defect-page-caps-v2}
 r\le4,3,2\quad\hbox{in types }Z,S_a,S_d.
\end{equation}
With \(s=0\), the front says \(r+q=7\), so
\eqref{eq:C39-defect-page-caps-v2} and \(q\le3\) force type
\(Z\) and \((r,q)=(4,3)\).  If its zero edge is \(uv\), the four pages are
\(uvw\), one for every other outsider \(w\).  The high split gives
\(e_X=-2+2y\), hence \(y=1\): there is one outsider \(K_5\).  It must omit
\(u\) or \(v\), since otherwise it shares an outsider triple \(uvw\) with
a page, and it is consequently all-double.  Suppose it omits \(u\), and
let \(\gamma_i\) be the colour of the page \(uvw\).  The adjacent double
edges \(uw,vw\) determine two distinct vertices of \(D_i\); triple
ownership makes both host matchings omit \(i\).  The vertex coming from
\(vw\) lies on the all-double \(K_5\) trace and is its unique centre
\(\delta_i\) omitting \(i\), whereas the page trace contains
\(\gamma_i,\delta_i\) and the second vertex coming from \(uw\).  This is
exactly the tangent-separation obstruction in (K2.4).  The case in which
the \(K_5\) omits \(v\) is symmetric.  This proves the claim.

The singleton graph on the \(m\) proper five-circles consequently has
minimum degree one.  If \(m=5\), it has at least three singleton edges,
whereas \(E=1\), a contradiction.  Let \(m=6\).  Now it has at least three
edges and \eqref{eq:C39-singleton-v2} gives at most \(E=3\), so it is a
perfect matching and all three high points have \(z_p=2\).

Let \(A_j\) count pairs of five-circles meeting in \(j\) points, and put
\(Q=\sum_\Delta q_\Delta\).  Since \(A_1=3\),
\[
 A_0+A_1+A_2=15,\qquad
 \sum_p\binom{z_p}{2}=A_1+2A_2=27-2A_0.
\]
The three high points contribute six to the high incidence sum, hence
\begin{equation}\label{eq:C39-m6-G-v2}
 G=48-4A_0.
\end{equation}
Every disjoint pair contributes three to \(q_\Delta\) at each endpoint, so
\(Q\ge6A_0\).  The global moment and \(H_\Delta\le28\) give
\(G+Q\le48\).  Together with \eqref{eq:C39-m6-G-v2}, this forces
\begin{equation}\label{eq:C39-m6-equality-v2}
 A_0=0,\qquad Q=0,\qquad H_\Delta=28
 \quad\hbox{for every }\Delta.
\end{equation}
For any \(\Delta\), the perfect matching gives \(s=1\), while
\eqref{eq:C39-m6-equality-v2} gives
\(q=0\); the front therefore gives \(r=4\).  But an \(A_{14}\) outsider
four-set forces the two deficiency units to be the disjoint perfect
matching on that set.  Every \(A_{13}\) triangle then contains exactly one
of its two single edges, and (K2.1) permits at most one clean page on each.
Thus \(r\le2\), the final contradiction for \(h=28\).

If \(h=29\), the deficiency graph of a maximum-host circle is one single
edge \(uv\).  An \(A_{14}\) support is impossible, and every \(A_{13}\)
triangle contains \(uv\).  The edge has one host chord; (K2.1) assigns that
chord to a unique page colour, and triple ownership permits at most one page
of that colour through \(u,v\).  Hence \(r\le1\).  The front, however, gives
\(r+q=6\), so \(q\le3\) forces \(r\ge3\).

Finally, if \(h=30\), every outsider edge is double.  Clause (K2.1) forbids
an \(A_{13}\) block, and any \(A_{14}\) block contains a forbidden
all-double outsider triangle.  Thus \(r=s=0\), while the front gives
\(q=5>3\).  This exhausts \(26\le h\le30\).
\end{proof}


\begin{lemma}[K2 routing of the $C=39$ layer]\label{lem:n11-C39}
There is no $n=11$, maximum-circle-five configuration with $C=39$.
\end{lemma}

\begin{proof}
By \cref{lem:n11-C39-front-v2}, all generalized blocks have size at most
five and \(L\in\{12,15\}\).  The value \(L=15\) is excluded directly by
the added-centre slack and harmonic cap in
\cref{lem:n11-C39-L15-v2}.  If \(L=12\),
\cref{lem:n11-C39-low-L12-v2} supplies a proper five-circle of host load at
least \(26\).  Choose one of maximum load.  The host matching gives
\(H\le30\), while \cref{lem:n11-C39-maxhost-v2} excludes every integer in
that interval.  Hence \(C=39\) is impossible.
\end{proof}

\begin{lemma}[$C=40$ pivot spectrum]\label{lem:n11-C40-spectrum}
Suppose that an admissible eleven-point set has forty proper circles and no
proper circle larger than five.  At every point \(p\),
\[
 d_3(p)\in\{6,9,12\},\qquad d_5(p)\le4,\qquad
 d_3(p)+1\ge2d_5(p).
\]
If, in addition, \(L=11\), then
\begin{equation}\label{eq:C40-pivot-domains}
\begin{array}{c|c}
d_3(p)&\text{possible }d_5(p)\\ \hline
6&1,2,3\\
9&3,4.
\end{array}
\end{equation}
Consequently, if \(k=B_5\) and \(n_i=\#\{p:d_3(p)=i\}\), then
\begin{equation}\label{eq:C40-pivot-counts}
 n_9=2k-9,\qquad n_6=20-2k .
\end{equation}
\end{lemma}

\begin{proof}
Write \(c_3(p)\) and \(\ell_3(p)\) for the numbers of three-circles and
three-lines through \(p\).  The set \(P\setminus\{p\}\) is admissible:
if its ten points were collinear, the remaining point would determine
\(\binom{10}{2}=45\) distinct proper three-circles with pairs on that line;
ten concyclic points would violate the assumed circle cap.  Deleting \(p\)
destroys only the three-circles through it, so the ten-point theorem gives
\[
 c_3(p)\le C(P)-f(10)=40-33=7.
\]
The three-lines through \(p\) use disjoint pairs of the other ten points,
and hence \(\ell_3(p)\le5\).  Thus \(d_3(p)\le12\).

Invert about \(p\), and denote the image of \(P\setminus\{p\}\) by \(Q_p\).
This set too is admissible.  Indeed, a line or circle containing all of
\(Q_p\) inverts either to a ten-point line in \(P\), or to a circle
containing all of \(P\); both alternatives contradict the preceding caps.
The ordinary image lines are precisely the three-blocks through \(p\).
Kelly--Moser therefore gives \(d_3(p)\ge5\), while the exact link equation
\[
 d_3(p)+3d_4(p)+6d_5(p)=45
\]
makes \(d_3(p)\) divisible by three.  Thus \(d_3(p)\ge6\), and this proves
\(d_3(p)\in\{6,9,12\}\).

By \cref{lem:n11-n12-outer-router}, every original line has at most five
selected points; together with the circle cap, every image line therefore
has at most four of the ten points.  Since
\(4\le2|Q_p|/3\), Langer's incidence inequality applies and yields
\[
 2d_3(p)+3d_4(p)+4d_5(p)\ge44.
\]
Subtracting the link equation gives
\begin{equation}\label{eq:C40-pivot-langer}
 d_3(p)+1\ge2d_5(p).
\end{equation}

It remains to prove the absolute cap \(d_5(p)\le4\).  Suppose first that
five four-point image lines \(L_1,\ldots,L_5\) exist, and let \(r_x\) be
the number of them through \(x\in Q_p\).  Then
\[
 \sum_x r_x=20,\qquad
 \sum_x\binom{r_x}{2}\le\binom52=10.
\]
Convexity gives the reverse inequality, with equality only when
\(r_x=2\) at all ten points.  Thus the points of \(Q_p\) are exactly the
ten intersections \(L_i\cap L_j\), labelled by the two-subsets
\(\{i,j\}\subset[5]\).  A further line distinct from all \(L_i\) contains
at most two of these points: among three two-subsets of a five-set, either
two share an index, in which case their joining line is the corresponding
\(L_i\), or all three would have to be pairwise disjoint, which is
impossible.  Hence a sixth four-point image line cannot exist.  If there
are exactly five, the already proved range of \(d_3\) and the link equation
give
\[
 d_4(p)=5-\frac{d_3(p)}3\ge1,
\]
but the resulting additional three-point image line is equally impossible.
Therefore \(d_5(p)\le4\).

Now assume \(L=11\).  There are \(B=C+L=51\) generalized blocks.  Under
inversion, the circles of \(Q_p\) are in bijection with the blocks of \(P\)
which avoid \(p\); hence
\[
 C(Q_p)=51-\beta(p)\ge f(10)=33,\qquad
 \beta(p):=d_3(p)+d_4(p)+d_5(p)\le18.
\]
Eliminating \(d_4(p)\) by the link equation gives
\begin{equation}\label{eq:C40-pivot-beta}
 \beta(p)=15+\frac{2d_3(p)}3-d_5(p).
\end{equation}
For \(d_3=12\), this forces \(d_5\ge5\), already excluded.  For
\(d_3=6\), \eqref{eq:C40-pivot-beta} and
\eqref{eq:C40-pivot-langer} give \(1\le d_5\le3\); for \(d_3=9\) they give
\(3\le d_5\le4\).  This is \eqref{eq:C40-pivot-domains}.

Finally, the block equations give \(B_3=13+2k\).  Double-counting incidences
with three-blocks and using \(n_6+n_9=11\), we obtain
\[
 6n_6+9n_9=3B_3=39+6k,\qquad n_6+n_9=11,
\]
whose solution is \eqref{eq:C40-pivot-counts}.
\end{proof}

\begin{lemma}[$C=40,L=11$ singleton exclusion]
\label{lem:n11-C40-no-singletons}
In the $C=40,L=11$ layer, two distinct five-blocks meet in zero or two
points, never in exactly one.
\end{lemma}

\begin{proof}
Suppose two five-blocks meet only at $p$.  After inversion at $p$, their
four-point link lines are disjoint.  The spectrum
\eqref{eq:C40-pivot-domains} and $d_5(p)\ge2$ leave
\[
 (d_3,d_5)=(6,2),(6,3),(9,3),(9,4).
\]
Both degree-three cases are excluded by \cref{lem:n11-K3}(K3.1).  In the
last case the link equation gives $d_4=4$, and (K3.2) makes the four base
lines a complete quadrangle, so no pair is disjoint.

It remains to exclude $(d_3,d_5)=(6,2)$, whose link profile is $(6,9,2)$.
Dualize the two disjoint four-point lines.  The ten dual lines split into
families of sizes $4,4,2$ through the two resulting fourfold points; after
removing pairs already meeting there, the four pair classes have sizes
\[
 16,\qquad8,\qquad8,\qquad1.
\]
The nine three-point link lines become nine triple vertices.  If none uses
the sole pair in the last class, all nine require one edge from each
eight-edge side class.  If one uses that pair, the other eight triples
exhaust both side classes before the exceptional triple uses two further
edges on one side.  Both alternatives are impossible.
\end{proof}

\begin{lemma}[The $C=40,L=11$ defect--harmonic faces]
\label{lem:n11-C40-smallfaces}
In the $n=11$, maximum-circle-five layer with $C=40$ and $L=11$, neither
$B_5=5$ nor $B_5=6$ is possible.
\end{lemma}

\begin{proof}
By \cref{lem:n11-C40-no-singletons}, two five-blocks meet in zero or two
points.
For a five-block $E$ put
\[
 \delta_E=\sum_{F\ne E}(2-|E\cap F|).
\]
Thus every summand is zero or two and
\[
 D:=\sum_{\{E,F\}}(2-|E\cap F|)
   =2\binom{B_5}{2}-\sum_p\binom{d_5(p)}2
\]
is twice the number of disjoint block pairs.

\smallskip
\noindent\emph{Five blocks.}
There are ten low points with $d_3=6$ and one high point with $d_3=9$.
Let $a\in\{0,1\}$ indicate whether the high point has $d_5=4$, and let
$n_i$ count low points with $d_5=i$ for $i=1,2$; all other low points have
degree three.  The incidence and defect sums are
\[
 n_2+2n_1=a+8,\qquad D=3-a-n_1.
\]
Since $D$ is nonnegative and even, the complete set of faces is
\[
\begin{array}{c|c}
D&(a,n_1,n_2)\\ \hline
2&(0,1,6),(1,0,9)\\
0&(0,3,2),(1,2,5).
\end{array}
\]
If $c_E,d_E$ count in $E$ the points of degrees two and one and $a_E$
indicates the degree-four high point, then the row sum gives
\begin{equation}\label{eq:C40-five-row-defect}
 \delta_E=c_E+2d_E-a_E-2.
\end{equation}

On the face $(a,n_1,n_2)=(1,0,9)$, the unique disjoint pair $F,G$
satisfies $c_E-a_E=4$ at both ends.  A block containing the high point has
$c_E-a_E\le3$, so the high point lies outside $F\cup G$.  Each of the two
disjoint blocks would then need four degree-two points and the unique
degree-three point, impossible.

On the other $D=2$ face the degrees are $(3^4,2^6,1)$.  If $U$ is the
degree-one point and $F,G$ the disjoint pair, \eqref{eq:C40-five-row-defect}
forces, after exchanging $F,G$,
\[
 F:\ U,\ 2^2,\ 3^2,\qquad G:\ 2^4,\ 3.
\]
The fourth degree-three point $q$ lies outside $F\cup G$ and hence belongs
to all three remaining blocks $A,B,C$; their other pairwise intersections
are the three degree-three points in $F\cup G$.  Invert at $q$.  The images
of $A,B,C$ are three four-lines, each split $2+2$ across the partition
$F\sqcup G$.  If $q$ were low, the seven residual triple lines would each
use two cross-pairs, but the three four-lines have already used twelve of
the $25$ cross-pairs, leaving only thirteen instead of fourteen.  Thus $q$
is the high point.

Rebase now at $\Gamma=F$.  All its points are low, and
\eqref{eq:C40-five-row-defect} gives
\[
 \sum_{p\in F}d_5(p)=11,\qquad
 \sum_{p\in F}d_4(p)=5\cdot13-2\cdot11=43.
\]
No four-block is disjoint from $F$, since four points of $G\cup\{q\}$
contain a triple of $G$.  Hence, if $n_{22},n_{13}$ count four-blocks
meeting $F$ in two and one points,
\[
 n_{22}+n_{13}=23,\qquad2n_{22}+n_{13}=43,
 \quad\text{so}\quad(n_{22},n_{13})=(20,3).
\]
The ten edges within $G$ receive seventeen hosts from the
$2F+2G$ four-blocks and three from $A,B,C$, so every such edge is
double-hosted.  The five $qG$ edges receive nine hosts; four are
 double-hosted.  Each of the three $q+F+2G$ four-blocks therefore contains
 a double edge of $G$ and an adjacent double $qG$ edge.  Both host matchings
 omit the colour of the page: a host chord through that point would repeat
 the triple already owned by the page.  This contradicts the
 golden-separation clause \cref{lem:five-conic-rigidity}(K2.4).

It remains to dispose of the two $D=0$ faces simultaneously.  Choose a
five-block $\Gamma$ avoiding the high point.  Every point on $\Gamma$ is
low, every other five-block meets it twice, and
\[
 \sum_{p\in\Gamma}d_5(p)=13,\qquad
 \sum_{p\in\Gamma}d_4(p)=39.
\]
Let $t$ count four-blocks disjoint from $\Gamma$, and let $n_{22},n_{13}$
count the other four-block types.  Triple ownership gives
\begin{align}
 n_{22}&=16+t,&n_{13}&=7-2t,\notag\\
 H&=n_{22}+3\cdot4=28+t,&
 n_{12}&=4t-2,\qquad n_{03}=9-2t.\label{eq:C40-five-routing}
\end{align}
Here $H$ is the host load on the fifteen outsider edges, and $n_{12},n_{03}$
count triple blocks of the indicated types.  Host matching gives $H\le30$,
while $n_{12}\ge0$, so $t=1$ or $2$.  The graph of double-hosted outsider
edges is then $K_6$ minus one edge, with sixteen triangles, or $K_6$, with
twenty.  By \cref{lem:five-conic-rigidity}(K2.1), every such triangle is
owned either by one of the four $2\Gamma+3X$ five-blocks or by a block
disjoint from $\Gamma$.  Their total triangle capacity is only
\[
 4+(n_{03}+4t)=13+2t,
\]
namely fifteen or seventeen.  Both faces are impossible.

\smallskip
\noindent\emph{Six blocks.}
There are eight low and three high points.  Let $a$ count high points with
$d_5=4$, and again let $n_i$ count low points of degree $i$.  Then
\[
 n_2+2n_1=a+3,\qquad D=3-a-n_1,
\]
so the complete face list is
\[
\begin{array}{c|c}
D&(a,n_1,n_2)\\ \hline
2&(0,1,1),(1,0,4)\\
0&(1,2,0),(2,1,3),(3,0,6).
\end{array}
\]
If $h_E,c_E,d_E$ count in $E$ the points of degrees four, two, and one,
respectively, then
\begin{equation}\label{eq:C40-six-row-defect}
 \delta_E=c_E+2d_E-h_E.
\end{equation}
For $(0,1,1)$ the two ends of the disjoint pair require total defect four,
but all degree-two and degree-one points have total weight only three in
\eqref{eq:C40-six-row-defect}.  For $(1,2,0)$, zero defect would give
$2d_E=h_E$ in every block; with only one high point this forces
$d_E=h_E=0$, so no block could contain either degree-one point.

Consider $(1,0,4)$.  Its degrees are $(4,3^6,2^4)$.  Equality at the
disjoint pair $F,G$ places the high point $h$ outside their union, all four
degree-two points inside it, and two in each block.  The other four
five-blocks are precisely the blocks through $h$.  After inversion at $h$
they are four lines in complete-quadrangle position, with six pair-points
$v_{ij}$ of degree three and four private points $u_i$ of degree two;
$F,G$ partition these ten points.  Colour $ij$ and $i$ by $F$ when
$v_{ij}$ and $u_i$ lie in $F$.  The equality
\[
 \deg_F(i)+\mathbf1_{\{u_i\in F\}}=2
\]
and $|F|=5$ show that the $F$-edges form a Hamilton path and its two private
points are the endpoints; the $G$-data are complementary.  If $ij$ is the
endpoint pair and $kl$ the internal pair of this path, then the residual
triples $S_{ij}=\{u_i,u_j,v_{kl}\}$ and
$S_{kl}=\{u_k,u_l,v_{ij}\}$ lie in $F$ and $G$.  The triangle-plus-pendant
conclusion of \cref{lem:n11-K3}(K3.2) selects at least one of the two
complementary indices.  That triple would then belong both to $F$ or $G$
and to a distinct four-block through $h$, violating triple ownership.

For $(2,1,3)$ the degrees are $(4^2,3^5,2^3,1)$.  Let the high points be
$h_1,h_2$ and the degree-one point be $U$.  Equation
\eqref{eq:C40-six-row-defect} forces the block $Z$ through $U$ to contain
both high points and no degree-two point.  The cap
\cref{lem:n11-K3}(K3.3), followed by the incidence and pair rows, forces
one further common block $W$ and four side blocks.  Relabeling their
supports gives
\begin{equation}
\begin{array}{c|l}
Z&h_1,h_2,U,m,n\\
W&h_1,h_2,r,\ell_1,\ell_2\\
A_1&h_1,m,r,\ell_3,s\\
A_2&h_1,n,\ell_1,s,t\\
B_1&h_2,n,r,\ell_3,t\\
B_2&h_2,m,\ell_2,s,t.
\end{array}\label{eq:C40-harmonic-skeleton}
\end{equation}
For clarity, the support is forced without a catalogue: $Z$ has the two
high points; among the other five blocks their remaining six high
incidences split as one double and four single incidences.  Pair ownership
then places the three degree-two incidences and successively the five
degree-three incidences, yielding the displayed crossed pattern.

Invert first at $h_1$ and then at $h_2$.  In either link, K3.2 gives the
unique real triangle-plus-pendant orientation.  On the common block $W$,
the first orientation says that $h_2$ is the harmonic conjugate of $r$
with respect to the unordered pair $\{\ell_1,\ell_2\}$; the second says the
same of $h_1$.  A harmonic conjugate is unique, so $h_1=h_2$, a
contradiction.

Finally, on $(3,0,6)$ the degrees are $(4^3,3^2,2^6)$.  Zero row defect and
K3.3 imply that every block contains exactly two high, two degree-two, and
one degree-three point.  Calling the highs $h_1,h_2,h_3$, each high pair
labels two blocks:
\[
 X_0,X_1\supset h_1h_2,\qquad
 Y_0,Y_1\supset h_1h_3,\qquad
 Z_0,Z_1\supset h_2h_3.
\]
The two degree-three points choose, after relabeling, the triples
$X_0Y_0Z_0$ and $X_1Y_1Z_1$; the six opposite-index block pairs own the six
degree-two points.  Invert at $h_1$ with bases $X_0,X_1,Y_0,Y_1$.  The two
complementary residual triples satisfy
\[
 S_{X_0Y_0}\subset Z_1,\qquad S_{X_1Y_1}\subset Z_0.
\]
The complement of the triangle-plus-pendant graph in K3.2 is a pair of
adjacent edges, so it cannot omit both complementary indices.  One of
these two triples is therefore also owned by a four-block through $h_1$,
again contradicting triple ownership.  All nine faces are excluded.
\end{proof}

\begin{lemma}[Seven five-blocks]\label{lem:n11-seven-five}
In the \(n=11\), maximum-circle-five layer with \(C=40\) and \(L=14\),
the row \(B_5=7\) is impossible.
\end{lemma}

\begin{proof}
Let \(F_1,\ldots,F_7\) be the five-blocks and put \(d(p)=d_5(p)\).  Since
two distinct maximal blocks meet in at most two points,
\[
 \sum_p d(p)=35,\qquad
 I:=\sum_p\binom{d(p)}2
   =\sum_{i<j}|F_i\cap F_j|\le42.
\]
With \(x_p=d(p)-3\), these identities become
\begin{equation}\label{eq:seven5-moment}
 \sum_p x_p=2,\qquad I=38+\frac12\sum_p x_p^2,
 \qquad \sum_p x_p^2\le8.
\end{equation}
By \cref{lem:n11-C40-spectrum}, \(0\le d(p)\le4\); degree zero alone would
contribute nine to the last sum, so \(1\le d(p)\le4\).  If \(a,b,c,e\)
count the points of degrees \(4,3,2,1\), respectively, then
\[
 a-c-2e=2,\qquad c+3e\le3.
\]
Thus either \(e=0\) and \(c=0,1,2,3\), or \(e=1,c=0\).  This gives exactly
\begin{equation}\label{eq:seven5-profiles-lemma}
\begin{array}{c|c|c}
 & (d(p):p\in P)&D:=42-I\\ \hline
A&4^2\,3^9&3\\
B&4^3\,3^7\,2&2\\
C&4^4\,3^5\,2^2&1\\
D&4^5\,3^3\,2^3&0\\
E&4^4\,3^6\,1&0.
\end{array}
\end{equation}
This is an algebraic exhaustion, not a search over supports.

The triple partition gives \(B_3=31\) and \(B_4=16\).  If
\(u=\#\{p:d_3(p)=12\}\), the spectrum lemma and
\(\sum_p d_3(p)=3B_3=93\) give
\begin{equation}\label{eq:seven5-d3-counts}
 n_6=u+2,\qquad n_9=9-2u,\qquad n_{12}=u,\qquad0\le u\le4.
\end{equation}

\smallskip
\noindent\emph{The external-trace obstruction.}
Fix \(H\) with \(d(H)=4\).  Remove \(H\) from its four five-blocks and,
on the remaining ten points, let \(r_x\) be the number of the resulting
four-sets through \(x\).  Then
\[
 \sum_xr_x=16,\qquad \sum_x\binom{r_x}{2}\le6.
\]
Convexity gives the reverse inequality.  Equality therefore holds: there
are four private points \(u_i\) and six pair-points \(v_{ij}\), one for
each pair of the four bases.

Let \(G\) be one of the three five-blocks avoiding \(H\), and suppose that
it contains \(k\) pair-points.  It contains \(5-k\) private points, and its
total intersection with the four bases is \(5+k\).  Each base meets \(G\)
at most twice, while only four private points exist; hence
\(1\le k\le3\).  Let \(U\subset[4]\) index the private points in \(G\), and
let \(E_G\subset\binom{[4]}2\) index its pair-points.  Some
\(kl\in E_G\) has complementary pair \(\{i,j\}\subset U\).  Indeed, this
is immediate for \(k=1\).  For \(k=2\), otherwise the two edges of \(E_G\)
would be a matching on the three indices in \(U\); for \(k=3\), otherwise
three edges would have to use the total residual capacity two of the two
indices in \(U\).  Consequently
\begin{equation}\label{eq:seven5-external-R}
 R_{ij}:=\{u_i,u_j,v_{kl}\}\subset G,\qquad\{i,j,k,l\}=[4].
\end{equation}
The three external blocks give three distinct \(R_{ij}\), since two
maximal blocks cannot share a triple.

If \(H\) had \((d_3,d_5)=(9,4)\), then
\cref{lem:n11-K3}(K3.2) would select four of the six \(R_{ij}\), in the
triangle-plus-pendant pattern, as traces of four-blocks through \(H\).
A four-subset and a subset of at least three among six candidates intersect.
The corresponding triple would belong both to such a four-block and to an
external five-block, contrary to triple ownership.  Thus no point has type
\((9,4)\).  Since \eqref{eq:C40-pivot-langer} forbids \((6,4)\), every
degree-four point has type \((12,4)\), and there are at most four of them.

We shall also use the following two immediate carrier facts.  A degree-four
point cannot carry a singleton intersection: its four bases form the
complete quadrangle just obtained, so every two share a second point.
A degree-three point carries at most one singleton, since two singleton
pairs among its three bases would make their three four-point traces occupy
at least \(4+4+4-1=11\) of the ten link points.  Finally,
\cref{lem:n11-K3}(K3.1) says that a degree-three singleton carrier cannot
have \(d_3=6\) or \(9\).  Hence
\begin{equation}\label{eq:seven5-singleton-carrier}
 d_5(p)=3\text{ and \(p\) carries a singleton}\quad\Longrightarrow\quad
 d_3(p)=12.
\end{equation}

\smallskip
\noindent\emph{Three immediate profiles.}
A disjoint pair is impossible in profile \(A\).  Indeed, if \(a_F\) counts
degree-four points on \(F\), its row defect is
\[
 \delta_F:=\sum_{G\ne F}(2-|F\cap G|)=2-a_F.
\]
The two ends of a disjoint pair would both have \(a_F=0\), although their
ten-point union cannot omit both degree-four points.  Thus \(D=3\) consists
of three singleton pairs.  By \eqref{eq:seven5-singleton-carrier} their
three distinct carriers have \(d_3=12\); together with the two degree-four
points this contradicts \(u\le4\).

In profile \(C\), the unique defect is one singleton.  Its carrier cannot
have degree four, and by \eqref{eq:seven5-singleton-carrier} it cannot have
degree three, for the four degree-four points already exhaust \(u\le4\).
Call the carrier \(q\), and call the other degree-two point \(r\).  If
\(a_F\) counts the degree-four points of \(F\), the row identity is
\[
 a_F=2+\mathbf1_{\{r\in F\}}.
\]
The two blocks through \(r\) therefore contain \(r\) and three of the four
degree-four points; they share at least three points, a contradiction.
Profile \(D\) is still shorter: its five degree-four points would all have
\(d_3=12\), again contradicting \(u\le4\).

\smallskip
\noindent\emph{The path residue \(B\).}
Here the row defect is
\(\delta_F=2-a_F+c_F\), where \(c_F\) records the unique degree-two point.
The same ten-point-union argument excludes a disjoint pair.  The two units
of defect are therefore singleton pairs.  At most one is carried by the
degree-two point; \eqref{eq:seven5-singleton-carrier} and the three
degree-four points force equality in \eqref{eq:seven5-d3-counts}:
\[
 u=4,
\]
with three points of type \((12,4)\), one singleton carrier of type
\((12,3)\), and the degree-two point carrying the other singleton.

Label the seven blocks by a set \(V\).  Let \(T_i\) be the three-element
complement in \(V\) of the support of the \(i\)-th \((12,4)\) point, and
let \(E,Q\in\binom V2\) be the two singleton pairs, carried respectively
by the degree-three and degree-two points.  Pair multiplicity and the row
defect give, if \(t_v=\#\{i:v\in T_i\}\),
\begin{equation}\label{eq:seven5-path-row}
 t_v=1+\deg_{E\cup Q}(v)-\mathbf1_{\{v\in Q\}}.
\end{equation}
Thus precisely the two ends of \(E\) have \(T\)-degree two and all other
labels have degree one.  Moreover \(|T_i\cap T_j|\le1\): at a degree-four
pivot the six complete-quadrangle pair-points are distinct, so two
complement triples cannot reuse a pair of block labels.  The three linear
triples consequently form a path;
after relabelling,
\[
 T_1=\{a,x_1,x_2\},\qquad
 T_2=\{a,b,y\},\qquad
 T_3=\{b,z_1,z_2\},\qquad E=\{a,b\}.
\]
Because \(Q\) is itself a singleton, it meets every \(T_i\).
The two-point transversals of this path, apart from \(E\), allow us to take
\(Q=\{a,z_1\}\).

Writing \(w\) for the third block through the carrier of \(E\), pair
multiplicity leaves, up to the path reflections, \(w=x_1,y,z_2\).
The choice \(w=y\) forces both \(ax_1z_2,ax_2z_2\) and leaves the pair
\(by\) uncovered; the choice \(w=z_2\) forces
\(az_2x_1,yz_2x_2\) and then reuses \(x_2y\).  Hence \(w=x_1\), and all
positive residual pair degrees force the six supports
\begin{equation}\label{eq:seven5-path-six}
 ayz_2,\quad ax_2z_2,\quad bx_1z_2,\quad
 x_1yz_1,\quad byx_2,\quad bx_2z_1.
\end{equation}
At the middle pivot the four forbidden \(S\)-indices are
\[
 x_1x_2,\quad x_1z_1,\quad x_2z_2,\quad z_1z_2,
\]
a four-cycle; the only safe indices are the matching
\(\{x_1z_2,x_2z_1\}\).  A \((12,4)\) pivot has three further local
three-lines, each of type \(T\) or \(S\), and at most one is of type \(T\).
With no \(T\), three safe \(S\)'s are required, but only two exist.  With
one \(T\), the other two safe \(S\)'s must meet its complementary vertex,
whereas the safe matching has degree one at every vertex.  Profile \(B\)
is impossible.

\smallskip
\noindent\emph{The degree-one residue \(E\).}
Let \(q\) be the degree-one point, \(E\) its five-block, and
\(D_1,\ldots,D_4\) the four points of type \((12,4)\).  Zero defect gives
\[
 E=\{q,D_1,D_2,D_3,D_4\}.
\]
Index the other six blocks by the edges of \(K_4\).  The residual pair
multigraph on these six labels has multiplicity two on the three pairs of
opposite edges and multiplicity one on the twelve adjacent pairs.  Its six
degree-three supports each contain exactly one opposite pair.  Grouping
the labels into the three opposite pairs \(A,B,C\), all four cross-edges
between any two groups must be assigned in one direction: assignments in
both directions would repeat one cross-edge and omit another.  Since every
group anchors two triples, these directions form a directed three-cycle.
Up to reversal and relabelling, the six supports are therefore
\begin{equation}\label{eq:seven5-directed-six}
\begin{split}
&\{12,34,13\},\ \{12,34,24\},\
  \{13,24,14\},\ \{13,24,23\},\\
&\{14,23,12\},\ \{14,23,34\}.
\end{split}
\end{equation}
This derived orientation has two consequences visible directly in
\eqref{eq:seven5-directed-six}: at each \(D_i\) the forbidden \(S\)-indices
form the triangle on the three bases other than \(E\), and the complements
of the six incidence triples cover all fifteen pairs of the six
degree-three points.

Since \(d_3(D_i)=12\) and \(d_5(D_i)=4\), the link equation gives
\(d_4(D_i)=3\).  Avoiding the forbidden triangle forces its three
four-blocks to be either the \(S\)-star centred at \(E\), or one \(T\) and
two edges of that star.  Every such \(T/S\) completion contains exactly one
of the four points \(D_i\); hence it is counted at only one pivot.  Thus each
\(D_i\) supplies three distinct four-blocks, at least two through \(q\).
These are twelve distinct blocks
and exhaust all four-block incidences with the \(D_i\).  Globally
\(B_4=16\), while the link equation gives \(d_4(q)\le11\).  Hence among the
four remaining four-blocks at least one avoids both \(q\) and all \(D_i\);
it consists of four of the six degree-three points.  Its complementary pair
is contained in the complement of one incidence triple by the covering
property above.  Equivalently, one of the five-blocks in
\eqref{eq:seven5-directed-six} shares three points with this four-block,
contrary to triple ownership.  Profile \(E\), and hence every row in
\eqref{eq:seven5-profiles-lemma}, is impossible.
\end{proof}

\begin{lemma}[The $C=40$ endpoint dispatch]\label{lem:n11-C40}
There is no $n=11$, maximum-circle-five configuration with $C=40$.
\end{lemma}

\begin{proof}
Put $k=B_5$ and $B=C+L=40+L$.  The block and triple partitions give
\[
 N_3+N_4+k=B,\qquad N_3+4N_4+10k=165,
\]
and hence
\begin{equation}\label{eq:C40-global-router}
 3N_4+9k=125-L,\qquad L\equiv2\pmod3.
\end{equation}
Moreover
\[
 \sum_p\beta(p)=3N_3+4N_4+5k
               =\frac{8L+485}{3}-k,
\]
where $\beta(p)$ is the number of blocks through $p$.  Inversion at $p$
turns precisely the $B-\beta(p)$ blocks avoiding $p$ into proper circles
of the admissible ten-point image.  Thus
$B-\beta(p)\ge f(10)=33$, so $\beta(p)\le L+7$.  Summing over the eleven
pivots yields
\begin{equation}\label{eq:C40-inversion-router}
 25L+3k\ge254.
\end{equation}
By \cref{lem:n11-C40-spectrum}, $5k=\sum_pd_5(p)\le44$, hence $k\le8$;
therefore \eqref{eq:C40-inversion-router} and
\eqref{eq:C40-global-router} give $L\ge11$.

On the other hand, the line-pair partition and Melchior give
\[
 3L+4L_4+9L_5\le52,
\]
so $L\le17$.  The congruence in \eqref{eq:C40-global-router} leaves
$L\in\{11,14,17\}$.  If $L=17$, the last inequality forces
$L_4=L_5=0$, hence $L_3=17$, contrary to the eleven-point ceiling
$L_3\le16$ in \cref{lem:real-nine-link}.  Consequently
\[
 L\in\{11,14\}.
\]
For $L=11$ the four mutually exclusive rows are
\begin{equation}
\begin{array}{c|c|c}
B_5&(n_6,n_9,n_{12})&(N_3,N_4,N_5)\\ \hline
5&(10,1,0)&(23,23,5)\\
6&(8,3,0)&(25,20,6)\\
7&(6,5,0)&(27,17,7)\\
8&(4,7,0)&(29,14,8).
\end{array}\label{eq:C40-L11}
\end{equation}
The cases $B_5=5,6$ are excluded by
\cref{lem:n11-C40-smallfaces}.  For $B_5=7$, let $a$ be the number of
points of type $(d_3,d_5)=(9,4)$, and put
\[
 b=\sum_{p:d_3(p)=6}(3-d_5(p)).
\]
The domains in \cref{lem:n11-C40-spectrum} and $\sum_pd_5(p)=35$ give
$a-b=2$.  By K3.3, double-counting harmonic certificates over the seven
five-blocks gives $4a\le14$.  Hence $(a,b)=(2,0)$ or $(3,1)$, with degree
multisets
\[
(4^2,3^9)\quad\text{or}\quad(4^3,3^7,2).
\]
Put
\[
 D=2\binom72-\sum_p\binom{d_5(p)}2.
\]
The two displayed profiles give respectively $D=3$ and $D=2$.
By \cref{lem:n11-C40-no-singletons}, the first profile cannot realize its
odd defect.

In the second profile, the same lemma says that the defect must be one
disjoint pair $F,G$.  If $a_E$ and $c_E$ count the degree-four and
degree-two points in a block $E$, its row defect is
\[
 \delta_E=\sum_{J\ne E}(2-|E\cap J|)=2-a_E+c_E.
\]
The disjoint pair exhausts the total defect, so
$\delta_F=\delta_G=2$ and $a_F=c_F$, $a_G=c_G$.  Since $F,G$ are disjoint
and there is only one degree-two point, their union contains at most one of
the three degree-four points.  But $|F\cup G|=10$ misses only one of the
eleven points and therefore contains at least two, a contradiction.

For $B_5=8$, the profile is
\[
7(9,4,4)+4(6,7,3).
\]
Thus
\[
 S_{55}=7\binom42+4\binom32=54,
 \qquad I_\omega=7\cdot4=28.
\]
Since $58-N_4-3B_5=20$, \eqref{eq:K1-global} becomes
\[
 \sum_\Delta H_\Delta=240+\sum_\Delta q_\Delta.
\]
The bounds $H_\Delta\le30$ and $q_\Delta\ge0$ force
$H_\Delta=30$ and $q_\Delta=0$ for every five-block $\Delta$.  Choose a
proper five-circle $\Gamma$.  By
\cref{lem:five-conic-rigidity}(K2.1), the all-double triangle obstruction gives
$n_{41}=n_{51}=0$, while $q_\Gamma=0$ gives $n_{40}=n_{50}=0$.  Hence
$n_{42}=14,n_{52}=7$ and
$H_\Gamma=14+3\cdot7=35$, contradiction.  This closes all of
\eqref{eq:C40-L11}.

Now let $L=14$.  By \cref{lem:n11-K4}, only $B_5=7,8$ remain.
The seven-block row is excluded by \cref{lem:n11-seven-five}.  For
$B_5=8$, write $u=n_{12}$.  The pivot equations give
\[
 (n_6,n_9,n_{12})=(u,11-2u,u).
\]
Points with $d_3=6$ have $d_5\le3$, and all points have $d_5\le4$.
If $a$ counts the points of type $(d_3,d_5)=(9,4)$, the degree sum
$\sum d_5=40$ gives $a\ge7-u$.  The harmonic cap K3.3 gives
$4a\le2B_5=16$, so $u\ge3$; the degree maximum $44-u\ge40$ gives
$u\le4$.  Equality in these bounds forces in both cases the common degree
row
\[
 (d_5(p))_{p\in P}=(4^7,3^4),
\]
and every degree-three point has $d_3\in\{6,9\}$.

 The total intersection defect is
 $2\binom82-(7\binom42+4\binom32)=2$.  A singleton pair is impossible.
 Indeed, if its carrier has degree four, delete that carrier from its four
 incident bases and let $r_x$ be the resulting multiplicity of each of the
 other ten points.  The sixteen incidences give
 \[
  \binom42\ge\sum_x\binom{r_x}{2}
  \ge\sum_x(r_x-1)_+\ge16-10=6.
 \]
 Thus equality holds throughout and every pair of those four bases has a
 second intersection point, contradicting singleton intersection.  A
 degree-three carrier is excluded by K3.1.  If instead the defect were one
 disjoint pair $F,G$, and $a_E$
counted degree-four points in $E$, then
\[
 \delta_E=\sum_{J\ne E}(2-|E\cap J|)=4-a_E.
\]
The disjoint pair exhausts the defect, so $a_F=a_G=2$.  Their ten-point
union would contain only four of the seven degree-four points although it
can omit only one.  This contradiction excludes $B_5=8$ and closes the
$L=14$ layer.
\end{proof}

\subsection{The eleven-point six-circle branch}

\begin{lemma}\label{lem:n11-gamma6}
If eleven admissible points have a largest proper circle of size six, then
$C\ge41$.
\end{lemma}

\begin{proof}
Fix the six-circle $\Gamma$ and let $X$ be the five outsiders.  The rich-line
bound excludes lines of size at least six.  For a block type $(g,x)$ write
$c_{gx},\ell_{gx}$, and let $N_s$ count blocks with $s$ outsiders.  Apart
from $\Gamma$, the full type universe is
\[
 \mathcal U=\{(0,3),(0,4),(0,5),(1,2),(1,3),(1,4),(1,5),
                 (2,1),(2,2),(2,3),(2,4)\},
\]
with one circle and one line variable for each type (the already excluded
six-line variables may harmlessly be retained as zeros).  Triple ownership
gives the exact rows
\begin{align}
c_{60}&=1,\notag\\
\sum_{\mathcal U}\binom{x}{3}(c_{gx}+\ell_{gx})&=10,\notag\\
\sum_{\mathcal U}g\binom{x}{2}(c_{gx}+\ell_{gx})&=60,\label{eq:n11-g6-triple-rows}\\
\sum_{\mathcal U}\binom g2x(c_{gx}+\ell_{gx})&=75.\notag
\end{align}
Each outsider triple has a unique owner block, whose footprint on $X$ has
 size three, four, or five.  The collections of three-subsets contained in
 distinct footprints are disjoint, by triple ownership.  Two four-subsets
 of the five-set $X$ share a triple, so at most one four-footprint occurs;
 a five-footprint contains all ten triples and excludes every other
 footprint.  Consequently the ten outsider triples have exactly one of the
 three profiles
\[
(N_3,N_4,N_5)=(10,0,0),(6,1,0),(0,0,1),
\]
 called A, B, C.  Put
\[
 W=c_{22}+3c_{23}+6c_{24}
   =\sum_{e\in\binom X2}q_e.
\]
The nonnegative slacks used below are, explicitly,
\begin{align}
\sigma_6={}&30-(2\ell_{12}+3\ell_{13}+4\ell_{14}+5\ell_{15}
 +2\ell_{21}+4\ell_{22}+6\ell_{23}+8\ell_{24}),\notag\\
\sigma_{10}={}&3c_{03}+2c_{12}+c_{21}-15,\notag\\
\sigma_{17}={}&52-(3\ell_{03}+7\ell_{04}+12\ell_{05}
 +3\ell_{12}+7\ell_{13}+12\ell_{14}+18\ell_{15}\notag\\
&\hspace{38mm}
 +3\ell_{21}+7\ell_{22}+12\ell_{23}+18\ell_{24}),\label{eq:n11-g6-slacks}\\
\sigma_{63}={}&28-W,\notag\\
\sigma_L={}&2-(\ell_{03}+\ell_{13}+\ell_{23}
 +2\ell_{04}+2\ell_{05}+2\ell_{14}+2\ell_{15}+2\ell_{24}).\notag
\end{align}
Here $\sigma_6$ is the unused capacity among the thirty
$\Gamma$--outsider pairs in line blocks, $\sigma_{10}$ is the summed
ordinary-block inequality over the outsider group, and $\sigma_{17}$ is
 the global Melchior slack.  More explicitly, the pivot bridge
 \cref{cor:pivot-lenchner} supplies at least three proper three-circles
 through each outsider.
Summing over the five outsiders counts a $c_{03}$ block three times, a
$c_{12}$ block twice, and a $c_{21}$ block once.  Hence
\[
 3c_{03}+2c_{12}+c_{21}\ge15,
\]
which proves $\sigma_{10}\ge0$.

For $\sigma_L$, consider the outsider footprints of maximal lines which
contain at least three points of $X$.  Distinct footprints meet in at most
one point.  On a five-set, a footprint of size four or five therefore
excludes every other such footprint.  If two three-footprints occur, they
meet in one point and have union $X$; a third three-footprint would meet one
of them in at least two points.  Assigning weight one to a three-footprint
and weight two to a footprint of size at least four gives total weight at
most two, exactly the inequality $\sigma_L\ge0$ printed above.

Finally U17 on each of the five four-subsets of $X$ gives
$3W\le5\cdot17$, hence the integral bound $W\le28$ and
$\sigma_{63}\ge0$.

The triple rows, Melchior rows, and U17 give the following exact
nonnegative-slack identity:
\begin{align}
C-\left(32+\frac9{10}N_3+\frac35N_4\right)
={}&\frac{\sigma_6}{2}+\sigma_{63}+\sigma_L+c_{12}+2c_{24}\notag\\
&+\ell_{04}+\ell_{05}+\ell_{12}+\frac32\ell_{13}
+3\ell_{14}+\frac72\ell_{15}+\ell_{24}.\label{eq:n11-g6-F1}
\end{align}
In footprint C, where $N_4=0$ and $N_5=1$, there is the second exact
identity
\begin{align}
C-\frac{467}{12}
={}&\frac{\sigma_6}{8}+\frac{\sigma_{10}}2
+\frac{\sigma_{17}}{12}+2c_{04}+\frac32c_{13}+2c_{14}\notag\\
&+\frac34\ell_{03}+\frac{19}{12}\ell_{04}
+\frac12\ell_{12}+\frac{35}{24}\ell_{13}
+\frac52\ell_{14}+\frac98\ell_{15}\notag\\
&+\frac1{12}\ell_{22}+\frac34\ell_{23}
+\frac32\ell_{24}.\label{eq:n11-g6-F2}
\end{align}
For a coefficientwise certificate, denote the first and third partition
rows in \eqref{eq:n11-g6-triple-rows} by $T_0,T_2$.  The complete multiplier
data are
\[
\begin{array}{c|ccc|ccccc}
&\multicolumn{3}{c|}{\text{inequality slacks}}
 &c_{60}&T_0&T_2&N_3&N_4\\ \hline
\mathrm{F1}&\frac12\sigma_6&\sigma_{63}&\sigma_L
 &1&\frac1{10}&1&\frac9{10}&\frac35
\end{array}
\]
For F2, after imposing the footprint-C row
$N_4+6(N_5-1)=0$, the multiplier data are
\[
\begin{array}{c|ccc|ccccc}
&\multicolumn{3}{c|}{\text{inequality slacks}}
 &c_{60}&T_0&T_2&N_4&N_5\\ \hline
\mathrm{F2}&\frac18\sigma_6&\frac12\sigma_{10}&\frac1{12}\sigma_{17}
 &1&-\frac12&\frac12&1&6.
\end{array}
\]
Expanding these fixed linear combinations over the displayed universe
$\mathcal U$ gives exactly the nonnegative residues printed in
\eqref{eq:n11-g6-F1} and, in footprint C, \eqref{eq:n11-g6-F2}; thus the
identities are directly
checkable without a type search.  Therefore A gives $C\ge41$, B gives
$C\ge38$, and C gives
$C\ge39$.  At $C=38$ equality in B would give, on writing
$y=c_{13},q=c_{14},f=\ell_{22},w=\ell_{23}$ and $z=c_{23}$,
\[
 3y+6q+2f+6w=4.
\]
The remaining exact rows give
\[
\begin{gathered}
 c_{04}=1-q,\quad c_{03}=4-y-z,\quad \ell_{03}=2-w,\\
 c_{21}=4+3z,\quad c_{22}=28-3z,\quad
 \ell_{21}=15-2f-3w.
\end{gathered}
\]
Substitution in \eqref{eq:n11-g6-slacks} yields
$\sigma_{10}=1-3y\ge0$ and $\sigma_{17}=1-f\ge0$.  Hence
$y\le1/3$, $f\le1$, and the integer $r=q+w$ satisfies
\[
 \frac16\le r=\frac{4-3y-2f}{6}\le\frac23,
\]
which is impossible.

It remains to treat $C=39,40$.  In footprint B,
\eqref{eq:n11-g6-F1}, after substituting
$\sigma_{63}=28-W$, reads exactly
\begin{align}
C+W-66={}&\frac{\sigma_6}{2}+\sigma_L+c_{12}+2c_{24}
 +\ell_{04}+\ell_{05}+\ell_{12}+\frac32\ell_{13}\notag\\
&+3\ell_{14}+\frac72\ell_{15}+\ell_{24}.
\label{eq:n11-g6-B-boundary}
\end{align}
The unique four-host clause of \cref{lem:six-conic-events} gives $W\le26$.
For $C=39$, \eqref{eq:n11-g6-B-boundary} instead requires $W\ge27$.
For $C=40$, both inequalities are equalities.  Thus $W=26$ and every
summand on the right of \eqref{eq:n11-g6-B-boundary} vanishes; in particular,
\[
 \sigma_6=\ell_{12}=\ell_{13}=\ell_{14}=\ell_{15}=\ell_{24}=0.
\]
The definition of $\sigma_6$ now reduces to
\[
 2\ell_{21}+4\ell_{22}+6\ell_{23}=30.
\]
Since $\ell_{24}=0$, the left side divided by two is precisely the incidence
number $J$ of outsiders with original lines carrying a pair of points of
$\Gamma$.  Hence
\[
 J=\ell_{21}+2\ell_{22}+3\ell_{23}=15,
\]
contrary to $J\le14$ in \cref{lem:six-conic-events}.

For footprint C at $C=39$, multiplying \eqref{eq:n11-g6-F2} by $24$ gives
\[
2=3\sigma_6+12\sigma_{10}+2\sigma_{17}
+12\ell_{12}+27\ell_{15}+2\ell_{22}.
\]
All quantities are nonnegative integers, so
\[
 \sigma_6=\sigma_{10}=\ell_{12}=\ell_{15}=0,
 \qquad \sigma_{17}+\ell_{22}=1.
\]
There are no three- or four-outsider blocks in footprint C.  Consequently
the tight $\sigma_6$ row gives
\[
 \ell_{21}=15-2\ell_{22},
\]
and the definition of the Melchior slack becomes
\[
 \sigma_{17}=52-(12\ell_{05}+3\ell_{21}+7\ell_{22})
              =7-12\ell_{05}-\ell_{22}.
\]
Adding $\ell_{22}$ and using $\sigma_{17}+\ell_{22}=1$ yields
$12\ell_{05}=6$, impossible for an integer block count.

It remains to exclude footprint C at $C=40$.  Here the left side of
$24\,$\eqref{eq:n11-g6-F2} is $26$.  Its coefficient $27$ excludes an
$\ell_{15}$ host, so the unique five-outsider host has type
$c_{15}$, $\ell_{05}$, or $c_{05}$.  Put
\[
 a=c_{12},\quad b=c_{21},\quad c=c_{22}=W,
 \qquad d=\ell_{12},\quad e=\ell_{21},\quad f=\ell_{22}.
\]
Only these six inner variables can occur.  The $\Gamma$--outsider line-pair
capacity, namely $\sigma_6\ge0$, gives
\begin{equation}
 d+e+2f\le15.\label{eq:n11-g6-row6}
\end{equation}
For each outsider, invert at that point.  The remaining ten points are
noncollinear: otherwise restoring the pivot would put all eleven original
points on one generalized block, contrary to the established size cap.
Thus \cref{thm:kelly-moser} supplies at least five ordinary
lines.  Summing over the five outsiders counts each generalized three-block
once for each outsider on it.  Since footprint C has no three-outsider
block, this gives the second structural row
\begin{equation}
 2a+b+2d+e\ge25.\label{eq:n11-g6-row26}
\end{equation}

Suppose first that the host is $c_{15}$.  Circle count and the two mixed
triple rows in \eqref{eq:n11-g6-triple-rows} become
\begin{equation}
 a+b+c=38,\qquad a+2c+d+2f=50,
 \qquad b+2c+e+2f=75.\label{eq:n11-g6-c15-rows}
\end{equation}
Substitution in \eqref{eq:n11-g6-row26} and
\eqref{eq:n11-g6-row6}, respectively, gives
\[
 175-6c-6f\ge25,\qquad 87-3c-2f\le15;
\]
equivalently, $c+f\le25$ and $3c+2f\ge72$.  Since
$3c+2f=c+2(c+f)$, we get $W=c\ge22$.  But this host contains a point of
$\Gamma$, so \cref{lem:six-conic-events} gives $W\le20$.

If the host is $\ell_{05}$, it already owns all ten outsider pairs in the
line-pair row, whence $d=f=0$.  The exact rows are
\[
 a+b+c=39,\qquad a+2c=60,
 \qquad b+2c+e=75.
\]
Thus $b=c-21$ and
\[
 J=e=96-3c=96-3W.
\]
The disjoint-host bound $W\le26$ now gives $J\ge18$, contrary to the global
bound $J\le14$ in \cref{lem:six-conic-events}.

Finally suppose that the host is $c_{05}$.  The three exact rows are
\begin{equation}
 a+b+c=38,\qquad a+2c+d+2f=60,
 \qquad b+2c+e+2f=75.\label{eq:n11-g6-c05-rows}
\end{equation}
Substitution in \eqref{eq:n11-g6-row26} and
\eqref{eq:n11-g6-row6} gives
\[
 195-6c-6f\ge25,\qquad 97-3c-2f\le15,
\]
or $c+f\le85/3$ and $3c+2f\ge82$.  Hence
$c\ge76/3$, and integrality together with the disjoint-host bound forces
$W=c=26$.  The circle row gives $a+b=12$, while
$\sigma_{10}\ge0$ (here $c_{03}=0$) says $2a+b\ge15$; hence $a\ge3$.
The last row of \eqref{eq:n11-g6-c05-rows} now yields
\[
 J=e+2f=75-b-2c=11+a\ge14.
\]
All five outsiders lie on one fixed proper circle, so the localized clause
of \cref{lem:six-conic-events} gives $J\le13$, the final contradiction.
Thus all footprints and both boundary values are excluded.
\end{proof}

\begin{proposition}\label{prop:n11-value}
One has $f(11)=41$.
\end{proposition}

\begin{proof}
Assume $C\le40$.  By \cref{lem:n11-n12-outer-router}, the largest proper
circle has size five or six.  The five-circle branch is excluded by
\cref{lem:n11-g5-low,lem:n11-C39,lem:n11-C40}, and the six-circle branch by
\cref{lem:n11-gamma6}.  Construction \ref{con:circle-centre} attains $41$.
\end{proof}

\subsection{Twelve points: the five-circle branch}

We first isolate the scalar spine.  Its role is only to route the proof;
all geometric equality cases will be discharged explicitly afterwards.

\begin{lemma}[Five-circle spine at twelve points]\label{lem:n12-g5-spine}
Suppose that $|P|=12$, the largest proper circle has size five, and
$C\le50$.  Put $B_s=C_s+L_s$ and
\[
 \sigma_p=d_3(p)-3-d_5(p)-2d_6(p),\qquad
 K=\sum_p\sigma_p.
\]
Then every block has size $3,4,5$, or $6$, every six-block is a line,
$B_5\le12$, and
\begin{align}
B_3+4B_4+10B_5+20B_6&=220,\label{eq:n12-g5-triples}\\
12B_4+35B_5+72B_6&=624-K,\label{eq:n12-g5-block}\\
4L&=256-4C-B_5-4B_6+K,\label{eq:n12-g5-lines}\\
9B_5&\ge1776-36C+5K+72B_6.\label{eq:n12-g5-added}
\end{align}
At every point,
\begin{equation}\label{eq:n12-g5-local}
d_3+3d_4+6d_5+10d_6=55,\qquad
3d_4+7d_5+12d_6=52-\sigma,\qquad d_5\le6.
\end{equation}
If $B_6=L_5=0$, write
\begin{equation}\label{eq:n12-q-definition}
d_3=7+3j,\quad \epsilon=\ell_4,\quad
u=5-\ell_3-\ell_4,\quad q=u+2j+2-d_5.
\end{equation}
Then $j,u,q\ge0$.  Moreover,
\begin{equation}\label{eq:n12-d5-caps}
j=0\Longrightarrow d_5\le4,\qquad
j=1\Longrightarrow d_5\le5,\qquad
j\ge2\Longrightarrow d_5\le6.
\end{equation}
In addition, the two local types
\begin{equation}\label{eq:n12-grid-forbidden-types}
(j,u,\epsilon,q,d_5)=(2,0,0,0,6),\qquad(2,1,0,1,6)
\end{equation}
do not occur.
\end{lemma}

\begin{proof}
The line cap in \cref{lem:n11-n12-outer-router} leaves only blocks of the
stated sizes, and a six-block cannot be a proper circle.  Triple ownership
gives \eqref{eq:n12-g5-triples}.  Summed Melchior after inversion gives
\[
 K=3B_3-36-5B_5-12B_6.
\]
Eliminating $B_3$ proves \eqref{eq:n12-g5-block}; adding
$C+L=\sum_sB_s$ proves \eqref{eq:n12-g5-lines}.  Finally, the summed
added-centre inequality is
\[
 5L_3+12L_4+20L_5+28L_6\le4C-124+C_5.
\]
Its left side is at least $5L+23B_6$, while $C_5\le B_5$; substitution of
\eqref{eq:n12-g5-lines} gives \eqref{eq:n12-g5-added}.

For a five-block $F$, let
$y_F=\mathbf1_F-\frac5{12}\mathbf1$.  These vectors lie in the
eleven-dimensional space $\mathbf1^\perp$, and distinct ones satisfy
\[
 \langle y_F,y_G\rangle=|F\cap G|-\frac{25}{12}\le-\frac1{12}.
\]
A nonzero obtuse family in $\R^{11}$ has at most twelve members, proving
$B_5\le12$.

The first local equality in \eqref{eq:n12-g5-local} partitions the pairs
of the other eleven points; eliminating $d_3$ gives the second.  If seven
five-blocks passed through $p$, deleting $p$ would leave seven four-sets
on eleven points, pairwise meeting in at most one point.  Their $28$
incidences have pair moment at least
\[
 6\binom32+5\binom22=23>\binom72,
\]
which is impossible.  Hence $d_5\le6$.

Assume now $B_6=L_5=0$.  Kelly--Moser after inversion and the congruence
in \eqref{eq:n12-g5-local} give $d_3=7+3j$ with $j\ge0$.  The arms of the
rich lines through $p$ are disjoint, so $u\ge0$.  The added-centre slack is
\begin{equation}\label{eq:n12-kappa-q}
 \kappa=3(u+j)-\epsilon-d_5\ge0.
\end{equation}
If $j=0$, then $\sigma\ge0$ gives $d_5\le4$.  To prove the middle cap and
record its two later equality adapters, suppose $d_5=6$ and invert at $p$.
The six resulting four-lines have multiplicities $m_x$ on the remaining
eleven selected points.  Pair ownership gives
\[
 \sum_xm_x=24,\qquad \sum_x\binom{m_x}{2}\le\binom62=15.
\]
Writing $z_x=m_x-2$, we have $\sum z_x=2$ and
\[
 \sum_x\binom{m_x}{2}=14+\frac12\sum_xz_x^2\ge15.
\]
Thus equality holds: the profile is $(3,3,2^9)$ and every pair of the six
lines meets at a selected point.  On each four-line,
$\sum_{x\text{ on the line}}(m_x-1)=5$, so it contains exactly one of the
two triple points.  The six lines are therefore two three-line pencils,
and their other nine intersections form a $3\times3$ grid.  Every further
three-point image line avoids the pencil centres by pair ownership and is
a grid transversal.  Every image of an original line through $p$ is a
two-secant through the inversion centre $o$; apart from the possible line
through the two pencil centres, it is an ordinary grid secant.

If $j=1$, \eqref{eq:n12-g5-local} gives $d_4=3$, hence three grid
transversals, contrary to \cref{lem:real-grid}.  This proves
$d_5\le5$.  For the first type in \eqref{eq:n12-grid-forbidden-types}, the
local equations give $(d_3,d_4,\ell_3)=(13,2,5)$; the two transversals
normalize the grid, while at least four ordinary grid secants concur at
$o$.  For the second they give $(13,2,4)$ and at least three such secants.
Both contradict the final clause of \cref{lem:real-grid}.  The last cap in
\eqref{eq:n12-d5-caps} is the already proved $d_5\le6$.

It remains to check $q\ge0$.  For $j=0$, the cap settles $u\ge2$, while
\eqref{eq:n12-kappa-q} gives $d_5\le3u$ for $u\le1$.  For $j=1$, the cap
settles $u\ge1$, and \eqref{eq:n12-kappa-q} settles $u=0$.  For $j\ge2$,
$d_5\le6\le u+2j+2$.  This proves every assertion.
\end{proof}

We shall use once the following parity consequence of the local line
arrangement.  After inversion and restoration of the centre, dualize the
twelve-point configuration.  The product of the twelve defining linear
forms two-colours the faces.  Counting edge--face incidences shows that the
two colour sums of $|F|-3$ are congruent modulo three.  Consequently the
total Melchior excess cannot equal one; in the notation above,
\begin{equation}\label{eq:n12-even-cell}
 \kappa_p\ne1\qquad\text{for every }p.
\end{equation}

\begin{lemma}[The nonnegative $q$-front]\label{lem:n12-g5-qfront}
Assume the hypotheses of \cref{lem:n12-g5-spine}, and in addition
$B_6=L_5=0$ and $K=B_5=k$.  Put $h=L_4$.  Then
\begin{equation}\label{eq:n12-q-sums}
 \sum j=2k-16,\qquad
 \sum u=3C-132-h,\qquad
 \sum q=3C-140-k-h.
\end{equation}
In particular $k\ge8$ and $k+h\le3C-140$.
\end{lemma}

\begin{proof}
Equations \eqref{eq:n12-g5-triples}--\eqref{eq:n12-g5-lines} give
\[
 B_3=12+2k,\qquad B_4=52-3k,\qquad L=64-C.
\]
Now sum the four definitions in \eqref{eq:n12-q-definition}, using
$\sum d_3=3B_3$, $\sum d_5=5k$, $\sum\ell_3=3(L-h)$, and
$\sum\epsilon=4h$.  This gives \eqref{eq:n12-q-sums}.  Its first and last
members are sums of nonnegative integers.
\end{proof}

\begin{lemma}[Twelve-point five-circle closure]\label{lem:n12-gamma5}
If twelve admissible points contain a proper five-circle and no larger
proper circle, then $C\ge51$.
\end{lemma}

\begin{proof}
Suppose $C\le50$.  From \eqref{eq:n12-g5-added} and $B_5\le12$,
\[
108\ge1776-36C+5K+72B_6.
\]
Thus $C\ge47$.  Moreover \eqref{eq:n12-g5-block} gives
$B_5\equiv K\pmod {12}$.  At $C=47$ the preceding inequality and
$1\le B_5\le12$ leave only
$K=B_6=0,B_5=12$, and \eqref{eq:n12-g5-lines} then gives $L=14$.  For
$r_p=\sum_{s=3}^6\ell_s(p)$, the pointwise added-centre row gives
\[
 r_p\le3+\left\lceil\frac{\sigma_p}{3}\right\rceil,\qquad
 \sum_pr_p\le36+K.
\]
But $\sum_pr_p\ge3L=42$, a contradiction.

At $C=48$, the inequality is
$48+5K+72B_6\le9B_5\le108$.  Hence $B_6=0$, $K\le12$, and the congruence
$B_5\equiv K\pmod {12}$ gives the complete list
\[
 (K,B_5,B_6,L)=(0,12,0,13),\qquad(12,12,0,16).
\]
In the first row $\sigma=0$ pointwise, so
\eqref{eq:n12-g5-local} modulo three and $d_5\le6$ give
$d_5\in\{1,4\}$.  Thus $\sum d_5\le48<60=5B_5$.  In the second row the
total added-centre slack is
\[
 Q=\sum_p\kappa_p=-7L_4-16L_5.
\]
Hence $L_4=L_5=0$, and \cref{lem:n12-g5-qfront} gives
$\sum q=144-140-12=-8$, impossible.

At $C=49$, the corresponding inequality is
$12+5K+72B_6\le9B_5\le108$, so $B_6\le1$.  If $B_6=1$, the congruence
leaves only $K=0,B_5=12$; if $B_6=0$, it leaves that endpoint and
$K=B_5=k$, $3\le k\le12$.  The two $K=0,B_5=12$ endpoints are
excluded by the same congruence and degree sum.  Every remaining endpoint
has
\[
 B_6=0,\qquad K=B_5=k,\quad3\le k\le12,\qquad L=15,
\]
and total added-centre slack
\[
 Q=k-3-7L_4-16L_5\ge0.
\]
Thus $L_5=0$.  The two nonnegative sums in
\eqref{eq:n12-q-sums} say simultaneously $k\ge8$ and $k+L_4\le7$, a
contradiction.

We turn to $C=50$.  Here
$-24+5K+72B_6\le9B_5\le108$, hence $B_6\le1$.  Splitting the single
congruence $B_5\equiv K\pmod {12}$ over $1\le B_5\le12$ gives exactly five
families:
\begin{equation}\label{eq:n12-C50-five-fronts}
\begin{array}{c|c|c|c}
B_6&K&B_5&L\\ \hline
0&k&k\ (1\le k\le12)&14\\
0&0&12&11\\
0&B_5+12&9\le B_5\le12&17\\
1&0&12&10\\
1&12&12&13.
\end{array}
\end{equation}
The two $K=0$ rows again fail the congruence degree sum.  In the last row,
$Q=-7L_4-16L_5$, so $L_3=12,L_6=1$ and $\kappa=0$ pointwise.  It follows
that $\sigma=1$ at every point.  Equation \eqref{eq:n12-g5-local} then
gives $d_5\in\{0,3\}$ on the six-line and
$d_5\in\{0,3,6\}$ off it, whence
$\sum d_5\le6\cdot3+6\cdot6=54<60$.

In the $L=17$ row, write $k=B_5$.  The total slack forces
$L_4=L_5=0$.  The definitions in \eqref{eq:n12-q-definition}, with the
corresponding shifted value of $K$, give
\[
 \sum j=2k-12,\qquad \sum u=9,\qquad \sum q=9-k.
\]
Thus $k=9$, $Q=0$, all $q$ and $\kappa$ vanish, and exactly three points have
$(j,u,d_5)=(2,0,6)$; all others have $(j,u,d_5)=(0,1,3)$.  At any of the
three former points $\epsilon=q=0$, so this is the first forbidden type in
\eqref{eq:n12-grid-forbidden-types}.  We are left with
\begin{equation}\label{eq:n12-C50-main-front}
 B_6=0,\qquad K=B_5=k,\qquad L=14.
\end{equation}

We next remove a possible five-line.  The total slack in
\eqref{eq:n12-C50-main-front} is
$6+k-7L_4-16L_5$.  Hence $L_5>0$ would force
$L_5=1,L_4=0$ and $k\in\{10,11,12\}$.  Put
\[
 z_p=\ell_5(p),\qquad d_3(p)=7+3j_p,\qquad
 u_p=5-\ell_3(p)-2z_p,\qquad
 \widehat q_p=u_p+2j_p+2+z_p-d_5(p).
\]
Here $z_p\in\{0,1\}$.  The eleven-point inversion link is noncollinear by
the established block caps, so Kelly--Moser and the link congruence give
$j_p\ge0$; the line-pencil partition gives $u_p\ge0$.  The degree caps in
\eqref{eq:n12-d5-caps} require only $B_6=0$, not $L_5=0$.  Indeed, for
$j_p=0$ they follow from $\sigma_p=4-d_5(p)\ge0$, and for $j_p\ge2$ from
the universal bound $d_5(p)\le6$.  If $j_p=1$ and $d_5(p)=6$, the six
four-point lines in the inversion link have equality profile
$(3,3,2^9)$, exactly as in the proof of \eqref{eq:n12-d5-caps}; they form
two three-line pencils and a real $3\times3$ grid.  The local equations give
$d_4(p)=3$, hence three grid transversals, contrary to
\cref{lem:real-grid}.  This argument is unchanged if one of the six base
lines is the image of the unique five-line.  Thus the three caps are valid
throughout the present branch.

Restored-centre Melchior gives the nonnegative slack
\[
 \widehat\kappa_p=3(u_p+j_p)+z_p-d_5(p)\ge0.
\]
If $j_p=0$, the degree-four cap settles $u_p\ge2$, while
$\widehat\kappa_p\ge0$ settles $u_p\le1$.  If $j_p=1$, the degree-five cap
settles every case except $u_p=z_p=0$, which is again covered by
$\widehat\kappa_p\ge0$.  If $j_p\ge2$, use $d_5(p)\le6$.  Consequently
$\widehat q_p\ge0$ for every point.  On the other hand,
\[
 \sum_pz_p=5,\quad \sum_pu_p=11,\quad
 \sum_pj_p=2k-16,\quad \sum_pd_5(p)=5k,
\]
and therefore
\[
 \sum_p\widehat q_p
 =11+2(2k-16)+24+5-5k=8-k<0.
\]
This contradiction proves $L_5=0$.

Set $h=L_4$ and retain $q$ from \eqref{eq:n12-q-definition}.  The exact
sums now are
\begin{equation}\label{eq:n12-C50-six-sums}
\sum j=2k-16,\quad \sum u=18-h,\quad
\sum q=10-k-h,\quad \sum\epsilon=4h,\quad
\sum\kappa=6+k-7h.
\end{equation}
Their nonnegativity leaves precisely
\[
 (h,k)=(0,8),(0,9),(0,10),(1,8),(1,9),(2,8).
\]
Eliminating $d_5$ between $q$ and \eqref{eq:n12-kappa-q} gives
\begin{equation}\label{eq:n12-low-defect-identities}
d_5=u+2j+2-q,\qquad
\kappa=2u+j+q-\epsilon-2.
\end{equation}
Using \eqref{eq:n12-even-cell}, the caps
\eqref{eq:n12-d5-caps}, and the two exclusions
\eqref{eq:n12-grid-forbidden-types}, these identities imply
\begin{equation}\label{eq:n12-low-defect-table}
\begin{array}{c|c}
\text{local condition}&\text{forced conclusion}\\ \hline
\epsilon=0,q=0&j=0\\
\epsilon=0,q=1&j\le1\\
q=0,\epsilon=1&(j,u,d_5,\kappa)=(1,1,5,0)\\
q=0,\epsilon=2&(j,u,d_5,\kappa)=(0,2,4,0).
\end{array}
\end{equation}
For completeness, in the first two rows $j\ge1$ reduces either to
$\kappa=1$ or to $(j,u,q,d_5)=(2,0,0,6)$ or $(2,1,1,6)$; the latter two
are \eqref{eq:n12-grid-forbidden-types}.  In the last two rows,
\eqref{eq:n12-low-defect-identities} and $\kappa\ge0$ leave the displayed
value directly.

Now $(0,10)$ contradicts $\sum j=4$ and the first row of
\eqref{eq:n12-low-defect-table}; $(0,9)$ contradicts $\sum j=2$ because
only the unique point with $q=1$ can have $j>0$; $(1,9)$ puts $j=1$ at all
four points of the four-line although $\sum j=2$; and $(1,8)$ would require
positive $q$ at all four line points although $\sum q=1$.  At $(2,8)$,
$j=q=0$ pointwise, so every point of either four-line must have
$\epsilon=2$; the two distinct lines would have the same support.

Only $(h,k)=(0,8)$ remains.  Here $j=\epsilon=0$ pointwise,
$\sum q=2$, and $\sum u=18$.  At least ten points have $q=0$; at each of
them \eqref{eq:n12-low-defect-identities} and $d_5\le4$ force
$u\in\{1,2\}$.  If all had $u=2$, their contribution to $\sum u$ would
already be twenty.  Hence some point has
\[
 (j,u,\epsilon,q,d_5,\kappa)=(0,1,0,0,3,0).
\]
Inversion, restoration of the centre, and duality turn this point into a
twelve-line arrangement with multiplicity vector
$(t_2,t_3,t_4)=(6,14,3)$ and distinguished word $3D+4T$.  This is the first
forbidden completion in \cref{lem:gallery-completions}.  The last endpoint,
and therefore the entire five-circle branch, is impossible.
\end{proof}

\subsection{Twelve points: the six-circle branch}

\begin{lemma}[Universal six-circle spine]\label{lem:n12-g6-spine}
Assume $|P|=12$, $C\le50$, and the largest proper circle has size six.  Put
\[
 m=B_6,\quad k=B_5,\quad h=L_4,\quad g=L_5,\quad r=L_6,\quad e=50-C,
\]
and define
\[
 \sigma_p=d_3(p)-3-d_5(p)-2d_6(p),\qquad
 \kappa_p=11+\sigma_p-3\ell_3(p)-4\ell_4(p)
             -5\ell_5(p)-6\ell_6(p).
\]
If $K=\sum_p\sigma_p$ and $Q=\sum_p\kappa_p$, then, for an integer
$s\ge0$,
\begin{align}
K&=k+12s,\label{eq:n12-g6-K}\\
B_3&=12+4m+2k+4s,&
B_4&=52-6m-3k-s,\notag\\
B&=64-m+3s,&
L&=14-m+3s+e,\label{eq:n12-g6-normalization}\\
Q&=6+9m+k-15s-9e-7h-16g-27r.\label{eq:n12-g6-Q}
\end{align}
Moreover $1\le m\le4$, $m+k\le12$, and
\begin{equation}\label{eq:n12-g6-gram-caps}
m=4\Rightarrow k=0,\qquad
m=3\Rightarrow k\le3,\qquad
m=2\Rightarrow k\le8,\qquad
m=1\Rightarrow k\le11.
\end{equation}
Pointwise,
\begin{equation}\label{eq:n12-g6-residues}
3d_4+7d_5+12d_6=52-\sigma,\qquad
\sigma\equiv1-d_5\pmod3,\qquad
\kappa\equiv\ell_5-d_5-\ell_4\pmod3,
\end{equation}
where $\sigma,\kappa\ge0$ and $\kappa\ne1$.
\end{lemma}

\begin{proof}
The selected circle gives $m\ge1$, and the outer router bounds every block
by six.  For a six-block $A$ and a five-block $F$, set
\[
y_A=\mathbf1_A-\frac12\mathbf1,\qquad
z_F=\mathbf1_F-\frac5{12}\mathbf1.
\]
Distinct maximal blocks meet in at most two points, so
\[
\begin{array}{c|ccc}
&\|\cdot\|^2&\text{same-size inner product}&
 \text{mixed inner product}\\ \hline
y&3&\le-1&\le-\frac12\\
z&\frac{35}{12}&\le-\frac1{12}&\le-\frac12.
\end{array}
\]
Thus all $m+k$ vectors form an obtuse family in $\mathbf1^\perp$, giving
$m+k\le12$.  In particular $k\le11$.  Triple ownership and summed
Melchior give $K\equiv k\pmod {12}$; hence
$K=k+12s$ with $s\ge0$.  Substitution in the triple equation and the
circle--line count gives \eqref{eq:n12-g6-normalization}; summing the
definition of $\kappa$ gives \eqref{eq:n12-g6-Q}.

For the six-block vectors alone,
\[
0\le\left\|\sum_Ay_A\right\|^2\le m(4-m),
\]
so $m\le4$.  More sharply, applying constant positive coefficients to the
two size classes majorizes their Gram form by
\begin{equation}\label{eq:n12-two-class-gram}
G(m,k)=
\begin{pmatrix}
m(4-m)&-mk/2\\
-mk/2&k(36-k)/12
\end{pmatrix}.
\end{equation}
Because the off-diagonal entries are nonpositive, copositivity forces the
relevant determinant to be nonnegative.  For $m=3,2,1$ the determinants
factor, respectively, as
\[
 \frac{k(36-10k)}4,\qquad
 \frac{k(36-4k)}3,\qquad
 \frac{k(36-2k)}4.
\]
Together with $m+k\le12$ this gives all bounds in
\eqref{eq:n12-g6-gram-caps}, except that $m=2$ initially permits $k=9$.
At that null endpoint equality forces every rich-block intersection to
have size two and
\[
 9(y_{A_1}+y_{A_2})+4\sum_Fz_F=0,\qquad
 4d_5(p)+9d_6(p)=24.
\]
The two six-blocks meet in two points, leaving eight points with $d_6=1$;
there $4d_5=15$, impossible.  When $m=4$, the first diagonal entry of
\eqref{eq:n12-two-class-gram} is zero, so its negative mixed entry forces
$k=0$.

Finally, the pair partition at a pivot and the definitions of the slacks
give \eqref{eq:n12-g6-residues}.  The assertion $\kappa\ne1$ is precisely
the even-cell argument preceding \cref{lem:n12-g5-qfront}; it depends only
on the parity of the twelve-line dual arrangement, not on the selected
circle size.
\end{proof}

\begin{lemma}[Ordinary-direction inequality]\label{lem:n12-direction}
Under the hypotheses of \cref{lem:n12-g6-spine}, suppose $L_6=0$.  At every
point lying on a six-block,
\begin{equation}\label{eq:n12-direction}
 \sigma_p\ge2d_5(p)+6d_6(p)-8.
\end{equation}
If also $d_5(p)=\ell_4(p)=\ell_5(p)=\ell_6(p)=\kappa_p=0$ and the branch
under consideration has $d_6(p)\le2$, then $d_6(p)=2$, $\sigma_p\ge4$,
and $\ell_3(p)\ge5$.
\end{lemma}

\begin{proof}
Invert at $p$, restore the inversion centre, and dualize the resulting
twelve-point configuration.  A chosen proper six-block through $p$ becomes
a five-point line not containing the restored centre, hence supplies a
fivefold vertex $v$ of the dual twelve-line arrangement.  Let $t_i$ be the
number of its $i$-fold vertices, let
\[
 \delta=\sum_{w\notin\operatorname{star}(v)}
              (\operatorname{mult}(w)-1)
        =t_2+2t_3+3t_4+4(t_5-1)-35,
\]
and let
\[
 R=\#\{\text{vertices on the distinguished line}\}
   =\ell_2(p)+\ell_3(p)+\ell_4(p)+\ell_5(p).
\]
The distinguished line misses $v$.  Its intersections with the five lines
of $\operatorname{star}(v)$ are therefore distinct, and every further
vertex on it contributes at least one unit to the off-star defect.  Thus
$\delta\ge R-5$.

For completeness, the inversion dictionary at $p$ is
\[
 \begin{aligned}
 t_2&=d_3-\ell_3+\ell_2,&
 t_3&=d_4-\ell_4+\ell_3,\\
 t_4&=d_5-\ell_5+\ell_4,&
 t_5&=d_6+\ell_5,
 \end{aligned}
\]
while the pointwise pair partitions are
\[
 \ell_2+2\ell_3+3\ell_4+4\ell_5=11,
 \qquad
 d_3+3d_4+6d_5+10d_6=55.
\]
(All local quantities in these two displays are evaluated at $p$.)
Substituting these identities and
$\sigma_p=d_3(p)-3-d_5(p)-2d_6(p)$ gives the exact integral identity
\begin{equation}\label{eq:n12-direction-defect}
 \delta-(R-5)
 =\frac{11+\sigma_p-2d_5(p)-6d_6(p)}3.
\end{equation}

Equality $\delta=R-5$ is impossible.  Indeed, equality says that every
off-star vertex is ordinary and lies on the distinguished line.  Send $v$
to infinity and dualize back: the five lines in its star become five marked
directions, and every remaining join passes through the point dual to the
distinguished line.  Deleting that point leaves six distinct real points
whose joins use only those five directions.  They are noncollinear.  A
sixfold vertex on the distinguished line would correspond before inversion
to a six-line, excluded by $L_6=0$; a sixfold vertex away from it would
correspond to a proper seven-circle, excluded by maximality of the selected
six-circle.  This contradicts \cref{thm:ungar-six}.  Hence the integral
left-hand side of \eqref{eq:n12-direction-defect} is at least one, and
\eqref{eq:n12-direction} follows.

For the final assertion, the same inversion dictionary gives
\[
 t_2=d_3-\ell_3+\ell_2
    =3+d_5+2d_6+\kappa+\ell_4+\ell_5+\ell_6.
\]
Under the clean hypotheses this is $t_2=3+2d_6$.  Kelly--Moser gives
$t_2\ge6$, so $d_6(p)\ge2$; the branch hypothesis makes $d_6(p)=2$.  Now
\eqref{eq:n12-direction} gives $\sigma_p\ge4$, and $\kappa_p=0$ reads
$11+\sigma_p=3\ell_3(p)$.
\end{proof}

\begin{lemma}[Four six-blocks]\label{lem:n12-g6-m4}
The branch $B_6=4$ is impossible when $C\le50$.
\end{lemma}

\begin{proof}
By \cref{lem:n12-g6-spine}, $k=0$.  Equality in
$\|\sum_Ay_A\|^2\le m(4-m)=0$ says that every pair of six-blocks meets in
two points and $\sum_Ay_A=0$; coordinatewise, $d_6(p)=2$ at all twelve
points.  Equation \eqref{eq:n12-g6-residues} gives
$\sigma_p=1+3j_p$, so $s\ge1$.
Two distinct six-lines meet in only one point, whereas every pair of the
four six-blocks meets in two; consequently $r\le1$.

If $r=0$, \cref{lem:n12-direction} gives $\sigma_p\ge4$ everywhere and
therefore $s\ge4$.  On the other hand
$Q=42-15s-9e-7h-16g\ge0$ gives $s\le2$.

If $r=1$, nonnegativity in \eqref{eq:n12-g6-Q} forces
\[
s=1,\quad e=h=g=Q=0,\quad C=50,\quad L_3=12,\quad L_6=1.
\]
The residues force $\sigma_p=1$ and $\kappa_p=0$ at every point.  At a
point off the six-line the exact local data are
\[
(d_3,d_4,d_5,d_6)=(8,9,0,2),\qquad
(\ell_2,\ell_3,\ell_4,\ell_5,\ell_6)=(3,4,0,0,0).
\]
After inversion and duality this is a twelve-line arrangement with
$(t_2,t_3,t_5)=(7,13,2)$ and distinguished word $3D+4T$, the second
forbidden completion in \cref{lem:gallery-completions}.
\end{proof}

\begin{lemma}[Three six-blocks]\label{lem:n12-g6-m3}
The branch $B_6=3$ is impossible when $C\le50$.
\end{lemma}

\begin{proof}
Write $k=B_5$.  The Gram cap gives $0\le k\le3$.  The three six-blocks
have total incidence eighteen.  Since their union has at most twelve
points,
\[
 \sum_p(d_6(p)-1)_+\ge6,
\]
whereas pairwise intersection at most two gives
$\sum_p\binom{d_6(p)}2\le6$.  Since
$\binom d2\ge(d-1)_+$, equality holds throughout and
\begin{equation}\label{eq:n12-m3-d6}
 d_6=(2^6,1^6).
\end{equation}
The scalar equation is
\[
 K=k+12s,\qquad
 Q=33+k-15s-9e-7h-16g-27r,
\]
so $s\in\{0,1,2\}$.

For $s=0$, the degree rows follow from one scalar defect.  For $k=2,3$ let
$T$ be the total deficit from intersection size two among all rich-block
pairs.  Expanding the centred norm gives
\[
 0\le3-\frac{k^2}{12}-2T,
\]
so the integer $T$ is at most one.  If $r_p=d_5(p)+d_6(p)$, then
$\sum r_p=18+5k$ and
\[
 \sum r_p^2=\sum r_p+2\bigl(2\tbinom{3+k}{2}-T\bigr).
\]
For $k=2$, the value $T=1$ would give $66$, below the balanced minimum
$68$ for total degree $28$ on twelve points; for $k=3$ it would give
$91<93$, the balanced minimum for total degree $33$.  Thus $T=0$.
Taking first and second moments now gives the following degree rows and
minimum residue costs:
\[
\begin{array}{c|c|c}
k&(d_5(p))_{p\in P}&\min\sum_p\sigma_p\\ \hline
0&(0^{12})&12\\
1&(1^5,0^7)&7\\
2&(2^2,1^6,0^4)&8\\
3&(2^6,1^3,0^3)&15.
\end{array}
\]
The first two rows are immediate.  In the last column we used the least
nonnegative value in the class $1-d_5\pmod3$; every entry is larger than
$K=k$, excluding $s=0$.

For $s=1,2$, nonnegativity of $Q$ forces $r=0$.  Summing
\eqref{eq:n12-direction} and using
$\sum d_5=5k$, $\sum d_6=18$ gives
\begin{equation}\label{eq:n12-m3-direction-sum}
 K\ge10k+12.
\end{equation}
At $s=1$ this forces $k=0$; then \eqref{eq:n12-m3-d6}, the direction
bound, and the residue give $K\ge6\cdot4+6\cdot1=30>12$.
At $s=2$, \eqref{eq:n12-m3-direction-sum} gives $k\le1$.  The case $k=0$
again costs at least thirty.  If $k=1$, the minimum sigma cost is twenty
five and equality forces the unique five-block to use five of the six
points with $d_6=1$.  Here $e=h=g=r=0$ and $Q=4$; those five points have
$\kappa\equiv-1\pmod3$, hence $Q\ge10$, a contradiction.
\end{proof}

\begin{lemma}[Two six-blocks]\label{lem:n12-g6-m2}
The branch $B_6=2$ is impossible when $C\le50$.
\end{lemma}

\begin{proof}
Put $k=B_5$; by \cref{lem:n12-g6-spine}, $0\le k\le8$ and
$s\in\{0,1,2\}$.  Let $I\le2$ be the intersection size of the two
six-blocks.

First take $s=0$.  If $r=0$, sum \eqref{eq:n12-direction} over the union
of the six-blocks.  If $A,B,D$ denote the sums of $d_5$ on the classes
$d_6=2,1,0$, respectively, then $D\le Ik$ and
\[
 k=K\ge2(A+B)+8I-24.
\]
Since $A+B=5k-D$, this gives, for $I=0,1,2$ respectively,
\begin{equation}\label{eq:n12-m2-s0-three}
 k\ge10k-24,\qquad k\ge8k-16,\qquad k\ge6k-8.
\end{equation}
All three fail for $k\ge3$.  If $k=0$, the sigma residue costs at least
twelve; if $k=1$, $(d_5)=(1^5,0^7)$ costs at least seven.  For $k=2$, the
five-block pair moment and the sigma budget force
\[
 d_5=(1^{10},0^2),\qquad \sigma=(0^{10},1^2).
\]
A point with $d_6=2$ violates \eqref{eq:n12-direction}, so $I=0$; then
$\sum d_5d_6=10>4k$, although the four mixed block pairs meet in at most
two points each.

If $s=0$ and $r=1$, \eqref{eq:n12-g6-Q} gives
$e=h=g=0$, $k\ge3$, and $Q=k-3$.  Counting triples with two points on the
six-line, while \eqref{eq:n12-g6-normalization} gives
\[
(C_3,C_4,C_5,C_6)=(9+2k,40-3k,k,1),\qquad L_3=11.
\]
Hence
\[
90\le C_3+2C_4+3C_5+4C_6=93-k,
\]
so $k=3$ and equality holds.  Every proper circle then meets the six-line
twice.  Among the twenty triples off the line, the three five-circles
cover at most three, the proper six-circle at most four, and the eleven
triple-lines at most eleven.  Thus at most $3+4+11=18$ can be owned, a
contradiction.

Next let $s=1$.  Here $r=0$ and
\begin{equation}\label{eq:n12-m2-s1-spine}
 K=k+12,\qquad Q=k+9-9e-7h-16g.
\end{equation}
When $g=0$, the residue for $\kappa$ implies pointwise
\begin{equation}\label{eq:n12-m2-defect}
 \kappa_p+2\binom{d_5(p)}2\ge2d_5(p)-2\ell_4(p).
\end{equation}
It is immediate for $d_5\ge3$ and follows from the three residues for
$d_5=0,1,2$.  Since
$\sum_p\binom{d_5(p)}2\le k(k-1)$, summing gives
\begin{equation}\label{eq:n12-m2-scalar-defect}
 2k^2-11k+9-9e+h\ge0.
\end{equation}
This removes $k=2,3,4$ in the core $e=g=0$ for every allowed $h$, and
$k=1,\ldots,5$ when $(e,h,g)=(1,0,0)$.

Here is the complete scalar routing, recorded to make exhaustiveness
transparent.  Since $k\le8$, \eqref{eq:n12-m2-s1-spine} gives
\[
\begin{array}{c|c}
\text{nonzero parameter}&\text{only values not already killed by
\eqref{eq:n12-m2-scalar-defect}}\\ \hline
g=1&(e,h,k)=(0,0,7)\\
e=1,h=1&(g,k)=(0,7)\\
e=1,h=0&g=0,\ k\in\{0,6,7,8\}\\
e=g=0&h\in\{0,1,2\},\ 0\le k\le8.
\end{array}
\]
Indeed $g\ge2$, $e\ge2$, or $h\ge3$ makes $Q<0$; in the first two rows
$k=8$ would give the forbidden value $Q=1$.  Thus $k=6,7,8$ below cover
every rich scalar row, $k=5$ covers the sole remaining middle gap, and the
clean argument for $k=0,1$ covers all tiny rows.

For $k\ge5$, the two six-blocks must meet twice.  Indeed $I=0$ would give
$\sum d_5d_6=5k>4k$.  If $I=1$, the three degree-class sums satisfy
\[
A+B+D=5k,\qquad2A+B\le4k,\qquad D\le k.
\]
Hence $D-A\ge k$, so necessarily $D=k,A=0,B=4k$.  Summing the direction
inequality on the ten middle points gives $k+12\ge2B-20$, which forces
$k\le4$.  Thus $I=2$ for
$k\ge5$.  On the classes of sizes $2,8,2$ we then have
\begin{equation}\label{eq:n12-m2-groups}
A+B+D=5k,\qquad D-A\ge k,\qquad B\le\frac{k+28}{2}.
\end{equation}
Let $\Phi_t(S)$ be the balanced minimum of
$\sum_{i=1}^t\binom{x_i}{2}$ over nonnegative integers of sum $S$.
Writing $S=qt+r$, $0\le r<t$, balancing gives
\[
 \Phi_t(S)=t\binom q2+rq,\qquad
 \Phi_t(S+1)-\Phi_t(S)=\floor{S/t}.
\]
Convexity and \eqref{eq:n12-m2-groups} give
\[
\begin{array}{c|c|c|c}
k&(A,B,D)&\Phi_2(A)+\Phi_8(B)+\Phi_2(D)&k(k-1)\\ \hline
8&(7,18,15)&70&56\\
7&(5,17,13)\text{ or }(6,16,13)&50&42\\
6&(4,16,10)&30&30.
\end{array}
\]
Here is a supporting-slope certificate for the three minima.  Write
$(A,B,D)=(A_0+a,B_0+b,D_0+c)$ at the displayed row, so
$a+b+c=0$.  For $k=8$, the active constraints are
$b\le0$ and $b+2c\ge0$; the three supporting slopes are $3,2,7$, and the
increase is at least
\[
 3a+2b+7c=-b+4c\ge0.
\]
For $k=7$, use the first displayed minimizer.  The constraints are
$b\le0$ and $b+2c\ge-1$, while slopes $2,2,6$ give an increase at least
$4c$.  Thus $c\ge0$; when $c=0$, only $b=0,-1$ are possible, giving exactly
the two displayed minimizers.  For $k=6$, the constraints are
$b\le1$ and $b+2c\ge0$.  Slopes $2,2,5$ give an increase at least $3c$.
If $c=0$, the only point besides the displayed minimizer has
$(a,b,c)=(-1,1,0)$ and objective $31$; if $c<0$, the two constraints
conflict.  This proves the table without an integer search.
The first two rows are impossible.  Equality at $k=6$ forces
$d_5=(2^{10},5^2)$; every point then has minimum sigma two, so
$K\ge24>18$.  This also disposes of the only scalar five-line case
$(g,k)=(1,7)$ and the only positive-$e$ or positive-$h$ cases not already
removed by \eqref{eq:n12-m2-scalar-defect}.

For $k=5$, necessarily $e=g=0$.  Set
\[
G_p=\kappa_p+2\binom{d_5(p)}2-2d_5(p)+2\ell_4(p)\ge0.
\]
Equations \eqref{eq:n12-m2-s1-spine} and the pair moment give
$\sum G_p\le4+h\le6$.  If $h\le1$, any $d_5\ge4$ has $G_p\ge6$, exceeding
the relevant bound; if $h=2$, then $Q=0$ and the residue improves this to
$G_p\ge8$.  Hence $d_5\le3$ everywhere.  In
\eqref{eq:n12-m2-groups}, $D\le6$ and $D-A\ge5$, so $A+D\le7$ and
$B\ge18>(5+28)/2$, impossible.

Only $k=0,1$ remain.  If $h=1$, every point of the four-line contributes
at least two to $Q$, whereas \eqref{eq:n12-m2-s1-spine} gives $Q=2$ or
$3$.  In the clean core $e=g=h=0$, the $k=0$ row has at least nine points
with $d_5=\kappa=0$, and the $k=1$ row has seven; the clean clause of
\cref{lem:n12-direction} would put all of them in the at most two-point
intersection of the six-blocks.  The sole extra row
$(e,h,g,k)=(1,0,0,0)$ has all twelve points clean and is identical.
This exhausts $s=1$.

Finally let $s=2$.  The scalar formula forces
\[
e=h=g=r=0,\qquad k\in\{6,7,8\},\qquad Q=k-6.
\]
The value $k=7$ has forbidden total slack one.  For $k=8$, summing
\eqref{eq:n12-direction} on the six-block union gives, according as
$I=1$ or $2$, $K\ge8k-16=48$ or $K\ge6k-8=40$, while $K=32$;
$I=0$ already violates the mixed moment.  For $k=6$, $Q=0$ makes every
$d_5$ a multiple of three.  Writing $d_5=3a$ and using the pair moment,
\[
 \sum a=10,\qquad
 \frac{9\sum a^2-30}{2}\le30,
\]
so $d_5=(3^{10},0^2)$.  The direction and residue costs are made explicit
by
\[
\begin{array}{c|ccc}
&d_6=0&d_6=1&d_6=2\\ \hline
d_5=3&1&4&10\\
d_5=0&1&1&4.
\end{array}
\]
If the two six-blocks meet in $I=0,1,2$ points, the $d_6$ classes have
sizes $I,12-2I,I$ in the order $2,1,0$.  Were all twelve points of
five-degree three, their minimum sigma cost would be $48+3I$.  Placing the
two zero-degree points optimally saves at most
$6I+3(2-I)=6+3I$.  Thus the actual cost is at least $42>30=K$, and
$s=2$ is impossible as well.
\end{proof}

\begin{lemma}[One six-block]\label{lem:n12-g6-m1}
The branch $B_6=1$ is impossible when $C\le50$.
\end{lemma}

\begin{proof}
Let $\Gamma$ be the unique six-block and put $k=B_5$.  It is a proper
circle, so $r=0$.  The scalar spine becomes
\begin{equation}\label{eq:n12-m1-spine}
K=k+12s,\qquad Q=15+k-15s-9e-7h-16g,\qquad s\in\{0,1\}.
\end{equation}
Write $A=\sum_{p\in\Gamma}d_5(p)$.  Every five-block has an outsider
footprint of size three, four, or five.  Distinct footprint families own
disjoint outsider triples.  A size-five footprint is unique and excludes
all size-four footprints; otherwise the complements of the size-four
footprints are disjoint pairs, hence form a matching of size at most three.
Consequently there is $u\in\{0,1,2,3\}$ such that
\begin{equation}\label{eq:n12-m1-footprint}
A=2k-u,\qquad \sum_{p\notin\Gamma}d_5(p)=3k+u.
\end{equation}

At a point of $\Gamma$, the direction inequality and the sigma residue
give the lower envelope
\begin{equation}\label{eq:n12-m1-envelope}
\sum_{p\in\Gamma}\sigma_p\ge f(A),\qquad
f(A)=\begin{cases}6-A,&A\le6,\\2A-12,&A\ge6.
\end{cases}
\end{equation}
For $k\ge5$, equations \eqref{eq:n12-m1-footprint}--
\eqref{eq:n12-m1-envelope} give $K\ge4k-18$.  Hence
\begin{equation}\label{eq:n12-m1-kcaps}
s=0\Rightarrow k\le6,\qquad s=1\Rightarrow k\le10.
\end{equation}

The two boundary values are already contradictory.  At $(s,k)=(0,6)$,
equality forces $A=9,u=3$ and uses the entire sigma budget on $\Gamma$.
The six outsider degrees, all $1$ modulo three, sum to $21$ and have pair
moment at least $30$; the six positive degrees on $\Gamma$, of sum nine,
add at least three.  This exceeds $k(k-1)=30$.  At $(s,k)=(1,10)$ one has
$A=17,u=3$, and discrete convexity gives
\[
 \sum_p\binom{d_5(p)}2
 \ge\Phi_6(17)+\Phi_6(33)=16+75>90=k(k-1).
\]

We next remove five-lines.  The local ordinary-point identity is
\begin{equation}\label{eq:n12-m1-t2}
t_2(p)=3+d_5(p)+2\mathbf1_{p\in\Gamma}+\kappa_p
       +\ell_4(p)+\ell_5(p).
\end{equation}
If $g=1$, necessarily $s=0$.  Set
\[
b_p=\frac{d_5(p)+\ell_4(p)+\kappa_p-\ell_5(p)}3.
\]
The residue makes $b_p$ integral, and Kelly--Moser applied to
\eqref{eq:n12-m1-t2} gives $b_p\ge1$, except possibly at the at most two
points of $\Gamma$ on the five-line.  Thus $\sum b_p\ge10$, while
\[
 \sum b_p=2k-2-3e-h.
\]
Together with $k\le6$, this forces $k=6,e=h=0$, the boundary value already
excluded.  Two five-lines make $Q<0$.  Hence
\begin{equation}\label{eq:n12-m1-gzero}g=0.\end{equation}

Define the integer and the convex energy
\[
a_p=\frac{d_5(p)+\ell_4(p)+\kappa_p}{3},\qquad
H(d)=\frac{(d-4)(d-5)}2=\binom d2-4d+10.
\]
By \eqref{eq:n12-m1-t2} and Kelly--Moser,
\begin{equation}\label{eq:n12-m1-a}
a_p\ge1,\qquad
\sum_pa_p=2k+5-5s-3e-h.
\end{equation}
If $N_1=|\{p:a_p=1\}|$, then
\begin{equation}\label{eq:n12-m1-N1}
N_1\ge19+5s+3e+h-2k.
\end{equation}
The five-block pair moment gives
\begin{equation}\label{eq:n12-m1-energy}
E:=\sum_p\bigl(\sigma_p+H(d_5(p))\bigr)
 \le k^2-20k+120+12s.
\end{equation}

There is one apparent exception to the useful pointwise lower bound.  If
$p\notin\Gamma$ and
\begin{equation}\label{eq:n12-typeA-local}
(d_5,\ell_4,\ell_5,\kappa,\sigma)=(3,0,0,0,1),
\end{equation}
then inversion and duality give
$(t_2,t_3,t_4)=(6,14,3)$ with distinguished word $3D+4T$.  The first part
of \cref{lem:gallery-completions} excludes this Type-A pivot.  In its
absence, $a_p=1$ implies
\begin{equation}\label{eq:n12-m1-local-energy}
 \sigma_p+H(d_5(p))\ge5.
\end{equation}
Indeed $d_5\le1$ gives $H\ge6$; $d_5=2$ forces
$(\ell_4,\kappa)=(1,0)$ and $\sigma\ge2$; and $d_5=3$ forces
$\ell_4=\kappa=0$, where \cref{lem:n12-direction} gives $\sigma\ge4$ on
$\Gamma$ and \eqref{eq:n12-typeA-local} is the only smaller off-circle
residue.  Combining \eqref{eq:n12-m1-N1}--
\eqref{eq:n12-m1-local-energy} yields the single inequality
\begin{equation}\label{eq:n12-m1-master}
 (k-5)^2\ge13s+15e+5h.
\end{equation}

For $s=0$, \eqref{eq:n12-m1-a} gives $2k\ge7+3e+h$.  With the boundary
$k=6$ gone, \eqref{eq:n12-m1-master} leaves only
$k\in\{4,5\}$ and $e=h=0$.  Let $z$ count the $a=1$ points with
$(d_5,\kappa)=(3,0)$.  Every other $a=1$ point has energy at least six, so
\[
 z\ge6N_1-E.
\]
For $k=4$, \eqref{eq:n12-m1-N1}--\eqref{eq:n12-m1-energy} give
$N_1\ge11,E\le56,z\ge10$; for $k=5$ they give
$N_1\ge9,E\le45,z\ge9$.  Each of these $z$ points has $\sigma\ge4$ by
the exclusion of Type A, contradicting $K=k$.

For $s=1$, \eqref{eq:n12-m1-a} gives $k\ge6$, and the boundary $k=10$ is
gone.  Inequality \eqref{eq:n12-m1-master} leaves only
\[
k=9,\qquad e=h=g=0,\qquad K=21,\qquad Q=9.
\]
Here $\sum a_p=18$.  Equations \eqref{eq:n12-m1-N1}--
\eqref{eq:n12-m1-energy} force exactly six points with $a=1$ and six with
$a=2$.  Let $q_i$ be the sum of $\kappa$ on the class $a=i$; then
$q_1+q_2=9$.  Since $d_5=3i-\kappa$, one checks
\[
a=1:\ 2H(d_5)\ge2+5\kappa,\qquad
a=2:\ 2H(d_5)\ge2-\kappa;
\]
the difference in either inequality is $\kappa(\kappa-2)\ge0$.
Consequently
\[
2\sum H(d_5)\ge15+6q_1,\qquad
\sum H(d_5)\le k(k-1)-60=12.
\]
If $q_1>0$, then $q_1\ge2$ and these bounds conflict.  Thus $q_1=0$;
all six $a=1$ points have $(d_5,\kappa)=(3,0)$ and, with Type A excluded,
contribute at least $24>K$ to the sigma sum.  This closes $s=1$ and the
one-six-block branch.
\end{proof}

\begin{lemma}[Twelve-point six-circle closure]\label{lem:n12-gamma6}
If twelve admissible points have largest proper circle of size six, then
$C\ge51$.
\end{lemma}

\begin{proof}
Assume $C\le50$.  By \cref{lem:n12-g6-spine},
$B_6\in\{1,2,3,4\}$.  These four values are excluded respectively by
\cref{lem:n12-g6-m1,lem:n12-g6-m2,lem:n12-g6-m3,lem:n12-g6-m4}.
\end{proof}


\begin{proposition}\label{prop:n12-value}
One has $f(12)=51$.
\end{proposition}

\begin{proof}
Assume $C\le50$.  By \cref{lem:n11-n12-outer-router}, the largest proper
circle has size five or six.  These alternatives are excluded by
\cref{lem:n12-gamma5,lem:n12-gamma6}.  Construction
\ref{con:circle-centre} attains $51$.
\end{proof}

\begin{theorem}\label{thm:finite-n10-n12}
The three remaining finite values are
\[
 f(10)=33,\qquad f(11)=41,\qquad f(12)=51.
\]
\end{theorem}

\begin{proof}
Combine \cref{prop:n10-value,prop:n11-value,prop:n12-value}.
\end{proof}

\section{The thirteen-point case and the direct large range}
\label{sec:n13-infinity}

Write
\[
  F(n)=1+\binom{n-1}{2}-\floor{\frac{n-1}{2}}.
\]
We finish the finite range at $n=13$ and then prove the lower bound
$C\geq F(n)$ directly for every $n\geq14$.  The large-range argument is
not recursive: it uses neither a value at a preceding integer nor a
distinguished base value.

For a point set $P$, let $C_s$ and $L_s$ be the numbers of maximal proper
circle blocks and maximal line blocks of size $s$, respectively, and put
$B_s=C_s+L_s$ and $C=\sum_s C_s$.  Every triple has a unique maximal
generalized-block owner.  We shall repeatedly use
\begin{align}
 \mathsf T&:=\sum_{s\geq3}\binom{s}{3}B_s=\binom n3,
                                                        \label{eq:large-T}\\
 \mathsf P&:=\sum_{s\geq3}s(4-s)B_s\geq3n,             \label{eq:large-P}\\
 \mathsf D&:=\sum_{s\geq3}s(s-4)C_s
       +\sum_{s\geq3}2s(s-2)L_s\leq n(n-4).           \label{eq:large-D}
\end{align}
Indeed, inversion at a selected point followed by Melchior's inequality
gives
\[
 \sigma_p=d_3(p)-3-\sum_{s\geq5}(s-4)d_s(p)\geq0.
\]
Summing this inequality gives \eqref{eq:large-P}.  Applying Melchior once
more after restoring the centre of inversion gives
\(\sum_s s\ell_s(p)\leq n-1+\sigma_p\); summing over $p$ and eliminating
the $\sigma_p$ gives \eqref{eq:large-D}.  These are geometric consequences
of the real projective-plane inequality of Melchior~\cite{Melchior1940},
not merely formal identities.

\subsection{The complete thirteen-point dispatch}

Suppose, towards a contradiction, that $|P|=13$ and $C\leq60$.  Let $m$
be the largest size of a proper circle block.

\begin{lemma}[Outer routing at thirteen points]
\label{lem:n13-outer}
Under the preceding assumptions, $m\in\{5,6\}$.  Moreover every line block
has size at most six.
\end{lemma}

\begin{proof}
If a line has $s$ selected points and $q=13-s$, the rich-line pencil gives
\begin{equation}
 q\binom{s}{2}-\binom q2\floor{\frac s2}.           \label{eq:n13-line-pencil}
\end{equation}
For $7\leq s\leq12$ its minimum is $66$.  Thus no such line occurs, while
$s=13$ is excluded by noncollinearity.
For a proper $m$-circle, the corresponding Bonferroni pencil is
\begin{equation}
 1+q\left(\binom m2-\floor{\frac m2}\right)
   -\binom q2\floor{\frac m2},\quad q=13-m.          \label{eq:n13-circle-pencil}
\end{equation}
For $m=7,8,9,10,11,12$ this equals, respectively,
\[
 64,81,105,106,96,61,
\]
so $m\geq7$ is impossible.  The case $m=13$ is excluded by the standing
non-concyclicity hypothesis.

It remains to exclude $m\leq4$.  Besides \eqref{eq:large-T} and
\eqref{eq:large-P}, global Melchior for line blocks gives
\[
 \mathsf G:=3L_3+\sum_{s\geq4}
 \left(\binom s2+s-3\right)L_s\leq\binom{13}{2}-3.
\]
Since line sizes are at most six, the fixed combination
\[
 \frac14\mathsf T+\frac7{60}\mathsf P-\frac15\mathsf G
\]
has coefficients
\[
\begin{array}{c|cc|rrrr}
 &C_3&C_4&L_3&L_4&L_5&L_6\\ \hline
 &3/5&1&0&-2/5&-29/60&0.
\end{array}
\]
It is therefore at most $C$, whereas its numerical lower bound is
\[
 \frac14\binom{13}{3}+\frac7{60}(3\cdot13)
 -\frac15\left(\binom{13}{2}-3\right)
 =\frac{1221}{20}>60.
\]
Thus $m\leq4$ is impossible as well.
\end{proof}

\subsubsection{The six-circle branch}

Fix a six-point circle $\Gamma$ and let $X=P\setminus\Gamma$, so
$|X|=7$.  If $e=xy\in\binom X2$, let $w_e$ be the number of proper
circles through $x,y$ and two points of $\Gamma$, and put
\begin{equation}
 W=\sum_{e\in\binom X2}w_e
   =\sum_j\binom j2c(2,j).                          \label{eq:n13-W}
\end{equation}
Here $c(i,j)$ counts nonselected circle blocks meeting $\Gamma$ in $i$
points and $X$ in $j$ points; $l(i,j)$ denotes the analogous number of
line blocks.
For either class $Y\in\{\Gamma,X\}$, define its pivot row by
\begin{equation}
 M_Y=\sum_B |B\cap Y|(4-|B|),                       \label{eq:n13-class-pivot}
\end{equation}
where the sum is over all maximal blocks, including $\Gamma$.
Summing the pointwise Melchior slack over $p\in Y$ gives
\[
 M_Y=3|Y|+\sum_{p\in Y}\sigma_p.
\]
Consequently
\begin{equation}
 M_\Gamma\geq18,\quad M_X\geq21.                    \label{eq:n13-class-pivot-bounds}
\end{equation}

\begin{lemma}[The opening inequality]
\label{lem:n13-gamma6-opening}
If $C\leq60$, then $W\geq48$.
\end{lemma}

\begin{proof}
For thirteen points, \eqref{eq:large-D} reads $\mathsf D\leq117$.
Let $O=3C_3$ be ordinary-circle incidence.  For every pivot $p$, the
inverted arrangement in \cref{cor:pivot-lenchner} has
$q=13-1=12\neq6$ lines.
Thus at least three ordinary proper circles pass through each pivot.
Summing over the thirteen pivots, and using
\eqref{eq:n13-class-pivot-bounds}, gives
\[
 O\geq39,\quad M_X\geq21.
\]
Weight the four triple partitions, in order of their number of
$\Gamma$-points, by
\[
 (141,340,539,141).
\]
Add $79O+199M_X$ and subtract $1344C+414W$.  For a nonselected block
write $g=|B\cap\Gamma|$, $x=|B\cap X|$, and $s=g+x$.  Its coefficient is
\begin{align*}
 A(g,x)={}&141\binom x3+340g\binom x2
  +539\binom g2x+199x(4-s)\\
 &+237\mathbf{1}_{\{B\text{ is a }3\text{-circle}\}}
  -1344\mathbf{1}_{\{B\text{ is a circle}\}}\\
 &-414\mathbf{1}_{\{B\text{ is a circle and }g=2\}}\binom x2.
\end{align*}
Direct factorisation, recorded in
\cref{subsec:n13large-certificates}, gives $A(g,x)\leq123D_B$ for every
geometrically possible block, including equality at the selected circle.
Since the four triple totals are $35,126,105,20$, summing gives
\begin{equation}
 123\mathsf D\geq33810-414W.                       \label{eq:n13-opening-dual}
\end{equation}
Together with $\mathsf D\leq117$, this first yields $W\geq47$.

At $W=47$, retaining every nonnegative slack in the same calculation gives
\[
 R+79(O-39)+199(M_X-21)+1344(60-C)\leq39,
\]
where every positive local residual contributing to $R$ is at least $96$.
Thus all four terms vanish.  For the zero-residual circle types
\[
 c(0,3),c(1,2),c(1,5),c(2,1),c(2,2),c(2,3),
\]
write their multiplicities, in order, as $(a,b,q,d,e,f)$.  The three
zero-residual line types and their multiplicities are
\[
\begin{array}{c|ccc}
\text{type}&l(0,3)&l(1,2)&l(2,1)\\ \hline
\text{multiplicity}&p_0&p_1&p_2.
\end{array}
\]
The circle count, ordinary-circle count, $W$, and the three triple
partitions reduce to
\begin{gather*}
 a+b+d=13,\quad q+e+f=46,\quad e+3f=47,\quad q=2f-1,\\
 p_0=45-a-21f,\quad p_1=42-b-20f,\quad p_2=11-d+3f.
\end{gather*}
Consequently
\[
 M_X=3a+2b+d+3p_0+2p_1+p_2-10q-3f=240-123f.
\]
The equality $M_X=21$ would force $f=73/41$, a contradiction.
\end{proof}

We now invoke the six-marked-conic interface
\cref{lem:six-conic-events}.  In the present notation, its first clause
gives $w_e\leq3$ and at most four projective involution signatures for the
\emph{full} edges $w_e=3$.  Each signature support is a matching in $X$,
and at most three full-signature centres lie on one trace line.  A pair of
support edges of the same signature is a \emph{repetition event}.
The trace dictionary \cref{lem:circle-lift-trace} and the second clause of
\cref{lem:six-conic-events} give every event a unique maximal
line-or-circle host.  If that host has $b$ outsider points, its event
capacity is
\begin{equation}
 k_0(b)=\min\left\{2\binom b4,
       3\binom{\floor{b/2}}2\right\};\quad
 k_0(4),k_0(5),k_0(6)=2,3,9.                        \label{eq:n13-host-cap}
\end{equation}
This is exactly the capacity \eqref{eq:conic-event-cap}.  The host cannot be
a line meeting $\Gamma$ or a circle meeting $\Gamma$ once.  A host circle
meeting $\Gamma$ in two points is admissible only for the signature
containing that marked pair.  These assertions follow by uniqueness of the
lift plane in \cref{lem:circle-lift-trace} and by the fact that two proper
circles sharing three points coincide.

\begin{lemma}[Capacity-gap lemma]
\label{lem:n13-capacity-gap}
Put $r=W-48$.  Then $0\leq r\leq6$.  If $R_{\rm rep}$ is the number of
repetition events, there is a nonnegative integer
\[
 E=2h+3k+9A_6+2\ell+2g,
\quad
 (h,k,A_6,\ell,g)=(c(0,4),c(0,5),c(0,6),l(0,4),c(2,4)),
\]
such that
\begin{equation}
 R_{\rm rep}\leq E,\quad 4E\leq6r+7+8A_6,\quad A_6=0.
                                                               \label{eq:n13-capacity-gap}
\end{equation}
\end{lemma}

\begin{proof}
By \cref{lem:n13-gamma6-opening}, $W\geq48$.  If $F_3$ denotes the number
of full edges, then
\[
 W\leq2\binom72+F_3.
\]
The four signature supports are matchings of size at most three, so
$F_3\leq12$ and hence $W\leq54$.  This proves $0\leq r\leq6$.

Two possible host types must first be removed explicitly.  Write
$\lambda_5=l(0,5)$ and $\lambda_6=l(0,6)$.  On a line of type $l(0,x)$,
the second-dual residual is $x(x-3)(14-x)$, hence it is $90$ at $x=5$
and $144$ at $x=6$.  Before forced-zero variables are suppressed, the
full second added-centre slack inequality therefore contains
\[
 90\lambda_5+144\lambda_6+\text{other nonnegative terms}
 \leq12+6r\leq48.
\]
Hence $\lambda_5=\lambda_6=0$.  Host uniqueness and admissibility now leave
only
\[
 c(0,4),\ c(0,5),\ c(0,6),\ l(0,4),\ c(2,4).
\]
Their capacities from \eqref{eq:n13-host-cap} are respectively
$2,3,9,2,2$, so this exhaustive list gives $R_{\rm rep}\leq E$.

For the numerical inequality, introduce
\[
 \begin{gathered}
 y=c(1,3),\ z=c(1,4),\ f=c(2,3),\quad
 w=l(1,3),\ t=l(2,2),\ p=l(2,3),\\
 d=60-C,\quad j=C_3-13,\quad
 u=M_X-21,\quad v=M_\Gamma-18,
 \end{gathered}
\]
all nonnegative integers.  The second added-centre calculation has the
exact retained-slack budget
\begin{align}
 &12h+26k+36A_6+3y+5z+2f+40\ell+31w+22t+48p
       +6u+3v+36d\leq12+6r.                       \label{eq:n13-second-budget}
\end{align}
Let $N=j+y+z$, $a=y+z$, and define
\begin{align*}
 K_0={}&N+2u+v+4h+8k+12A_6+10\ell+8d+4t+7w,\\
 Q_0={}&N+u+4h+8k+12A_6+8\ell+8d+2t+5w+7p.
\end{align*}
The eight conservation rows displayed in
\cref{subsec:n13large-certificates} give an integer $\mu\geq0$ and
$\delta=7\mu-4j$ satisfying
\begin{align}
 K_0+7j-2r&=14\mu,& Q_0+4\mu-2j&\leq3+2r,
                                                               \label{eq:n13-mu}\\
 3\delta&\leq6-k-5\ell-6d-5t-5w-24p.             \label{eq:n13-delta}
\end{align}
They also give, with
\[
 J=21+18r-(24g+24h+36k+84A_6+24\ell),
\]
the exact identity
\begin{align}
7J={}&49f+12j-60\delta+14a+7y+63u+14v+105k
 +168\ell+252d-161p+42t+189w.                    \label{eq:n13-J}
\end{align}
 Its right side is nonnegative.  If $p>0$, \eqref{eq:n13-delta} gives
\[
 \delta\leq2-8p,\quad
 -60\delta-161p\geq319p-120\geq199>0.
\]
Thus the only apparently negative pair of terms is already positive before
the remaining nonnegative terms are restored.
If $p=0$ and $\delta\leq0$, this is immediate.  For $\delta=1$, the
congruence $\delta=7\mu-4j$ gives $j\equiv5\pmod7$, hence
$12j-60\geq0$.  For $\delta=2$, \eqref{eq:n13-delta} forces
$k=\ell=d=t=w=0$ and $j\equiv3\pmod7$.  If $j\geq10$, then
$12j-120\geq0$.  For $j=3$ one has $\mu=2$, and
\eqref{eq:n13-mu} gives
\[
 a+u+4h+12A_6\leq2r-2.
\]
The two eliminated conservation identities then read $v\geq6-u$ and
\[
 f=(y-a)-2\bigl(v-(6-u)\bigr).
\]
Since $f\geq0$ and $y\leq a$, equality holds throughout:
$f=0$, $y=a$, and $v=6-u$.  The right side of
\eqref{eq:n13-J} is therefore $21a+49u\geq0$.  Hence $J\geq0$ in every
case.

Since
\[
 12E-24A_6=24g+24h+36k+84A_6+24\ell,
\]
$J\geq0$ is precisely the middle inequality of
\eqref{eq:n13-capacity-gap}.

It remains to remove $A_6$.  A first fixed added-centre row gives
\[
 2388A_6\leq453+414r,
\]
and hence $A_6=0$ for $0\leq r\leq4$.  For $r=5,6$, the number of full
edges is at least $6+r\geq11$.  Four matching supports, each of size at
most three, must therefore all have size at least two.  If a six-outsider
circle existed, every other event host would share at least three of its
six outsider points and hence would coincide with it.  All four signature
centres would then lie on its trace line, contradicting the three-centre
cap.  Thus $A_6=0$ also for $r=5,6$.
\end{proof}

The number $F_3$ of full outsider edges satisfies
\(W\leq2\binom72+F_3\), so $F_3\geq6+r$.  Balancing these edges among four
matching supports of capacity three gives
\begin{equation}
 R_{\rm rep}\geq\rho(r),\quad
\begin{array}{c|rrrrrrr}
r&0&1&2&3&4&5&6\\ \hline
\rho(r)&2&3&4&6&8&10&12.
\end{array}                                        \label{eq:n13-rho}
\end{equation}
Together, \eqref{eq:n13-capacity-gap} and \eqref{eq:n13-rho} immediately
exclude $r=0,4,5,6$.  At $r=1$, \eqref{eq:n13-second-budget} forces
$k=A_6=\ell=0$, so $E=2(h+g)$ is even; the capacity gap gives $E\leq3$,
and hence $E\leq2<\rho(1)$.  Only $r=2,3$ remain.

\begin{lemma}[Three-residue terminal]
\label{lem:n13-three-residue}
The equality layers $r=2$ and $r=3$ are empty.
\end{lemma}

\begin{proof}
For $r=2$, \eqref{eq:n13-rho} and the capacity gap give
\[
 4\leq E\leq\floor{\frac{19}{4}}=4;
\]
for $r=3$ they give
\[
 6\leq E\leq\floor{\frac{25}{4}}=6.
\]
Thus $E=4,6$, respectively.  Since $A_6=0$, the definition of $J$ becomes
$J=21+18r-12E$, and hence $J=9$ for $r=2$ and $J=3$ for $r=3$.

We next justify the reduction of variables.  At $r=2$, the right side of
\eqref{eq:n13-second-budget} is $24$; its coefficients force
\[
 k=\ell=d=w=p=0.
\]
At $r=3$, its right side is $30$, so
$\ell=d=w=p=0$ and $k\leq1$.  But here
\[
 E-3k=2(h+g)
\]
is even, whereas $E=6$; hence $k$ cannot equal one and again $k=0$.
Together with $A_6=0$, this leaves precisely the nonnegative integers
$a,y,u,v,h,t,f,j,\mu$, with $0\leq y\leq a$.  They satisfy
\begin{align}
 \delta&=7\mu-4j,                                  \label{eq:n13-res-delta}\\
 v&=2\delta+2r-a-2u-4h-4t,\nonumber\\
 f&=2j+\mu+y-v+4h-5t-2r,                           \label{eq:n13-res-vf}\\
 a+u+4h+2t+4\mu-j&\leq3+2r,                       \label{eq:n13-res-budget}\\
 J&=7a+48h+90j-123\mu-24r-9t+19u+8y.             \label{eq:n13-res-J}
\end{align}
The capacity inequality gives $\delta\leq2$.  Negative $\delta$ is
impossible: for $r=3$, the term $-60\delta$ in the uneliminated identity
\eqref{eq:n13-J} already exceeds $7J=21$; for $r=2$, the value
$\delta=-1$ forces $j\equiv2\pmod7$, and then
$-60\delta+12j\geq84>63=7J$.  More negative values only increase this
lower bound.  Thus $\delta\in\{0,1,2\}$.

Moreover,
\[
 4\mu-j=\frac{4\delta+9j}{7}.
\]
Dropping the other nonnegative terms in \eqref{eq:n13-res-budget} gives
$4\delta+9j\leq49$ for $r=2$ and at most $63$ for $r=3$.  The congruence
$j\equiv5\delta\pmod7$ therefore leaves
\begin{equation}
 (\delta,j,\mu)\in\{(0,0,0),(1,5,3),(2,3,2)\}.    \label{eq:n13-three-residues}
\end{equation}
For completeness, the sole apparent extra triple is $(0,7,4)$ when
$r=3$, but then the nonnegative right side of \eqref{eq:n13-J} contains
$12j=84>21$.

For $(0,0,0)$, the two inequalities $v,f\geq0$ give
\[
 a+2u+4h+4t\leq2r,\quad
 a+y+2u+8h-t\geq4r.
\]
Since $y\leq a$, subtraction yields
\[
 2r\leq y+4h-5t\leq a+4h-5t\leq2r-2u-9t.
\]
Hence $u=t=0$, $y=a=2r-4h$, and
$J=6r-12h$.  The pairs $(r,J)=(2,9),(3,3)$ would respectively make
$h=1/4,5/4$.

For $(1,5,3)$, the two relevant rows become
\begin{align*}
 a+u+4h+2t&\leq2r-4,\\
 J+24r-81&=7a+48h-9t+19u+8y.
\end{align*}
When $r=2$, the budget makes every variable zero, whereas the left side
of the identity is $-24$.  When $r=3$, the budget is two: for $t=0$ the
right side is nonnegative, and for $t=1$ it equals $-9$, while the left
side equals $-6$.

Finally, for $(2,3,2)$,
\begin{align*}
 a+u+4h+2t&\leq2r-2,\\
 J+24r-24&=7a+48h-9t+19u+8y. 
\end{align*}
For $r=2$ the left side is $33$ and the budget is two.  Thus $h=0$.
If $t=1$, the right side is $-9$; if $t=0$, parity forces
$a+u=1$, and the right side is at most $19$.

For $r=3$ the left side is $51$ and the budget is four.  The value $h=1$
gives $48$, while $t=2$ gives $-18$ and $t=1$ gives at most $29$.
Thus $h=t=0$.  Parity leaves $a+u\in\{1,3\}$.  The first value gives at
most $19$; for the second, substituting $u=3-a$ gives
$4(2y-3a)=-6$.  This is impossible.
\end{proof}

\begin{proposition}
\label{prop:n13-gamma6}
If a thirteen-point set has a largest proper circle of size six, then
$C\geq61$.
\end{proposition}

\begin{proof}
Assume $C\leq60$.  The capacity-gap lemma places
$r=W-48$ in $\{0,\ldots,6\}$ and excludes $r=0,1,4,5,6$, while
\cref{lem:n13-three-residue} excludes $r=2,3$.
\end{proof}

It remains to close the five-circle value left by \cref{lem:n13-outer}.

\begin{lemma}[The five-circle transfer]
\label{lem:n13-gamma5}
If a thirteen-point set has no proper circle larger than five, then
$C\geq61$.
\end{lemma}

\begin{proof}
Assume $C\leq60$.  By \eqref{eq:n13-line-pencil}, every line has size at
most six, so only $B_3,B_4,B_5,B_6$ occur and $B_6=L_6$.  Set
\begin{align*}
 K&=\sum_{p\in P}\sigma_p,\quad
 Q=3B_3-5B_5-12B_6=39+K,\\
 S&=3L_3+4L_4+10L_5+12L_6,\\
 A&=5L_3+12L_4+20L_5+28L_6.
\end{align*}
The triple partition and the two applications of Melchior described above
give
\begin{align}
 B_3+4B_4+10B_5+20B_6&=286,                       \label{eq:g5-triples}\\
 12B_4+35B_5+72B_6&=819-K,                         \label{eq:g5-block}\\
 A&\leq4C-169+C_5,\quad A\geq\frac53S,             \label{eq:g5-A}\\
 24\cdot286+27Q&\leq35(S+3C).                     \label{eq:g5-transfer}
\end{align}
The last inequality is coefficientwise:
\[
\begin{array}{c|rrrr}
s&3&4&5&6\\ \hline
\text{left side}&105&96&105&156\\
\text{right side, circle}&105&105&105&--\\
\text{right side, line}&105&140&350&420.
\end{array}
\]
The dash records $C_6=0$ in this branch.  Thus, on every type that can
occur, no separation or optimisation theorem is hidden in
\eqref{eq:g5-transfer}.

From \eqref{eq:g5-block}, $C_5\leq B_5\leq23$.  Equations
\eqref{eq:g5-A}--\eqref{eq:g5-transfer} now give
\[
 377-5C\leq A\leq4C-146,
\]
and hence $C\geq523/9$, so only $C=59,60$ remain.

For these two endpoints let $L=\sum_{s=3}^6L_s$.  Exact elimination from
\eqref{eq:g5-triples}--\eqref{eq:g5-A} gives
\begin{align}
 4L&=325-4C-B_5-4B_6+K,                            \label{eq:g5-L}\\
 9B_5&\geq2301-36C+5K+72B_6.                      \label{eq:g5-B5}
\end{align}
Combining \eqref{eq:g5-block}, $B_4\geq0$, and \eqref{eq:g5-B5} yields
\[
 \begin{array}{ll}
 C=59:&23K+396B_6\leq147,\\
 C=60:&46K+792B_6\leq609.
 \end{array}
\]
Thus $B_6=0$.  Reducing \eqref{eq:g5-block} modulo $12$, then using
\eqref{eq:g5-B5} and $B_4\geq0$, gives exactly
\begin{equation}
 (B_3,B_4,B_5,B_6)=(48+2K, 7-3K, 21+K, 0),
 \quad K\in\{0,1,2\}.                              \label{eq:g5-three-endpoints}
\end{equation}
At every point $p$, pair partition and the definition of $\sigma_p$ give
\begin{equation}
 7d_5(p)+3d_4(p)=63-\sigma_p.                      \label{eq:g5-point}
\end{equation}

If $K=0$, \eqref{eq:g5-point} and $d_4(p)\leq B_4=7$ force
$d_4(p)\in\{0,7\}$.  Since $\sum_p d_4(p)=4B_4=28$, four points lie in
all seven four-blocks, making two distinct maximal blocks share four
points.  If $K=1$, twelve zero-slack points have $(d_5,d_4)=(9,0)$ and
the remaining point has $(8,2)$; hence $\sum_p d_5(p)=116$, rather than
$5B_5=110$.  Finally, if $K=2$, any point with positive slack must have
$d_4(p)\geq2$, while \eqref{eq:g5-three-endpoints} gives $B_4=1$.
All three cases are impossible.
\end{proof}

\begin{theorem}\label{thm:n13-value}
The thirteen-point value is
\[
 f(13)=61.
\]
\end{theorem}

\begin{proof}
For a hypothetical set with $C\leq60$, \cref{lem:n13-outer} leaves only
$m=5$ and $m=6$.  These are excluded by \cref{lem:n13-gamma5} and
\cref{prop:n13-gamma6}, respectively.  Hence $C\geq61=F(13)$.  The
circle-with-centre construction from the foundational section has exactly
$F(13)$ proper circles, proving equality.
\end{proof}

\subsection{A uniform half-cap and one master for
  \texorpdfstring{$n\geq15$}{n >= 15}}

\begin{lemma}[Uniform half-cap]
\label{lem:large-half-cap}
Let $n\geq15$ and suppose $C\leq F(n)-1$.  Every maximal line or proper
circle block then has size at most
\[
 H=\floor{\frac n2}.
\]
\end{lemma}

\begin{proof}
Let a block have size $s$, put $q=n-s\geq1$, and write
$h=\floor{s/2}$.  The rich-circle and rich-line pencils are
\begin{align}
 R(n,s)&=1+q\left(\binom s2-h\right)-\binom q2h,   \label{eq:large-rich-circle}\\
 Q(n,s)&=q\binom s2-\binom q2h=R(n,s)-1+qh.        \label{eq:large-rich-line}
\end{align}
We show $R(n,s)\geq F(n)$ whenever $s\geq H+1$.  The following four rows
cover all parity choices; in each row $c\geq0$.
\[
\begin{array}{c|c|c|l}
s&q&\text{parameter relation}&R(n,s)-F(n)\\ \hline
2a&2b&a=b+1+c&
(4b-2)c^2+(6b^2-3b)c+2b^3-5b^2+b\\
2a&2b+1&a=b+1+c&
b\{4c^2+(6b-1)c+2b^2-5b-3\}\\
2a+1&2b&a=b+c&
(4b-2)c^2+(6b^2-7b+2)c+2b^3-7b^2+4b\\
2a+1&2b+1&a=b+1+c&
b\{4c^2+(6b+3)c+2b^2-b-1\}.
\end{array}                                                    \tag{*}
\]
The condition $n\geq15$ is, row by row,
\(2b+c\geq7,6,7,6\).  For $b\geq3$ all relevant coefficients in the
first three rows are nonnegative, and the fourth row is nonnegative for
every $b\geq1$.  The remaining possibilities reduce to
\[
\begin{array}{c|cc}
&b=1&b=2\\ \hline
\text{row 1}&2c^2+3c-2\ (c\geq5)&6c^2+18c-2\ (c\geq3)\\
\text{row 2}&4c^2+5c-6\ (c\geq4)&2(4c^2+11c-5)\ (c\geq2)\\
\text{row 3}&2c^2+c-1\ (c\geq5)&6c^2+12c-4\ (c\geq3),
\end{array}
\]
all strictly positive.  In the two odd-$q$ rows, $b=0$ gives equality
$R=F$.  Thus a circle of size $s\geq H+1$ already determines at least
$F(n)$ circles.  Equation \eqref{eq:large-rich-line} and $qh\geq1$ give
the same conclusion for a line.  Either conclusion contradicts
$C\leq F(n)-1$.
\end{proof}

Under the cap $s\leq M$, inversion leaves at most $M-1$ collinear images.
For $M=\floor{n/2}$,
\[
 M-1\leq\frac{2(n-1)}3.
\]
The Langer--de Zeeuw inequality \cref{thm:langer}, applied after each
inversion and summed over all centres, therefore gives
\begin{equation}
 \mathsf L:=\sum_{s\geq3}s(s-1)B_s
 \geq\frac{n(n-1)(n+2)}3.                          \label{eq:large-L}
\end{equation}

\begin{lemma}[The $S_M$ master]
\label{lem:large-SM}
Assume every maximal block has size at most $M\geq3$ and that
\eqref{eq:large-L} holds.  Then
\begin{equation}
 C\geq A(n,M):=
 \frac{\binom n3}{2M}+\frac{n(M+3)}{6M}
 +\frac{(M-2)n(n-1)(n+2)}{36M}-\frac{n(n-4)}9.
                                                               \label{eq:large-A}
\end{equation}
\end{lemma}

\begin{proof}
Consider the single functional
\begin{equation}
 \mathsf S_M=\frac{\mathsf T}{2M}
 +\frac{M+3}{18M}\mathsf P
 +\frac{M-2}{12M}\mathsf L-\frac19\mathsf D.      \label{eq:large-SM}
\end{equation}
If $k_C(s)$ and $k_L(s)$ are its coefficients on an $s$-circle and an
$s$-line, respectively, direct factorisation gives
\begin{align}
 1-k_C(s)&=\frac{(M-s)(s-4)(s-3)}{12M},            \label{eq:large-circle-slack}\\
 k_L(s)&=\frac{s(s-3)(-7M+3s-12)}{36M}.            \label{eq:large-line-slack}
\end{align}
Thus $k_C(3)=k_C(4)=1$, $k_C(s)\leq1$ for $5\leq s\leq M$, and
$k_L(s)\leq0$ for every $3\leq s\leq M$.  Hence
$\mathsf S_M\leq C$.  Substitution of
\eqref{eq:large-T}, \eqref{eq:large-P}, \eqref{eq:large-D}, and
\eqref{eq:large-L}, with attention to the negative coefficient of
$\mathsf D$, gives \eqref{eq:large-A}.
\end{proof}

Taking $M=\floor{n/2}$ in \eqref{eq:large-A} yields
\begin{align}
 A(2r,r)-F(2r)&=\frac{(2r-1)(r^2-9r+13)}9,
                                                               \label{eq:large-even-gap}\\
 A(2r+1,r)-F(2r+1)&=
 \frac{4r^4-32r^3+37r^2+9}{18r}.                  \label{eq:large-odd-gap}
\end{align}
The first is positive for $r\geq8$; after $r=8+t$ its quadratic factor is
$t^2+7t+5$.  In the second, after $r=7+t$, the numerator is
\[
 4t^4+80t^3+541t^2+1302t+450,
\]
so it is positive for $r\geq7$.  Together with
\cref{lem:large-half-cap,lem:large-SM}, this proves $C\geq F(n)$ for
every $n\geq15$.

\subsection{The exact dichotomy at fourteen points}

\begin{proposition}[Seven-circle dichotomy]
\label{prop:n14-dichotomy}
Every admissible fourteen-point set determines at least $73$ proper
circles.
\end{proposition}

\begin{proof}
Assume $C\leq72$.  The pencil bounds
\eqref{eq:large-rich-circle}--\eqref{eq:large-rich-line} show that proper
circle blocks have size at most seven and line blocks have size at most
six.  If $C_7=0$, all blocks have size at most six.  Every inverted line
then contains at most five of the thirteen image points, and
$5\leq2\cdot13/3$, so \cref{thm:langer} supplies \eqref{eq:large-L}.
The $S_M$ master with $M=6$ gives
\[
 C\geq A(14,6)=\frac{3899}{54}>72.
\]

It remains to consider a selected seven-circle $\Gamma$ with seven-point
complement $X$.  For every block put
$g=|B\cap\Gamma|$, $x=|B\cap X|$, and $s=g+x$, and define
\begin{equation}
 W=\sum_{\substack{B\text{ a proper circle}\\g=2}}\binom x2.
                                                               \label{eq:n14-W}
\end{equation}
The four triple totals according to their number of $\Gamma$-points are
$\binom7i\binom7{3-i}$.  Put
\[
 M_\Gamma=\sum_B g(4-s),\quad M_X=\sum_B x(4-s),
\]
where the sums include the selected block.  Summed pivot Melchior over each
of the two seven-point classes gives $M_\Gamma,M_X\geq21$, while
\eqref{eq:large-D} gives $\mathsf D\leq140$.

For every block set
\begin{align*}
 A_B={}&6\binom x3+9g\binom x2+12\binom g2x+6\binom g3
 +(3g+6x)(4-s)\\
 &-36\mathbf{1}_{\{B\text{ is a circle}\}}
 -6\mathbf{1}_{\{B\text{ is a circle and }g=2\}}\binom x2.
\end{align*}
For a nonselected block one has $g\leq2$, and the residual $4D_B-A_B$
factors as
\begin{equation}
\begin{array}{c|ccc}
&g=0&g=1&g=2\\ \hline
\text{circle}&(x-3)(-x^2+10x-12)&
 (5-x)(x-2)(2x-3)/2&(4-x)(x-2)(x-1)\\
\text{line}&x(14-x)(x-3)&
 (x-2)(-2x^2+21x+17)/2&(x-1)(-x^2+7x+12).
\end{array}                                                     \label{eq:n14-residuals}
\end{equation}
Under the block caps every line residual is nonnegative; a negative circle
residual has magnitude at most $3\binom x3$.  Outsider triples have unique
owners, so the total correction is at most $3\binom73$.  The selected
seven-circle contributes the remaining defect $27$.  Summing the block
inequality and using $C\leq72$ gives
\begin{align*}
 6W\geq{}&3\binom73+9\cdot7\binom72
 +12\binom72\cdot7+6\binom73+3\cdot21+6\cdot21\\
 &-36\cdot72-4\cdot140-27=412.
\end{align*}
Hence $W\geq69$.  On the other hand, for a fixed outsider pair the
$\Gamma$-pairs used by circles counted in $W$ are disjoint.  It contributes
at most $\floor{7/2}=3$, and therefore
\[
 W\leq3\binom72=63,
\]
a contradiction.
\end{proof}

\begin{theorem}\label{thm:large-range}
For every integer $n\geq14$,
\[
 f(n)=F(n)=1+\binom{n-1}{2}-\floor{\frac{n-1}{2}}.
\]
\end{theorem}

\begin{proof}
The lower bound for $n=14$ is \cref{prop:n14-dichotomy}.  For $n\geq15$,
\cref{lem:large-half-cap} supplies the common block cap, Langer supplies
\eqref{eq:large-L}, and the one functional in \cref{lem:large-SM} gives
the strict contradictions \eqref{eq:large-even-gap} and
\eqref{eq:large-odd-gap} under $C\leq F(n)-1$.  The foundational
circle-with-centre construction attains $F(n)$ for every $n\geq14$.
\end{proof}

\subsection{Fixed-certificate appendix and verification boundary}
\label{subsec:n13large-certificates}

For clarity, we record which parts of this section are fixed arithmetic
certificates and which parts are mathematical inputs.

\paragraph{The opening certificate.}
In \cref{lem:n13-gamma6-opening}, the residual
$123D_B-A(g,x)$ has the following line residuals:
\begin{align*}
 R_L(0,x)&=\frac{x(x-3)(890-47x)}2,\\
 R_L(1,x)&=\frac{(x-2)(-47x^2+597x+246)}2,\\
 R_L(2,x)&=\frac{x(x-1)(304-47x)}2.
\end{align*}
The circle residuals are
\begin{align*}
 R_C(0,x)&=237+\frac{(x-3)(-47x^2+644x-738)}2
                 -237\mathbf{1}_{\{x=3\}},\\
 R_C(1,x)&=\frac{(5-x)(47x^2-210x+390)}2
                 -237\mathbf{1}_{\{x=2\}},\\
 R_C(2,x)&=\frac{(x-2)(x-3)(284-47x)}2
                 -237\mathbf{1}_{\{x=1\}}.
\end{align*}
It is enough---and slightly stronger than the outer cap requires---to use
the intervals
\[
\begin{array}{c|ccc}
&g=0&g=1&g=2\\ \hline
\text{line}&3\leq x\leq7&2\leq x\leq6&1\leq x\leq5\\
\text{circle}&3\leq x\leq6&2\leq x\leq5&1\leq x\leq4.
\end{array}
\]
Each displayed expression is nonnegative on its indicated interval.  Thus
the opening is a coefficientwise proof, not a conclusion inferred from
sampled incidence patterns.

\paragraph{Affine certificate.}
For reference, retain the zero- and positive-slack variables
\[
 c_{03},c_{12},c_{15},c_{21},c_{22},l_{03},l_{12},l_{21}
\]
in addition to the variables in \cref{lem:n13-capacity-gap}.  With every
row written as left side minus right side, the eight conservation equations
are
\begin{align*}
R_C={}&1+c_{03}+h+k+A_6+c_{12}+y+z+c_{15}+c_{21}+c_{22}+f+g-(60-d),\\
R_O={}&c_{03}+c_{12}+c_{21}-(13+j),\\
R_0={}&c_{03}+4h+10k+20A_6+y+4z+10c_{15}+f+4g
       +l_{03}+4\ell+w+p-35,\\
R_1={}&c_{12}+3y+6z+10c_{15}+2c_{22}+6f+12g
       +l_{12}+3w+2t+6p-126,\\
R_2={}&c_{21}+2c_{22}+3f+4g+l_{21}+2t+3p-105,\\
R_W={}&c_{22}+3f+6g-(48+r),\\
R_X={}&3c_{03}-5k-12A_6+2c_{12}-4z-10c_{15}+c_{21}-3f-8g\\
 &\quad+3l_{03}+2l_{12}+l_{21}-3p-(21+u),\\
R_\Gamma={}&-12+c_{12}-z-2c_{15}+2c_{21}-2f-4g
       +l_{12}+2l_{21}-2p-(18+v).
\end{align*}
All eight equal zero.

\smallskip
\noindent\emph{The second-dual inequality.}
For a nonselected block $B$, put $g=|B\cap\Gamma|$, $x=|B\cap X|$ and
$s=g+x$.  Let $D_B=s(s-4)$ for a circle and $D_B=2s(s-2)$ for a line.
Weight the four triple rows by $6,9,12,6$, add $6M_X+3M_\Gamma$, and
subtract $36C+6W$.  The coefficient of $B$ is
\begin{align*}
 A_2(g,x)={}&6\binom x3+9g\binom x2+12\binom g2x
 +(6x+3g)(4-g-x)\\
 &-36\mathbf1_{\{B\text{ is a circle}\}}
 -6\mathbf1_{\{B\text{ is a circle and }g=2\}}\binom x2.
\end{align*}
For $g=0,1,2$, the residual $4D_B-A_2(g,x)$ is
\[
\begin{array}{c|ccc}
&g=0&g=1&g=2\\ \hline
\text{line}&
x(x-3)(14-x)&
(x-2)(-2x^2+21x+17)/2&
(x-1)(-x^2+7x+12)\\
\text{circle}&
(x-3)(-x^2+10x-12)&
(5-x)(x-2)(2x-3)/2&
(4-x)(x-2)(x-1).
\end{array}
\]
The line ranges are respectively $3\leq x\leq7$, $2\leq x\leq6$,
$1\leq x\leq5$, and the circle ranges are $3\leq x\leq6$,
$2\leq x\leq5$, $1\leq x\leq4$.  Every displayed factor is nonnegative;
the selected circle separately has equality $48=4\cdot12$.  Summing and
using $\mathsf D\leq117$ gives exactly
\begin{align}
 B_0:={}&12h+26k+36A_6+3y+5z+2f+40\ell+31w+22t+48p\notag\\
        &\quad+6u+3v+36d-(12+6r)\leq0.             \label{eq:n13-B0}
\end{align}
The zero-residual types are precisely
$c_{03},c_{12},c_{15},c_{21},c_{22},l_{03},l_{12},l_{21}$ and
$g=c(2,4)$; every omitted type has nonnegative residual and may safely be
discarded from the retained budget.
This proves the budget used in \eqref{eq:n13-second-budget}
coefficientwise.

\smallskip
\noindent\emph{The integral coordinate.}
In the row order $(R_C,R_O,R_0,R_1,R_2,R_W,R_X,R_\Gamma)$, define
\begin{align*}
 F_0={}&14f-\bigl(180A_6+24\ell-30r+36d+60h+36j+106k-140p\\
      &\qquad\quad-66t+2u-13v-21w+15y+z\bigr),\\
 G_0={}&336g-\bigl(294+342r-1716A_6-240\ell-276d-516h-276j\\
      &\qquad\quad-976k+1302p+618t-20u+123v+189w-87y+25z\bigr).
\end{align*}
Their fixed row certificates are
\begin{align}
 F_0={}&-36R_C+36R_O-6R_0+9R_1+24R_2-30R_W+2R_X-13R_\Gamma=0,
                                                               \label{eq:n13-F0}\\
 G_0={}&276R_C-276R_O+60R_0-83R_1-226R_2+342R_W
         -20R_X+123R_\Gamma=0.                                \label{eq:n13-G0}
\end{align}
Put $N=j+y+z$ and
\begin{align*}
 B_1={}&9N+36h+72k+108A_6+76\ell+49w+22t+49p\\
       &\quad+72d+11u+2v-(21+18r),\\
 H={}&2j+y-v+4h+7k+12A_6+\ell+2d-10p-5t-2w-2r.
\end{align*}
The right side subtracted from $14f$ in \eqref{eq:n13-F0} has the exact
decomposition
\[
 (K_0+7j-2r)+14H.
\]
Consequently
\[
 K_0+7j-2r=14(f-H).
\]
Set $\mu=f-H$ and $E_\mu=K_0+7j-2r-14\mu$.  Then $\mu$ is an integer,
$E_\mu=0$, and $\mu\geq0$: indeed the left side above is at least
$-2r\geq-12$, so it cannot be a negative multiple of $14$.

The remaining budget transformations are the fixed identities
\begin{align}
 7B_0-4B_1-F_0&=0,                                  \label{eq:n13-B01}\\
 B_1+21+18r&=2K_0+7Q_0,                             \label{eq:n13-B1KQ}\\
 B_1&=7\{Q_0+4\mu-2j-(3+2r)\}+2E_\mu.              \label{eq:n13-B1mu}
\end{align}
Thus $B_0\leq0$ and $F_0=E_\mu=0$ imply $B_1\leq0$ and precisely the two
relations in \eqref{eq:n13-mu}.

\smallskip
\noindent\emph{The $\delta$ and $J$ certificates.}
Let
\[
 \delta=7\mu-4j,\qquad
 L=3+2r-(Q_0+4\mu-2j)\geq0.
\]
Direct substitution in the displayed definitions gives the fixed identity
\begin{align}
 S&:=6-k-5\ell-6d-5t-5w-24p-21\mu+12j\notag\\
  &=f+2L+z\geq0,                                    \label{eq:n13-S}
\end{align}
and hence
\[
 3\delta=6-k-5\ell-6d-5t-5w-24p-S.
\]
This is \eqref{eq:n13-delta}.  The same substitution gives the two exact
equalities
\begin{align}
 v={}&2\delta+2r-a-2u-4h-8k-12A_6-10\ell-8d-4t-7w,
                                                               \label{eq:n13-v-cert}\\
 f={}&2j+\mu+y-v+4h+7k+12A_6+\ell+2d
       -10p-5t-2w-2r.                              \label{eq:n13-f-cert}
\end{align}

For the last elimination set $G=24g$ and
\begin{align*}
 G_*={}&21+42r+123\mu-228A_6-105\ell-90d-72h-83j-140k\\
      &\quad+93p+9t-19u-48w-8y-7N.
\end{align*}
Then the first fixed certificate is
\begin{equation}
 G_0-14(G-G_*)+123E_\mu=0,                          \label{eq:n13-Gstar}
\end{equation}
so $G=G_*$.  Let $E_\delta=\delta-(7\mu-4j)$, let $E_v$ be the left side
minus the right side of \eqref{eq:n13-v-cert}, and put
\begin{align*}
 E_{7f}=7f-\bigl(&18j+15\delta+7(y-a)-14v-7k-63\ell-42d\\
                 &-70p-63t-14u-63w\bigr).
\end{align*}
Equation \eqref{eq:n13-f-cert} gives $E_{7f}=0$.  Finally, let $E_J$ be
$7J$ minus the right side of \eqref{eq:n13-J}.  The terminal fixed
identity is
\begin{equation}
 E_J+7(G-G_*)-84E_v+7E_{7f}-123E_\delta=0.          \label{eq:n13-EJ}
\end{equation}
Every term except $E_J$ has just been shown to vanish, proving
\eqref{eq:n13-J}.  This explicit chain establishes
\eqref{eq:n13-mu}--\eqref{eq:n13-J}; no multiplier is chosen by
optimisation or incidence search.

\paragraph{What was checked arithmetically.}
During preparation, exact-arithmetic checks independently replayed the
opening residuals, the $W=47$
identities, all eight affine rows and their fixed eliminations, the seven
values of $\rho$, the capacity-gap conversion, the three-residue
calculations, the four half-cap polynomials, the coefficients of $S_M$, and
the fourteen-point residual table.  These checks manipulate integers, rationals,
and symbolic polynomials only.

Such arithmetic checks do not prove Melchior's theorem, Langer's inequality, the pivot form
of Lenchner's theorem~\cite{Lenchner2007}, the inversion dictionary, or the
real six-marked-conic signature and host lemmas.  Those are the stated
geometric inputs to the written proofs.  The classification-free
$\Gamma_5$ transfer in \cref{lem:n13-gamma5} is checked directly through
its displayed blockwise coefficients and endpoint equations.  The distributed
exact-arithmetic supplement
\path{supplement/verify_n13_gamma5_global_closure.py} checks an independent
global closure of that branch.  In every case the arithmetic is
corroborative, not a substitute for the theorem-level arguments above.


\begin{proof}[Proof of \cref{thm:main}]
The exact values for \(4\le n\le9\) are
\cref{thm:small-values}.  The values for \(n=10,11,12\) are
\cref{thm:finite-n10-n12}, and the value at \(n=13\) is
\cref{thm:n13-value}.  Finally, \cref{thm:large-range} proves the stated
formula for every \(n\ge14\).  The constructions in
\cref{sec:foundations-small} give equality in every case.
\end{proof}

\section*{Reproducibility statement}

The distributed supplementary checker uses exact integer, rational, and
symbolic arithmetic.  Its logical role is secondary verification of one
independent closure whose direct proof is already printed in the manuscript.
No optimizer, satisfiability result,
configuration catalogue, support search, or orbit search is a premise of
the proof.  This supplementary replay is not a Lean formalization.

\bibliographystyle{alpha}
\bibliography{references}

\end{document}
